% !TeX program = xelatex
% Standalone English volume. Use XeLaTeX until the references stabilize.
\documentclass[a4paper,11pt]{article}
\usepackage[T1]{fontenc}
\usepackage{lmodern}
\usepackage[english]{babel}
\usepackage[left=26mm,right=26mm,top=25mm,bottom=25mm,headheight=15pt,headsep=7mm,footskip=11mm]{geometry}
\usepackage{amsmath,amssymb,mathtools,amsthm}
\newtheorem{proposition}{Proposition}[section]
\newtheorem{lemma}[proposition]{Lemma}
\newtheorem{theorem}[proposition]{Theorem}
\newtheorem{corollary}[proposition]{Corollary}
\newtheorem{maintheorem}[proposition]{Main Theorem}
\theoremstyle{definition}
\newtheorem{definition}[proposition]{Definition}
\theoremstyle{remark}
\newtheorem{remark}[proposition]{Remark}
\makeatletter
\renewcommand*\l@section{\@dottedtocline{1}{0em}{2.4em}}
\renewcommand*\l@subsection{\@dottedtocline{2}{1.45em}{3.45em}}
\renewcommand*\l@subsubsection{\@dottedtocline{3}{3.8em}{4.55em}}
\makeatother
\usepackage{xcolor,fancyhdr,titlesec,needspace,xurl}
\usepackage{longtable,multicol}
\usepackage{hyperref}
\definecolor{ink}{HTML}{1A3542}
\definecolor{muted}{HTML}{4F606A}
\hypersetup{unicode=true,colorlinks=true,linkcolor=ink,filecolor=ink,urlcolor=ink,citecolor=ink,bookmarksnumbered=true,
 pdfdisplaydoctitle=true,pdflang={en},
 pdftitle={The generalized Riemann hypothesis for Dirichlet L-functions. Edition 930, Volume XXIV},
 pdfsubject={A natural-language reconstruction of the proposed proof},
 pdfauthor={A mathematical reconstruction}}
\titleformat{\section}{\Large\bfseries\color{ink}}{\thesection}{.65em}{}
\titlespacing*{\section}{0pt}{.8em}{.65em}
\setlength{\parindent}{1.2em}
\linespread{1.045}
\setlength{\parskip}{.22em plus .05em minus .03em}
\setlength{\emergencystretch}{2em}
\setcounter{topnumber}{2}
\setcounter{tocdepth}{3}\setcounter{secnumdepth}{3}
\widowpenalty=4500
\clubpenalty=4500
\pagestyle{fancy}\fancyhf{}
\fancyhead[L]{\small\color{muted}Dirichlet GRH \textperiodcentered\ Volume I}
\fancyhead[R]{\small\color{muted}Dirichlet zeros and power intervals}
\fancyfoot[C]{\small\thepage}
\renewcommand{\headrulewidth}{.3pt}
\newcommand{\lead}[1]{\par\smallskip\noindent{\bfseries\color{ink}#1.}\enspace}
\newcommand{\sour}[2]{\unskip\penalty150\hspace{.35em}{\small\hyperlink{ref#1}{[#2]}}}
\titleformat{\subsection}{\large\bfseries}{\thesubsection}{.6em}{}
\titlespacing*{\subsection}{0pt}{1.15em plus .15em}{.4em}
\setlength{\abovedisplayskip}{9pt plus 2pt minus 3pt}
\setlength{\belowdisplayskip}{9pt plus 2pt minus 3pt}
\newcommand{\RR}{\mathbb R}\newcommand{\CC}{\mathbb C}\newcommand{\QQ}{\mathbb Q}
\newcommand{\ZZ}{\mathbb Z}\newcommand{\ord}{\operatorname{ord}}
\newcommand{\Frob}{\operatorname{Frob}}\newcommand{\Gal}{\operatorname{Gal}}
\newcommand{\len}{\operatorname{length}}\newcommand{\Res}{\operatorname*{Res}}
\newcommand{\wmean}{\operatorname{av}}\newcommand{\Ker}{\operatorname{ker}}
\newcommand{\Lh}{\mathcal L}\newcommand{\Gg}{\mathcal G}
\newcommand{\real}{\operatorname{Re}}\newcommand{\imag}{\operatorname{Im}}
\newcommand{\abs}[1]{\left|#1\right|}\newcommand{\eps}{\varepsilon}
\newcommand{\anchoritem}[2]{\par\Needspace{2.8\baselineskip}\hypertarget{ref#1}{}\noindent{\bfseries\color{ink}#2.}\enspace}
\newcommand{\srcloc}[2]{{\small\nolinkurl{#1.lean}}, line~#2}
\numberwithin{equation}{section}
\raggedbottom

\newcommand{\Sel}{\mathcal S}
\newcommand{\Cl}{\operatorname{Cl}}
\newcommand{\Aut}{\operatorname{Aut}}
\newcommand{\GL}{\operatorname{GL}}
\newcommand{\OO}{\mathcal O}
\hypersetup{bookmarksdepth=3}


\addto\captionsenglish{\renewcommand{\partname}{Volume}}


\makeatletter
\expandafter\def\csname r@cor:allsample341\endcsname{{22.15}{74}{}{proposition.22.15}{Dirichlet-GRH-v930-Volume-IV.pdf}}
\expandafter\def\csname r@cor:asym181\endcsname{{104.10}{9}{}{proposition.104.10}{Dirichlet-GRH-v930-Volume-XX.pdf}}
\expandafter\def\csname r@cor:rank111\endcsname{{9.14}{119}{}{proposition.9.14}{Dirichlet-GRH-v930-Volume-I.pdf}}
\expandafter\def\csname r@cor:splitproduct334\endcsname{{83.20}{68}{}{proposition.83.20}{Dirichlet-GRH-v930-Volume-XVI.pdf}}
\expandafter\def\csname r@eq:letters111\endcsname{{9.1}{117}{}{equation.9.1}{Dirichlet-GRH-v930-Volume-I.pdf}}
\expandafter\def\csname r@eq:signedfibers331\endcsname{{99.1}{53}{}{equation.99.1}{Dirichlet-GRH-v930-Volume-XIX.pdf}}
\expandafter\def\csname r@lem:endpoint305\endcsname{{122.1}{56}{}{proposition.122.1}{Dirichlet-GRH-v930-Volume-XXIII.pdf}}
\expandafter\def\csname r@lem:finite167\endcsname{{72.3}{46}{}{proposition.72.3}{Dirichlet-GRH-v930-Volume-XIV.pdf}}
\expandafter\def\csname r@lem:gauss89\endcsname{{92.8}{8}{}{proposition.92.8}{Dirichlet-GRH-v930-Volume-XVIII.pdf}}
\expandafter\def\csname r@lem:highestline347\endcsname{{83.27}{74}{}{proposition.83.27}{Dirichlet-GRH-v930-Volume-XVI.pdf}}
\expandafter\def\csname r@lem:hook86\endcsname{{92.4}{5}{}{proposition.92.4}{Dirichlet-GRH-v930-Volume-XVIII.pdf}}
\expandafter\def\csname r@lem:oneweight355\endcsname{{120.10}{21}{}{proposition.120.10}{Dirichlet-GRH-v930-Volume-XXIII.pdf}}
\expandafter\def\csname r@lem:openflags353\endcsname{{120.7}{18}{}{proposition.120.7}{Dirichlet-GRH-v930-Volume-XXIII.pdf}}
\expandafter\def\csname r@lem:sphere60\endcsname{{66.8}{51}{}{proposition.66.8}{Dirichlet-GRH-v930-Volume-XIII.pdf}}
\expandafter\def\csname r@nav:I-28\endcsname{{6.6}{32}{}{subsection.6.6}{Dirichlet-GRH-v930-Volume-I.pdf}}
\expandafter\def\csname r@prop:adjtheta234\endcsname{{110.5}{41}{}{proposition.110.5}{Dirichlet-GRH-v930-Volume-XXI.pdf}}
\expandafter\def\csname r@prop:cartan244\endcsname{{95.2}{66}{}{proposition.95.2}{Dirichlet-GRH-v930-Volume-XVIII.pdf}}
\expandafter\def\csname r@prop:cmplaces335\endcsname{{67.13}{87}{}{proposition.67.13}{Dirichlet-GRH-v930-Volume-XIII.pdf}}
\expandafter\def\csname r@prop:column110\endcsname{{86.29}{23}{}{proposition.86.29}{Dirichlet-GRH-v930-Volume-XVII.pdf}}
\expandafter\def\csname r@prop:extend300\endcsname{{102.3}{90}{}{proposition.102.3}{Dirichlet-GRH-v930-Volume-XIX.pdf}}
\expandafter\def\csname r@prop:normalframe349\endcsname{{104.34}{29}{}{proposition.104.34}{Dirichlet-GRH-v930-Volume-XX.pdf}}
\expandafter\def\csname r@prop:ray50\endcsname{{6.30}{63}{}{proposition.6.30}{Dirichlet-GRH-v930-Volume-I.pdf}}
\expandafter\def\csname r@prop:res23\endcsname{{6.11}{32}{}{proposition.6.11}{Dirichlet-GRH-v930-Volume-I.pdf}}
\expandafter\def\csname r@prop:shift203\endcsname{{110.11}{47}{}{proposition.110.11}{Dirichlet-GRH-v930-Volume-XXI.pdf}}
\expandafter\def\csname r@prop:slice300\endcsname{{107.17}{88}{}{proposition.107.17}{Dirichlet-GRH-v930-Volume-XX.pdf}}
\expandafter\def\csname r@prop:source305\endcsname{{112.3}{76}{}{proposition.112.3}{Dirichlet-GRH-v930-Volume-XXI.pdf}}
\expandafter\def\csname r@prop:split173\endcsname{{72.15}{58}{}{proposition.72.15}{Dirichlet-GRH-v930-Volume-XIV.pdf}}
\expandafter\def\csname r@prop:split233\endcsname{{94.17}{51}{}{proposition.94.17}{Dirichlet-GRH-v930-Volume-XVIII.pdf}}
\expandafter\def\csname r@sec:char94\endcsname{{121.16}{51}{}{subsection.121.16}{Dirichlet-GRH-v930-Volume-XXIII.pdf}}
\expandafter\def\csname r@sec:covectorpaths352\endcsname{{73.1}{72}{}{subsection.73.1}{Dirichlet-GRH-v930-Volume-XIV.pdf}}
\expandafter\def\csname r@sec:fourierslopes354\endcsname{{120.5}{19}{}{subsection.120.5}{Dirichlet-GRH-v930-Volume-XXIII.pdf}}
\expandafter\def\csname r@sec:graphtransport337\endcsname{{71.20}{38}{}{subsection.71.20}{Dirichlet-GRH-v930-Volume-XIV.pdf}}
\expandafter\def\csname r@sec:lowerlattice347\endcsname{{83.13}{73}{}{subsection.83.13}{Dirichlet-GRH-v930-Volume-XVI.pdf}}
\expandafter\def\csname r@sec:missing\endcsname{{6.6}{32}{}{subsection.6.6}{Dirichlet-GRH-v930-Volume-I.pdf}}
\expandafter\def\csname r@seriesguide\endcsname{{1.1}{4}{}{subsection.1.1}{Dirichlet-GRH-v930-Volume-I.pdf}}
\expandafter\def\csname r@thm:adapted314\endcsname{{122.14}{72}{}{proposition.122.14}{Dirichlet-GRH-v930-Volume-XXIII.pdf}}
\expandafter\def\csname r@thm:archcolumn313\endcsname{{121.7}{37}{}{proposition.121.7}{Dirichlet-GRH-v930-Volume-XXIII.pdf}}
\expandafter\def\csname r@thm:charttrace320\endcsname{{122.5}{61}{}{proposition.122.5}{Dirichlet-GRH-v930-Volume-XXIII.pdf}}
\expandafter\def\csname r@thm:class61\endcsname{{65.11}{31}{}{proposition.65.11}{Dirichlet-GRH-v930-Volume-XIII.pdf}}
\expandafter\def\csname r@thm:classden312\endcsname{{122.13}{71}{}{proposition.122.13}{Dirichlet-GRH-v930-Volume-XXIII.pdf}}
\expandafter\def\csname r@thm:classtrace326\endcsname{{122.17}{75}{}{proposition.122.17}{Dirichlet-GRH-v930-Volume-XXIII.pdf}}
\expandafter\def\csname r@thm:closedperiod318\endcsname{{123.2}{78}{}{proposition.123.2}{Dirichlet-GRH-v930-Volume-XXIII.pdf}}
\expandafter\def\csname r@thm:common111\endcsname{{9.13}{118}{}{proposition.9.13}{Dirichlet-GRH-v930-Volume-I.pdf}}
\expandafter\def\csname r@thm:commonorbit305\endcsname{{121.5}{34}{}{proposition.121.5}{Dirichlet-GRH-v930-Volume-XXIII.pdf}}
\expandafter\def\csname r@thm:compact313\endcsname{{95.6}{71}{}{proposition.95.6}{Dirichlet-GRH-v930-Volume-XVIII.pdf}}
\expandafter\def\csname r@thm:compactrow335\endcsname{{90.14}{95}{}{proposition.90.14}{Dirichlet-GRH-v930-Volume-XVII.pdf}}
\expandafter\def\csname r@thm:conductor232\endcsname{{19.15}{12}{}{proposition.19.15}{Dirichlet-GRH-v930-Volume-IV.pdf}}
\expandafter\def\csname r@thm:covectorpaths352\endcsname{{73.1}{73}{}{proposition.73.1}{Dirichlet-GRH-v930-Volume-XIV.pdf}}
\expandafter\def\csname r@thm:criticalmarks339\endcsname{{22.12}{71}{}{proposition.22.12}{Dirichlet-GRH-v930-Volume-IV.pdf}}
\expandafter\def\csname r@thm:elliptic312\endcsname{{67.12}{86}{}{proposition.67.12}{Dirichlet-GRH-v930-Volume-XIII.pdf}}
\expandafter\def\csname r@thm:endpoint315\endcsname{{122.15}{73}{}{proposition.122.15}{Dirichlet-GRH-v930-Volume-XXIII.pdf}}
\expandafter\def\csname r@thm:finitecolumn313\endcsname{{122.10}{67}{}{proposition.122.10}{Dirichlet-GRH-v930-Volume-XXIII.pdf}}
\expandafter\def\csname r@thm:fourierslopes354\endcsname{{120.8}{20}{}{proposition.120.8}{Dirichlet-GRH-v930-Volume-XXIII.pdf}}
\expandafter\def\csname r@thm:framechange314\endcsname{{122.3}{58}{}{proposition.122.3}{Dirichlet-GRH-v930-Volume-XXIII.pdf}}
\expandafter\def\csname r@thm:frob244\endcsname{{102.5}{92}{}{proposition.102.5}{Dirichlet-GRH-v930-Volume-XIX.pdf}}
\expandafter\def\csname r@thm:global158\endcsname{{87.26}{50}{}{proposition.87.26}{Dirichlet-GRH-v930-Volume-XVII.pdf}}
\expandafter\def\csname r@thm:incidence338\endcsname{{86.37}{29}{}{proposition.86.37}{Dirichlet-GRH-v930-Volume-XVII.pdf}}
\expandafter\def\csname r@thm:intrinsicminor315\endcsname{{122.11}{68}{}{proposition.122.11}{Dirichlet-GRH-v930-Volume-XXIII.pdf}}
\expandafter\def\csname r@thm:last244\endcsname{{67.11}{84}{}{proposition.67.11}{Dirichlet-GRH-v930-Volume-XIII.pdf}}
\expandafter\def\csname r@thm:lastrow311\endcsname{{121.10}{41}{}{proposition.121.10}{Dirichlet-GRH-v930-Volume-XXIII.pdf}}
\expandafter\def\csname r@thm:lattice55\endcsname{{65.3}{23}{}{proposition.65.3}{Dirichlet-GRH-v930-Volume-XIII.pdf}}
\expandafter\def\csname r@thm:levels88\endcsname{{7.11}{97}{}{proposition.7.11}{Dirichlet-GRH-v930-Volume-I.pdf}}
\expandafter\def\csname r@thm:lowerlattice347\endcsname{{83.28}{74}{}{proposition.83.28}{Dirichlet-GRH-v930-Volume-XVI.pdf}}
\expandafter\def\csname r@thm:markedstrict337\endcsname{{22.11}{70}{}{proposition.22.11}{Dirichlet-GRH-v930-Volume-IV.pdf}}
\expandafter\def\csname r@thm:minors311\endcsname{{122.9}{65}{}{proposition.122.9}{Dirichlet-GRH-v930-Volume-XXIII.pdf}}
\expandafter\def\csname r@thm:mixedfourier334\endcsname{{83.19}{68}{}{proposition.83.19}{Dirichlet-GRH-v930-Volume-XVI.pdf}}
\expandafter\def\csname r@thm:moment286\endcsname{{105.3}{34}{}{proposition.105.3}{Dirichlet-GRH-v930-Volume-XX.pdf}}
\expandafter\def\csname r@thm:neighbour93\endcsname{{94.11}{46}{}{proposition.94.11}{Dirichlet-GRH-v930-Volume-XVIII.pdf}}
\expandafter\def\csname r@thm:neutral315\endcsname{{122.4}{59}{}{proposition.122.4}{Dirichlet-GRH-v930-Volume-XXIII.pdf}}
\expandafter\def\csname r@thm:norm58\endcsname{{65.25}{43}{}{proposition.65.25}{Dirichlet-GRH-v930-Volume-XIII.pdf}}
\expandafter\def\csname r@thm:one61\endcsname{{65.13}{33}{}{proposition.65.13}{Dirichlet-GRH-v930-Volume-XIII.pdf}}
\expandafter\def\csname r@thm:pathnorm352\endcsname{{119.11}{11}{}{proposition.119.11}{Dirichlet-GRH-v930-Volume-XXIII.pdf}}
\expandafter\def\csname r@thm:pointlattice353\endcsname{{99.20}{60}{}{proposition.99.20}{Dirichlet-GRH-v930-Volume-XIX.pdf}}
\expandafter\def\csname r@thm:poscost316\endcsname{{123.5}{81}{}{proposition.123.5}{Dirichlet-GRH-v930-Volume-XXIII.pdf}}
\expandafter\def\csname r@thm:posframe316\endcsname{{123.1}{76}{}{proposition.123.1}{Dirichlet-GRH-v930-Volume-XXIII.pdf}}
\expandafter\def\csname r@thm:prime209\endcsname{{93.9}{28}{}{proposition.93.9}{Dirichlet-GRH-v930-Volume-XVIII.pdf}}
\expandafter\def\csname r@thm:prime210\endcsname{{93.17}{34}{}{proposition.93.17}{Dirichlet-GRH-v930-Volume-XVIII.pdf}}
\expandafter\def\csname r@thm:projclass326\endcsname{{123.14}{95}{}{proposition.123.14}{Dirichlet-GRH-v930-Volume-XXIII.pdf}}
\expandafter\def\csname r@thm:pullback311\endcsname{{121.6}{35}{}{proposition.121.6}{Dirichlet-GRH-v930-Volume-XXIII.pdf}}
\expandafter\def\csname r@thm:range239\endcsname{{88.1}{52}{}{proposition.88.1}{Dirichlet-GRH-v930-Volume-XVII.pdf}}
\expandafter\def\csname r@thm:rangenew50\endcsname{{6.29}{62}{}{proposition.6.29}{Dirichlet-GRH-v930-Volume-I.pdf}}
\expandafter\def\csname r@thm:rankone354\endcsname{{104.36}{31}{}{proposition.104.36}{Dirichlet-GRH-v930-Volume-XX.pdf}}
\expandafter\def\csname r@thm:rationalperiod342\endcsname{{86.38}{30}{}{proposition.86.38}{Dirichlet-GRH-v930-Volume-XVII.pdf}}
\expandafter\def\csname r@thm:rationalrank302\endcsname{{102.4}{91}{}{proposition.102.4}{Dirichlet-GRH-v930-Volume-XIX.pdf}}
\expandafter\def\csname r@thm:realtrans300\endcsname{{107.3}{75}{}{proposition.107.3}{Dirichlet-GRH-v930-Volume-XX.pdf}}
\expandafter\def\csname r@thm:region198\endcsname{{20.13}{33}{}{proposition.20.13}{Dirichlet-GRH-v930-Volume-IV.pdf}}
\expandafter\def\csname r@thm:signedfibers331\endcsname{{99.12}{53}{}{proposition.99.12}{Dirichlet-GRH-v930-Volume-XIX.pdf}}
\expandafter\def\csname r@thm:signedrow335\endcsname{{90.12}{94}{}{proposition.90.12}{Dirichlet-GRH-v930-Volume-XVII.pdf}}
\expandafter\def\csname r@thm:splitcompact323\endcsname{{123.7}{84}{}{proposition.123.7}{Dirichlet-GRH-v930-Volume-XXIII.pdf}}
\expandafter\def\csname r@thm:splitflag343\endcsname{{9.20}{123}{}{proposition.9.20}{Dirichlet-GRH-v930-Volume-I.pdf}}
\expandafter\def\csname r@thm:splitinput315\endcsname{{123.4}{80}{}{proposition.123.4}{Dirichlet-GRH-v930-Volume-XXIII.pdf}}
\expandafter\def\csname r@thm:splitpath322\endcsname{{123.9}{87}{}{proposition.123.9}{Dirichlet-GRH-v930-Volume-XXIII.pdf}}
\expandafter\def\csname r@thm:transport59\endcsname{{65.9}{28}{}{proposition.65.9}{Dirichlet-GRH-v930-Volume-XIII.pdf}}
\expandafter\def\csname r@thm:twists\endcsname{{5.1}{19}{}{proposition.5.1}{Dirichlet-GRH-v930-Volume-I.pdf}}
\expandafter\def\csname r@thm:vac242\endcsname{{67.8}{81}{}{proposition.67.8}{Dirichlet-GRH-v930-Volume-XIII.pdf}}
\expandafter\def\csname r@thm:vectorchain353\endcsname{{73.2}{74}{}{proposition.73.2}{Dirichlet-GRH-v930-Volume-XIV.pdf}}
\expandafter\def\csname r@thm:weylperiod351\endcsname{{119.10}{10}{}{proposition.119.10}{Dirichlet-GRH-v930-Volume-XXIII.pdf}}
\makeatother
\begin{document}
\setcounter{part}{23}
\fancyhead[L]{\small\color{muted}Dirichlet GRH \textperiodcentered\ Volume XXIV}
\fancyhead[R]{\small\color{muted}Marked theta columns and covectors}
\thispagestyle{plain}\begin{center}
{\LARGE\bfseries The generalized Riemann hypothesis\\[2pt]for Dirichlet \(L\)-functions}\\[6pt]
{\small 27 November 2026 \quad\textperiodcentered\quad Edition 930}\end{center}
\part{Marked theta columns and terminal covectors}\label{part:XXIV}
\begin{abstract}\noindent We express the endpoint pairing in terms of signed tests on marked classes. Summing over the intermediate residue classes gives formulas for the composite terminal covectors, their Gram matrices and their integral moments.\end{abstract}
\setcounter{section}{124}
\noindent The \hyperref[seriesguide]{guide in Volume~I} lists all twenty-seven volumes.
\begingroup\small\setlength{\parskip}{0pt}\tableofcontents\endgroup
\clearpage

\section{Denominators and norms of the endpoint pairing}\label{nav:g-XII-3}\label{nav:volume-XXIV}

\subsection{Integral tests for the complex theta period}\label{nav:IX-61}
Keep the matching complex operator \(\Theta\) of
Theorem~\ref{thm:range239}, with its given real tests
and measures. For the covector \(b\) and the definite
endpoint test \(Q\), consider the scalar
\(\langle\Theta S b,Q\rangle\). It is also the ordinary
Bessel functional on the reverse lift \(\Theta^*Q\)
in the normalization of Proposition~\ref{prop:adjtheta234}.

\begin{theorem}[Two integral tests at the original finite levels]\leavevmode
\label{thm:pair305}
Choose a good orbit of Theorem~\ref{thm:commonorbit305}.
For each prescribed word on these framed data and each
terminal character nontrivial at every place above \(T\), there is a
literal integral character column \(b\) on that orbit and
an integral endpoint coefficient test \(Q\) at the original
level of the complex iteration such that
\[
\begin{split}
 |\langle\Theta S b,Q\rangle|^2 \ \ge\ &
     \alpha e^{-\Delta_x} M_{J,{\cal O}}\,
     c_{\rho,K^0}|Z|\,
        \frac{d_\rho B_s^{-2a_s}}{k^{e_s}}\,
        \frac{m_e^2 w_e B_e^{-4a_e}}{d_e h_e N^{2e_e}}>0 .
\end{split}
\]
The tests have norms at most
\((c_{\rho,K^0}|Z|)^{1/2}B_s^{a_s}\) and
\(m_e\sqrt{w_e}L_eB_e^{2a_e}\), respectively.
For a finite list of \(l\) such pairs with a common endpoint
coefficient space, integral lattice and finite level
character, one endpoint test serves all the selected
columns with at least half their respective lower bounds
on absolute value.
\end{theorem}
\begin{proof}
Take \(b=b(\ell)\) from Proposition~\ref{prop:source305}
and apply Lemma~\ref{lem:endpoint305} to \(F=\Theta S b\).
Multiplication gives the displayed inequality. The mean
\(M_{J,{\cal O}}\) is positive on these orbits for large
\(x\) by Theorem~\ref{thm:commonorbit305}. The source
norm bound follows from the upper dual frame bound and
the class norm of the unnormalized character sum. The
other norm bound is the one in the endpoint lemma.

For the last assertion call the individual endpoint tests
\(Q_i\) and the forward lifts \(F_i\). Consider
\(Q=\sum_{j=1}^l n_j Q_j\) for independent integers
\(0\le n_j\le l\) with uniform distribution. For a fixed
\(i\), condition on all coefficients but \(n_i\).
The \(l+1\) possible pairings are spaced by
\(|\langle F_i,Q_i\rangle|>0\). An open disk of half
this radius contains at most one. With probability at
least \(1-l/(l+1)>0\) all the lower bounds on absolute
value, divided by two, hold at once. Such a \(Q\) is
in the common integral lattice and has norm at most
\(l\sum_j\|Q_j\|\). This proves the assertion for the
same endpoints and levels.
\end{proof}

The two integrality statements refer to the coefficient
models of the input tests. Relating this Petersson scalar
to the arithmetic theta kernel and its distinguished
prime divisibilities requires the normalization of the
real tests, the measures and the integral projection.
The bounds above retain these factors explicitly.

\subsection{Archimedean and finite bounds for the same last-edge minors}\label{nav:IX-62}
Use the actual last-edge matrix \(Z\) of Theorem~\ref{thm:minors311}.
Write \(l=ns,\ D=fl,\ l_\tau=|\lambda_\tau|\), and
\(h_\lambda=\|P_*\|_{\rm F}^2\). Use a single Galois number
field \(J\) containing the arithmetic data, finite character
values, and the finite test lists. The archimedean volume
factors and the marked vacuum scalar have already been removed
in the definition of \(Z\). At an embedding \(\sigma:J\to\CC\)
let \(A_\sigma,\nu_\sigma\) be the area and the mean squared
elliptic theta norm of Theorem~\ref{thm:transport59}, and
\(c_\sigma=\pi/A_\sigma\). Put
\[
 \begin{split}
 H_l(a)&=\frac{(a+l-1)!}{(l-1)!}(e/a)^a\quad(a\ge1),
                         \qquad H_l(0)=1,\\
 B_\sigma^{\,2}
   &=b_\sigma^2\,\big[e\max(1,2l\,\operatorname{rd}(E)/\pi)
                     \,\nu_\sigma\big]^D
           c_\sigma^{\,d}\,h_\lambda\prod_\tau H_l(l_\tau).
 \end{split}
\]
The factor \(b_\sigma\) bounds the absolute coefficient sum of
each finite test in translates of its normalized lattice vacuum.

Take the fine representatives of Theorem~\ref{thm:class61}
away from the marked places. Move the right-level translations
from the two single-chart tests to the finite kernel. There
are only finitely many such translations. In a split position
model a term \(a\mathbf1_{y+D_v}\), with \(D_v\) a full lattice,
has coefficient sum at most
\[
             |\sigma(a/a_0)|[D_v:D_v\cap L_0]
\]
in translates and modulations of a reference position vacuum
of amplitude \(a_0\) on \(L_0\). Indeed split the coset by
\(D_v\cap L_0\). In one translate of \(L_0\) its subcoset
indicator is the finite Fourier sum on \(L_0/(D_v\cap L_0)\),
where each character contributes the reciprocal group order times
a root of unity. A position or frequency translate of the
reference vacuum only changes the phases and the cosets.
Multiply the bounds at different places, sum absolute
coefficients for a finite combination, and take the maximum
over the two level translations. This gives \(b_\sigma\).
For the split indicator of Proposition~\ref{prop:split173} in
its unit-volume position lattice the local factor is one.
An unchanged unit vacuum also has factor one. Retain the
central translate at the additional split place of
Proposition~\ref{prop:split233} as a moved vacuum, with its
splitting factor. The class transports of
Theorem~\ref{thm:one61} preserve the absolute coefficient
sum at every embedding. These statements use the same finite
kernel in the period and in \(b_\sigma\).

Let \(S\) be any finite set of the good primes of
Theorem~\ref{thm:frob244} for these tests. Each \(p\in S\)
is at least five, inert in the marked quadratic field, unramified
in \(E\), and prime to the finite CM level. The split indicator expansion
has denominators supported at its residue characteristic.
The transporting prime may be chosen completely split in
\(E\). At a good \(p\) the unmarked factors are the integral
vacua of that theorem. Use its common integral basis and
write \(g_p\) for its explicit lower bound on a last-edge
matrix valuation, using the same halves of the Frobenius
cycles for all entries. For each sign put \(B_\pm=1\) if
the unitary family is used, and \(B_-=n^{d/2}\) or
\(B_+=s^{d/2}\) when the corresponding integral family is used.

\begin{theorem}[Both bounds for one nonzero arithmetic minor]\leavevmode
\label{thm:bothlast311}
For the very minor \(z_j\) chosen in Theorem~\ref{thm:minors311},
\[
 \begin{aligned}
  \operatorname{ord}_{\mathfrak p} z_j
       &\ge j\,e(\mathfrak p/p)g_p
                 &&(\mathfrak p\mid p,\ p\in S),\\
  |\sigma z_j|
       &\le j^{j/2}\big(i_-i_+B_-B_+ B_\sigma\big)^j
                 &&(\sigma:J\to\CC).
 \end{aligned}
\]
At the original embedding it also satisfies
\(|z_j|\ge j!(i_-i_+L_*)^j\). In degree one the chosen
target can detect any specified nonzero incoming Bessel row
as in Theorem~\ref{thm:lastrow311}.

For the two integral families, let \(\mathfrak a\) be the
fractional ideal generated by the coordinate minor-product
derivatives for all last-edge class pairs and level translates.
These are \(M(e_I\otimes e_J)(\partial)G_{uv}(0)\) on the
coordinate exterior tensor bases. Take all finite factors in
the actual normalized class sections. With
\[
 b_{\rm den}=\frac{1}{[J:\QQ]}
      \sum_{\mathfrak p\nmid S}
        [-\operatorname{ord}_{\mathfrak p}\mathfrak a]_+
                                             \log N\mathfrak p ,
\]
one has on the same field
\[
 \sum_{p\in S}g_p\log p
 \le b_{\rm den}+\tfrac12\log j
       +\log(i_-i_+)+\tfrac d2\log(ns)
       +\frac{1}{[J:\QQ]}\sum_\sigma\log B_\sigma .
\]
\end{theorem}
\begin{proof}
At the chosen fine levels all stabilizers are trivial. The
class integral of Theorem~\ref{thm:last244} therefore gives
each entry as a sum of \(i_-i_+\) derivatives with finite
character factors of absolute value one at every embedding.
Within each term \(Y,Z'\) are the two exterior tensors from
the single charts. At every place above \(S\) these tensors
are integral for either family. We can keep the original
representatives by moving the two level elements onto the
finite test, as in the construction of \(b_\sigma\). At \(p\)
we can take these elements to be the identity since the unmarked
factor of the level is the same at the fine and original
levels, and the auxiliary split congruence for the fine
charts is away from \(S\). The common coordinate calculation
of Theorem~\ref{thm:frob244} now gives valuation at least
\(e(\mathfrak p/p)g_p\) in each term. Indeed each monomial
of the minor product has degree \(l_\tau\) at \(\tau\), and
each full-height minor uses every column. These are exactly
the two degree bounds used in that theorem. Taking sums and
then expanding the determinant proves the first inequality.

At \(\sigma\), Theorem~\ref{thm:transport59} and the triangle
inequality for torsion translates bound the squared metric
supremum of the conjugated theta section by
\[
        b_\sigma^2[e\max(1,2l\,\operatorname{rd}(E)/\pi)\nu_\sigma]^D .
\]
The block-sphere estimate of
Lemma~\ref{lem:sphere60} consequently bounds the Fischer
squared norm of its Taylor component by this expression
times \(c_\sigma^d\prod_\tau H_l(l_\tau)\). It applies to
these class sections in their conjugate marked differentials,
with the blocks permuted with the CM embeddings.
The unitary changes of complex frame in the
CM trace coordinates preserve the Fischer norm. Equivalently
one can first use the unitary change of variables and then
project to the indicated block degrees. The differential
is paired with the component by the monomial identity
\(\partial^\alpha Z^\beta|_0=\alpha!\,\mathbf1_{\alpha=\beta}\).
The isometry in Proposition~\ref{prop:cartan244} gives
\(\|M(\sigma Y\otimes\sigma Z')\|_{\rm F}
=\sqrt{h_\lambda}\|\sigma Y\|\|\sigma Z'\|\).
Complex conjugation in a Fischer pairing does not change
either norm. Cauchy--Schwarz bounds one term by
\(B_-B_+ B_\sigma\). Summing the level translates and using
Hadamard's determinant inequality gives the second assertion.

For integral tests each entry belongs to \(\mathfrak a\)
since the tensor entries are algebraic integers and the finite
character values are roots of unity. Thus \(z_j\in\mathfrak a^j\),
and this ideal is nonzero by the lower bound. Outside \(S\)
the total negative valuation of \(z_j\), weighted by
\(\log N\mathfrak p\), is at most \(j[J:\QQ]b_{\rm den}\).
The product formula and \(\sum_{\mathfrak p\mid p}
e(\mathfrak p/p)f(\mathfrak p/p)=[J:\QQ]\) give the final
inequality from the bound at the complex embeddings.
\end{proof}

For the paired fields with the prescribed mark of
Theorem~\ref{thm:norm58} the embedding average in the bound
can be evaluated further. Write \(d_K\) for the absolute
quadratic discriminant and \(D_\xi=L'(1,\xi_{d_K})/L(1,\xi_{d_K})\).
Set \(L_K=\log(\sqrt{d_K}/(2\pi))-\gamma+D_\xi\) and
\(V_K=\tfrac14\log d_K-\tfrac12\gamma+\tfrac12D_\xi\).
Writing a bar for the embedding mean, one has on the
original intervals
\[
 2\,\overline{\log B_\sigma}
 =2\,\overline{\log b_\sigma}+dL_K+D V_K
     +(k-1)\sum_\tau B_\tau(\log B_\tau-1)
     +O_{\psi,\varepsilon}(x^{C_0}) .
\]
Here \(B_\tau\) and \(k\) are the offsets and the weight
parameter of Corollary~\ref{cor:asym181}.
Indeed \(\overline{\log c_\sigma}=L_K\) and
\(\overline{\log\nu_\sigma}=V_K\) by Theorem~\ref{thm:norm58}.
Lemma~\ref{lem:hook86} gives
\[
 h_\lambda=D_-^{-1}\prod_{\tau,i}
                \frac{(\lambda_{\tau,i}+n-i)!}{(n-i)!}.
\]
The first \(k-1\) large parts at each \(\tau\) are
\(B_\tau+O_{\psi,\varepsilon}(x^{C_1})\), and all remaining
parts are polynomially bounded, as in
Corollary~\ref{cor:asym181}. Stirling's formula and
\(\log(2+\max l_\tau)=O_{\psi,\varepsilon}(x^\varepsilon+\log(2x))\)
give the displayed leading term for \(\log h_\lambda\) with
polynomial error. The Weyl dimension formula bounds
\(\log D_-\) by \(fn^2\log(2+n+\max l_\tau)\).
The sphere factor costs at most
\(\sum_\tau l(1+\log(1+l_\tau/l))\) in its logarithm by
Lemma~\ref{lem:sphere60}. Finally \(f,n,s,\log\operatorname{rd}(E)\)
are polynomially bounded on this family by the field and
rank bounds, so the remaining vacuum factor costs a power
of \(x\). Also \(\log(ns)=O_{\psi,\varepsilon}(\log(2x))\).
The denominators in
\(b_{\rm den}\) are those of the last-edge CM derivatives.
The preceding lifts still have to supply the arithmetic
coordinates of the actual incoming vector in the full
signed period. Its complex pairing with the chosen target
is controlled by Theorem~\ref{thm:lastrow311}.

\subsection{An input bound for the denominators of the arithmetic period}\label{nav:IX-63}
Take the two integral test families in Theorem~\ref{thm:bothlast311}.
In Theorem~\ref{thm:classden312} use all exponent vectors in the
coordinate polynomials \(M(e_I\otimes e_J)\) for this last edge,
and use its actual finite class sections. For \(S\) as in
Theorem~\ref{thm:bothlast311} define
\[
 h_S=\frac{1}{[J:\QQ]}\sum_{p\notin S}\sum_{\mathfrak p\mid p}
           [J_{\mathfrak p}:\QQ_p]\Gamma_{\mathfrak p}\log p.
\]
Write \(u=\lfloor d/2\rfloor\), \(\epsilon_3=\kappa_3\mathbf1_{3\nmid m}\),
and \(\epsilon_0=\mathbf1_{d_0\nmid m}\) for the input \(m\)
of Theorem~\ref{thm:classden312}.

\begin{theorem}[A common denominator for the endpoint and its minors]\leavevmode
\label{thm:periodden312}
On the fields and intervals of Section~\ref{sec:missing}, every
entry in the arithmetic array \(a_{v,I}(V)\) of
Theorem~\ref{thm:pullback311} becomes integral after multiplication
by the integer
\[
       A=N(mH)^d\,3^{u\epsilon_3}\,
                          d_0^{\epsilon_0\lceil u/2\rceil}.
\]
For the nonzero \(j\)-row minor \(z_j\) of
Theorem~\ref{thm:minors311}, \(A^jz_j\) is an algebraic integer.
For these same minors the comparison of
Theorem~\ref{thm:bothlast311} and its finite contribution satisfy
\[
 \begin{split}
  \sum_{p\in S}g_p\log p
  &\le h_S+\tfrac12\log j+\log(i_-i_+)
        +\tfrac d2\log(ns)
        +\frac1{[J:\QQ]}\sum_\sigma\log B_\sigma,\\
  h_S&\le \log N+d\log(mH)+u\epsilon_3\log3
                             +\tfrac12u\epsilon_0\log d_0 .
 \end{split}
\]
Here \(N,m,H\) and the sharper \(\Gamma_{\mathfrak p}\) depend
only on the finite presentations and the tangent coordinates.
\end{theorem}
\begin{proof}
Every coordinate polynomial is a product or an integral linear
combination of products of minors, hence has integer coefficients.
By Theorem~\ref{thm:classden312}, each of its evaluations on a
class section has valuation at least \(-\Gamma_{\mathfrak p}\).
The same lower bound holds for each one-chart matrix coefficient
and for each entry of \(a_{v,I}(V)\). Indeed the integral tests
have integral exterior tensor coefficients at every place and the
finite character projections in these formulas are unnormalized
sums with roots of unity as coefficients. The factors divided out
in defining the arithmetic matrix are exactly those specified in
Theorem~\ref{thm:pullback311}. The determinant expansion gives
\(v_{\mathfrak p}(z_j)\ge-j\Gamma_{\mathfrak p}\).

Let \(t_p=v_p(m)\). At any prime its contribution from
the translated elliptic factors is at most
\(d t_p+u(\kappa_p-2t_p)_+\) by Theorem~\ref{thm:classden312}.
Here \(t_p\) is an integer, and \(0\le\kappa_p\le1\). Thus the
second summand is \(u\epsilon_3\) at \(3\),
\(u\epsilon_0/2\) at \(d_0\), and zero elsewhere.
Adding \(v_p(N)+d\,v_p(H)\) gives a common bound for
\(\Gamma_{\mathfrak p}\) which vanishes off the finitely many
primes dividing \(3d_0NmH\). Rounding the possible half-integral
loss at \(d_0\) proves the assertions about \(A\). Summing the
unrounded bound, using \(\sum_{\mathfrak p\mid p}
[J_{\mathfrak p}:\QQ_p]=[J:\QQ]\), proves the bound for \(h_S\).

At the primes in \(S\) the same minor has normalized valuation
at least \(jg_p\) by Theorem~\ref{thm:bothlast311}. Outside \(S\)
we have just bounded its negative valuation before computing any
of the derivative values. Apply the product formula to this
nonzero \(z_j\), use the archimedean upper bound of that theorem
at every embedding, and divide by \(j\). This gives the claimed
comparison. No extra factor from \(A\) is charged at the primes
already counted in \(S\).
\end{proof}

The array here acts on the specific incoming row in
Theorem~\ref{thm:pullback311}. Its contraction still involves the
class coordinates of the preceding lifts. An estimate for that
full period on the scale required by Section~\ref{sec:missing}
must also control those coordinates and the growth of the finite
input costs along the interval.

\subsection{The relative different in the signed column Gram matrix}\label{nav:IX-64}
On the original interval use the period entries of
Theorem~\ref{thm:finitecolumn313} with a common directed tower,
degree profile, and marking. The indices \(a,b\) may specify
different signed tests of that family. Use one \(D_0\) as in that
theorem for the entire finite list. On edge \(i\) write \(T_i,Q_i\)
for the integers \(N,Q\) of Theorem~\ref{thm:adapted314} in the
fixed finite vacuum model. Use \(d_i=\sum_\tau d_{i,\tau}\) and
put
\[
 \rho_p=\frac{\operatorname{ord}_{w}\mathfrak D_{E/K}}{e(w/p)},
 \qquad J_{i,p}=v_p(T_i)+\Phi_p(d_i,v_p(Q_i))
                              +\tfrac12d_i\rho_p ,
\]
where \(w\) is any prime of \(E=E_x\) over \(p\). This is
independent of \(w\) since \(E/\QQ\) is abelian and \(K\)
is stable under its Galois group. Let \(S\) be any finite set
of inert good primes as in Theorem~\ref{thm:finitecolumn313}
away from all \(T_iQ_i\). Write \(G_S=\sum_{p\in S}
\sum_i g_{i,p}\log p\).

\begin{theorem}[A basis-independent finite estimate for the Gram products]\leavevmode
\label{thm:gramheight314}
In the original normalization, at each \(\mathfrak p\mid p\),
\[
 v_{\mathfrak p}(D_0^2\mathcal A_a\overline{\mathcal A_b})
 \ge
 \begin{cases}
       2\sum_i g_{i,p} &p\in S,\\
      -2\sum_i J_{i,p} &p\notin S .
 \end{cases}
\]
The same estimates hold for integral linear combinations of
these products after the indicated clearing factor. On all
sufficiently distant intervals of Section~\ref{sec:missing},
\[
 \begin{split}
 \sum_p J_{i,p}\log p
   &=\log T_i+d_i\log Q_i+
                 \tfrac12d_i\log\operatorname{rd}(E)-R_i,\\
 R_i&=\left(\tfrac14d_i-
  \tfrac12\lfloor d_i/2\rfloor\mathbf1_{d_0\nmid Q_i}\right)\log d_0
 \ \ge0 .
 \end{split}
\]
For each nonzero signed period entry with archimedean bounds
\(U_\sigma\) from Theorem~\ref{thm:archcolumn313} it follows that
\[
 G_S\le\log D_0+\sum_i
 \bigl(\log T_i+d_i\log Q_i+\tfrac12d_i\log\operatorname{rd}(E)-R_i\bigr)
               +\frac1{[J:\QQ]}\sum_\sigma\log U_\sigma .
\]
\end{theorem}
\begin{proof}
Fix a rational prime \(p\). On each edge make the single choice
\(B_p\) of Theorem~\ref{thm:adapted314} for all classes and
tests. Compute each signed period by applying the same formal
compact integration and stabilizer sums as in
Theorem~\ref{thm:finitecolumn313}. Its \(D_0\) still clears
the symbolic coefficients. Indeed the incident compact actions,
endpoint coefficients, and physical derivative symbols have been
kept fixed. Only the common normalization of the class germs
has changed. By Theorem~\ref{thm:framechange314} this change
leaves the product of one period with the conjugate of the other
exactly fixed. We may therefore estimate the actual Gram element
at \(p\) with these local choices, including when either factor
vanishes.

In the first factor the coordinate loss at edge \(i\) is at most
\(\sum_\tau d_{i,\tau}v_{\mathfrak p}(\bar\tau\mathfrak b)\)
by Theorem~\ref{thm:adapted314}. In the conjugate factor its
bound uses \(\tau\mathfrak b\). This follows by evaluating the
original bound at the place transported by the fixed complex
conjugation of \(J\). At an edge of opposite orientation the
two are interchanged. Their sum is \(d_i\rho_p\), since
\(\tau\mathfrak b\cdot\bar\tau\mathfrak b=
\tau\mathfrak D_{E/K}\). The two coefficient and torsion losses
are at most \(2v_p(T_i)+2\Phi_p(d_i,v_p(Q_i))\).
At \(p\in S\), the local basis is also integral in physical
coordinates, the finite coefficients are integral and the
product translations have order prime to \(p\). Thus the
Frobenius derivative bound in Theorem~\ref{thm:finitecolumn313}
applies to each factor at every place over \(p\). This proves
the local estimates. The argument can be repeated at each
rational prime on the same algebraic products.

For the sum of the coordinate losses, the norm of the relative
different and the discriminant formula in a tower give
\[
 \sum_p\rho_p\log p=
 \frac{\log N_{E/\QQ}\mathfrak D_{E/K}}{[E:\QQ]}
       =\log\operatorname{rd}(E)-\log\operatorname{rd}(K).
\]
Indeed each \(w\mid p\) contributes its different exponent times
its residue degree to the exponent in the rational norm, and the
sum of the local degrees over \(p\) is \([E:\QQ]\).
Here \(\operatorname{rd}(K)=\sqrt{d_0}\). The congruence
\(d_0\equiv-1\bmod3\) holds on these intervals by
Theorem~\ref{thm:twists}, whose range of prescribed small primes
contains \(3\). Thus \(\kappa_3=0\) in
Theorem~\ref{thm:elliptic312}. The sum of the torsion losses
\(\Phi_p(d_i,v_p(Q_i))\log p\) is exactly
\[
 d_i\log Q_i+\tfrac12\lfloor d_i/2\rfloor
                          \mathbf1_{d_0\nmid Q_i}\log d_0 .
\]
Adding the coordinate and finite coefficient costs gives the
formula with \(R_i\).

For a nonzero entry use its diagonal product. In a Galois field
the embeddings of the conjugate entry are a permutation of the
embeddings of the entry. The archimedean mean of the logarithm
of the absolute value of its diagonal product is therefore
twice the mean for the entry. Apply the product formula and the
first part, divide by two, and use the upper bound \(U_\sigma\)
at all embeddings. At \(p\in S\), \(J_{i,p}=0\). This proves
the comparison with the stated finite data and the actual field
discriminant.
\end{proof}

\section{Signed tests and terminal restrictions}\label{nav:chapter-signed-terminal-tests}
\subsection{The signed column estimate in terms of the original finite tests}\label{nav:IX-65}
Use the signed entries \(\mathcal A_a\), the edge degrees
\(d_i\), the common directed tower, and
\(\rho_p=\operatorname{ord}_w(\mathfrak D_{E/K})/e(w/p)\)
of Theorem~\ref{thm:gramheight314}. At each edge \(i\) take
\(T_i^0,Q_i^0\) on the original list in
Theorem~\ref{thm:neutral315}, and let \(a_{i,\pm}(p)\) be its
two exponents when \(p\) splits in \(K\). Define
\[
 \begin{aligned}
 H_{i,p}&=
 \begin{cases}
 d_i(a_{i,+}(p)+a_{i,-}(p))/2,&p\text{ split in }K,\\
 \Phi_p(d_i,v_p(Q_i^0)),&p\text{ nonsplit in }K,
 \end{cases}\\
 J^0_{i,p}&=v_p(T_i^0)+H_{i,p}+\tfrac12d_i\rho_p,\qquad
 B_i^0=\sum_p J^0_{i,p}\log p .
 \end{aligned}
\]
Choose \(B_{\rm end}\) as in Theorem~\ref{thm:endpoint315}.
Use the original stabilizer multiples \(A_v\) and degree bounds
\(L_{v,\tau}\) for the tensor-polynomial expansions of the
compact integrands before averaging. Put
\[
 D_{\rm intr}=B_{\rm end}\prod_v A_v
                         \prod_{v,\tau} D(n_v,L_{v,\tau}).
\]
Let \(S\) consist of any finite set of the good inert primes
in Theorem~\ref{thm:finitecolumn313}, with the original local
conditions and \(p\nmid\prod_i T_i^0Q_i^0\). Its Frobenius
weights are \(g_{i,p}\) and
\(G_S=\sum_{p\in S}\sum_i g_{i,p}\log p\).

\begin{theorem}[The original finite tests and the Gram products of the signed column]\leavevmode
\label{thm:originalgram315}
For the original projected definite period, at every place
\(\mathfrak p\mid p\), normalized by \(v_{\mathfrak p}(p)=1\),
\[
 v_{\mathfrak p}(D_{\rm intr}^2\mathcal A_a
                              \overline{\mathcal A_b})\ge
 \begin{cases}
    2\sum_i g_{i,p},&p\in S,\\
    -2\sum_i J^0_{i,p},&p\notin S .
 \end{cases}
\]
On the intervals of Theorem~\ref{thm:gramheight314}, the
per-edge expression is
\[
 \begin{aligned}
 B_i^0={}&\log T_i^0+d_i\log Q_i^0
      +\tfrac{d_i}2\log\operatorname{rd}(E)-R_i^0\\
 &-d_i\!\sum_{p\ {\rm split\ in}\ K}
  \left(v_p(Q_i^0)-\tfrac12(a_{i,+}(p)+a_{i,-}(p))\right)\log p,\\
 R_i^0={}&\left(\tfrac{d_i}4-
   \tfrac12\lfloor d_i/2\rfloor{\bf1}_{d_0\nmid Q_i^0}\right)\log d_0 .
 \end{aligned}
\]
For every nonzero entry, with the archimedean bounds
\(U_\sigma\) of Theorem~\ref{thm:archcolumn313} in a common
Galois coefficient field \(J\), this gives
\[
 G_S\le \log D_{\rm intr}+\sum_i B_i^0+
                    \frac1{[J:\QQ]}\sum_{\sigma:J\to\CC}\log U_\sigma .
\]
\end{theorem}
\begin{proof}
For a fixed \(p\) evaluate the original signed integral with the
representatives of Theorem~\ref{thm:neutral315}. This keeps the
period, its volumes, its restriction and its fixed projected
endpoint. In each edge choose the localized product basis of
Theorem~\ref{thm:adapted314}, common to the whole list of tests
and classes. The common phase theorem,
Theorem~\ref{thm:framechange314}, shows that this preserves the
Gram product exactly. An auxiliary finite algebraic extension used
to evaluate a particular representative or phase causes no change
in the normalized valuation of this original element.

By Theorem~\ref{thm:endpoint315} the symbolic coefficients are
integral over \(p\) after multiplication by \(B_{\rm end}\).
The compact Haar averages are coefficientwise, so
Theorem~\ref{thm:compact313} and the unchanged stabilizer orders
give \(D_{\rm intr}\) for each copy of the signed integral.
In every term there is one physical derivative of the prescribed
profile per edge. If \(p\in S\), the finite coefficients are
integral there, and all the shifts have order prime to \(p\)
at this place. The monomial estimate in
Theorem~\ref{thm:finitecolumn313} gives \(g_{i,p}\) for each
factor. This holds for all localizations over \(p\), hence also
for the conjugated factors.

For \(p\notin S\), the two finite coefficient losses on edge
\(i\) sum to at most \(2v_p(T_i^0)\) by
Theorem~\ref{thm:neutral315}. Use Theorem~\ref{thm:splitinput315}
at a split prime and the ordinary product derivative estimate
of Theorem~\ref{thm:classden312} at a nonsplit prime. Together
these give at most \(2H_{i,p}\) for translations and the marked
elliptic factors. The physical coordinate costs of the two
orientations sum to \(d_i\rho_p\), since
\(\mathfrak b\bar{\mathfrak b}=\mathfrak D_{E/K}\), exactly as
in Theorem~\ref{thm:gramheight314}. Products and sums now give
the asserted local bound. We may make this calculation separately
at each \(p\) on the same algebraic Gram element.

For the global sum, use the discriminant identity and the sum of
\(\Phi_p(d_i,v_p(Q_i^0))\log p\) in the proof of
Theorem~\ref{thm:gramheight314}. At split \(p\) that summand
had been \(d_i v_p(Q_i^0)\log p\). Its replacement by
\(H_{i,p}\log p\) gives precisely the displayed correction,
whose summands vanish outside \(Q_i^0\). Finally apply the
product formula to \(D_{\rm intr}^2\mathcal A_a
\overline{\mathcal A_a}\) for a nonzero entry. Conjugation
permutes its coefficient field embeddings, so the archimedean
mean for the two period factors is at most twice the mean of
\(\log U_\sigma\). At \(p\in S\) each \(J^0_{i,p}\) vanishes.
Division by two gives the final inequality.
\end{proof}

\subsection{Split position tests in the designated signed period}\label{nav:IX-66}
On the interval and in the notation of
Theorem~\ref{thm:originalgram315}, fix any finite set
\(\mathcal T\) of odd rational primes split in the marked field
\(K\). The extension \(E=FK\) splits over every \(F_v\) above
these primes, including when \(F_v/\QQ_p\) is ramified.
At each \(p\in\mathcal T\) use the actual finite input lists
of every edge, with the original right translates. Compute
\(J_{i,p}^{\rm pos}\) of Theorem~\ref{thm:poscost316}. Write
\[
     \Delta_i=\sum_{p\in\mathcal T}
               [J^0_{i,p}-J^{\rm pos}_{i,p}]_+\log p,
          \qquad [u]_+=\max(u,0).
\]

\begin{theorem}[A support bound for the signed period and the detected minor]\leavevmode
\label{thm:poscolumn316}
For the Gram products of the entire signed definite column, the
local bounds of Theorem~\ref{thm:originalgram315} hold with
\(\min(J^0_{i,p},J^{\rm pos}_{i,p})\) on each edge at
\(p\in\mathcal T\). Every nonzero entry satisfies
\[
 G_S\le\log D_{\rm intr}+\sum_i(B_i^0-\Delta_i)
                   +\frac{1}{[J:\QQ]}\sum_\sigma\log U_\sigma .
\]
For the nonzero last-edge minor selected in
Theorem~\ref{thm:intrinsicminor315}, use that theorem's
single-edge lists and write \(\Delta\) for their correction.
With its original archimedean lower bound, the very same
minor gives
\[
 \sum_{p\in S}g_p\log p
 \le B^0-\Delta+\tfrac12\log j+\log(i_-i_+b_-b_+)
           +\tfrac d2\log(ns)+\frac{1}{[J:\QQ]}\sum_\sigma\log B_\sigma .
\]
\end{theorem}
\begin{proof}
Fix \(p\in\mathcal T\). Use the representatives of
Theorem~\ref{thm:neutral315} for all the vertex classes. Choose
on each edge whichever of the bases of Theorems
\ref{thm:adapted314} and~\ref{thm:posframe316} gives the
smaller local loss. This is one choice for all the terms and
all the class tuples in the signed column. In the first case
the proof of Theorem~\ref{thm:originalgram315} supplies the
two derivative bounds. In the second,
Theorem~\ref{thm:poscost316} supplies them. When multiplying
the two signed sums, the derivative from an edge on one side
can be paired with the derivative from that edge on the other
side, since the directed towers and block degrees agree.
Both valuation estimates apply to every such pair of terms.
The external coefficients, compact averages and stabilizer
weights are cleared by \(D_{\rm intr}\) in each copy, exactly
as in Theorem~\ref{thm:originalgram315}.

Changing the basis multiplies every period entry by a common
algebraic scalar of complex absolute value one by
Theorem~\ref{thm:framechange314}. Thus the local calculation
applies to the original Gram product. It can be repeated at
another \(p\) using a different common set of edge bases.
The primes in \(S\) are inert, so all of their original
Frobenius bounds still apply. The product formula for a nonzero
diagonal entry now replaces the old sum of local losses on
edge \(i\) by \(B_i^0-\Delta_i\). The archimedean upper
bounds can be taken in the original calculation.

For the minor keep every row and column chosen by
Theorem~\ref{thm:intrinsicminor315}. Thus its nonzero lower
bound and the terminal test detecting the prescribed incoming
vector when \(j=1\) are fixed. Each entry of the normalized
matrix has integral local endpoint tensors. Each term in the
product of its determinant with the conjugate determinant uses
\(j\) derivatives in each orientation. Apply the same change of
basis on the one last edge to every matrix entry. Its common
scalar is raised to power \(j\) in the determinant and cancels
in the product. The lower bound at \(p\in\mathcal T\) is
therefore \(-2j\min(J^0_p,J^{\rm pos}_p)\). All other terms
in the product formula and the division by \(2j\) are those
of Theorem~\ref{thm:intrinsicminor315}. This proves the
second inequality.
\end{proof}

For the unit-amplitude congruence test of
Theorem~\ref{thm:poscost316}, with the other local factors
spherical and no other coefficient loss, the improvement over
the raw Weyl bound is \(\frac12hns\log p\) per edge. This
calculation uses the projection and the fields already fixed in
Section~\ref{sec:missing}. The comparison with the full
prescribed Bessel pairing and the residual strict cohomology
orders enter separately in the interval inequality.

\subsection{Exterior traces of the last projected theta operator}\label{nav:IX-67}
At the last definite edge use the normalization \(z(U,V)\)
of Theorem~\ref{thm:lastrow311}. Take the complete chart
tests of Theorem~\ref{thm:charttrace320} on both sides.
Let \(Z_*\) be their matrix of \(z(U,V)\), with row and
column in the marked conjugate spaces when required.
Write \(h_\pm,i_\pm,w_\pm,D_\pm\) for the numbers of
charts, level indices, chart volumes and compact type
dimensions there. Put \(s_\pm=i_\pm^2 w_\pm/D_\pm\) and
\[
  |\kappa_*|={w_-w_+|C|c^{-d/2}\over D_-D_+},\qquad
  E_j=\operatorname{tr}\big(\textstyle\bigwedge^j(Z_*Z_*^*)\big).
\]
At each real place \(\tau\) let \(L_\tau=|\lambda_\tau|\)
be the last partition degree. For \(R=n,s\) set
\[
 q_{R,\tau}=
 \frac{D(R,L_\tau)}
 {\gcd(D(R,L_\tau),\dim V_{R,\tau})},\qquad
 q=\prod_\tau\gcd(q_{n,\tau},q_{s,\tau}).
\]
Here \(D(R,L)\) is the integer in
Theorem~\ref{thm:compact313}, and \(V_{R,\tau}\) is the
single-place Cartan component in the marked exterior
tensor. Keep the last-edge \(S,g_p,J_p,B^0-\Delta\)
of Theorem~\ref{thm:poscolumn316} and the transported
Fischer bounds \(B_\sigma\) of Theorem~\ref{thm:bothlast311}.
Write \(k_0=\dim\mathcal U\) for the dimension of the
actual incoming subspace with squared norm lower bound \(a_*\).

\begin{theorem}[One compact denominator for an exterior trace]\leavevmode
\label{thm:lasttrace320}
For \(1\le j\le k_0\) the physical trace \(E_j\) is
algebraic and satisfies, at the original embedding,
\[
 E_j\ge {k_0\choose j}
      \left({s_-s_+a_*\over|\kappa_*|^2}\right)^{\!j}>0 .
\]
At every \(\mathfrak p\mid p\),
\[
   v_{\mathfrak p}(q^j E_j)\ge
   \begin{cases}2j g_p,&p\in S,\\ -2jJ_p,&p\notin S.\end{cases}
\]
Consequently on the same fields and intervals
\[
 \sum_{p\in S}g_p\log p\le
 B^0-\Delta+\tfrac12\log q+
 \log(i_-i_+\sqrt{h_-h_+})+
 {1\over[J:\QQ]}\sum_\sigma\log B_\sigma-
                         {\log(j!)\over 2j}.
\]
\end{theorem}
\begin{proof}
By the chart identity, the positive eigenvalues of
\(Z_*Z_*^*\) are \(s_-s_+/|\kappa_*|^2\) times the
positive eigenvalues of the output Gram of the physical
last operator. Its restriction to \(\mathcal U\) has
squared norm bounded below by \(a_*\) on unit vectors.
The argument by orthogonal complements in
Theorem~\ref{thm:closedperiod318} gives at least \(k_0\)
eigenvalues this large. Their elementary symmetric sum
gives the first assertion.

In a standard chart of rank \(R\) let \(P_{R,\tau}\) denote
orthogonal projection from the full exterior tensor at
\(\tau\) to its Cartan component. With \(t\) its unit
highest tensor, Schur orthogonality gives
\[
 P_{R,\tau}=(\dim V_{R,\tau})
          \int_{U(R)}(\rho(g)t)(\rho(g)t)^*\,dg .
\]
Use probability Haar. Indeed the average is scalar on the
cyclic irreducible span, has trace one there, and vanishes
on its orthogonal complement. In standard wedge coordinates
the integrand is integral polynomial of degree at most
\(L_\tau\) in each of \(g,\bar g\). Unitary determinant
characters cancel in the outer product. The compact moment
formula of Theorem~\ref{thm:compact313} gives rational
entries with \(q_{R,\tau}P_{R,\tau}\) integral. In these
coordinates the projectors are real symmetric because
their entries are rational and they are Hermitian.

For each class pair form the auxiliary matrix \(W\) by
using the raw product of matching coordinate minors on
the two ambient tensors and evaluating its derivative
on the actual translated last-edge section. Sum the
original signed level translates on both sides. The
individual polynomials here are the map \(M\) of
Proposition~\ref{prop:cartan244} on the full tensor.
Let \(P_-,P_+\) be the products over real places of the
Cartan projectors, acting on each respective class.
Theorem~\ref{thm:last244} gives \(Z_*=P_+WP_-\), since
the class coefficient of an arbitrary exterior tensor
depends only on its component in the compact cyclic span.
More precisely at any one \(\tau\),
\[
             P_{+,\tau}W=WP_{-,\tau}.
\]
To verify this, \(M\) is equivariant by Cauchy--Binet.
The Cauchy decomposition of the matrix polynomials pairs
the partition \(\lambda_\tau\) on one side only with that
same partition on the other. Each exterior tensor under
consideration contains that highest weight exactly once.
Indeed for a diagonal weight with strictly decreasing
entries the unique maximal tensor line is the product
of the leading wedges. Thus pairing a Cartan factor at
one end with its orthogonal complement at the other gives
zero for every matrix polynomial evaluation. This proves
the identity at each class pair. Finite translates preserve
it since the Fourier test changes in the finite section
and compact coordinate changes preserve the Cartan types.

Fix \(p\) and choose neutral representatives as in
Theorem~\ref{thm:neutral315}. At each \(\tau\) use whichever
end has smaller \(v_p(q_{R,\tau})\). Write \(P_L,P_R'\)
for the products of the chosen output and input projectors
respectively, using the identity for an empty product.
Then \(Z_*=P_L W P_R'\), with one projector per real place.
Taking exterior powers and using the Hilbert--Schmidt
inner product on rectangular matrices gives
\[
 E_j=\|\textstyle\bigwedge^j Z_*\|_{\rm HS}^2
     =\big\langle\textstyle\bigwedge^j W,
                         \textstyle\bigwedge^j Z_*\big\rangle_{\rm HS}.
\]
The equality holds since multiplication on the two sides
by orthogonal projectors is itself an orthogonal projector
on the matrix space, also after taking exterior powers.
In the last expression just one copy of the matrix of
minors contains compact projectors. Expand it by
Cauchy--Binet. A minor of each projector product is a
sum of products of \(j\) entries, so their combined
denominator at \(p\) divides \(q^j\) locally.

Each remaining term contains \(j\) raw derivative entries
of \(W\) and \(j\) conjugate entries. At the neutral charts
their polynomials are products of coordinate minors with
integer coefficients. Every full-height minor uses every
column, and the block degrees are \(L_\tau\), including
before a Cartan projector is applied. Thus the monomial
valuation calculation in Theorems~\ref{thm:frob244}
and~\ref{thm:intrinsicminor315} gives \(2g_p\) per pair
in \(S\) and at least \(-2J_p\) elsewhere, retaining the
position bound of Theorem~\ref{thm:poscolumn316}. No compact
average is left in \(W\). The trivial stabilizers and the
normalization of \(z\) remove the chart masses and the two
Schur dimension factors. These computations prove the
finite-place assertion. The auxiliary matrix can be chosen
in new charts for each prime because the first part
expresses \(E_j\) exactly as a fixed multiple of the
physical exterior trace.

At an embedding \(\sigma\) evaluate the actual matrix
\(Z_*\) in the original standard chart coordinates. On one
class pair and one translate, its Hilbert--Schmidt norm is
at most \(B_\sigma\). In fact in each block the map from
the two Cartan factors to matrix polynomials is an isometry
times \(\sqrt{h_{\lambda_\tau}}\) by
Proposition~\ref{prop:cartan244}. On the ambient product
the actual pairing first projects both factors
orthogonally onto these components. The adjoint of the
resulting map has norm at most \(\sqrt{h_\lambda}\), the
product of the single-place norm factors. Applying it to
the Fischer representative of the Taylor functional bounds
the Euclidean norm of the entire array by
\(\sqrt{h_\lambda}\) times the Fischer norm. This is exactly
the definition of \(B_\sigma\), including its finite vacuum
bound. Under transport the blocks are permuted or conjugated
with their indicated types, as in
Theorem~\ref{thm:bothlast311}. Summing at most \(i_-i_+\)
translates and then summing squares over \(h_-h_+\) class
pairs gives \(\|\sigma Z_*\|_{\rm HS}\le
A_\sigma=i_-i_+\sqrt{h_-h_+}B_\sigma\).

For any matrix of Hilbert--Schmidt norm \(A\), the sum of
squared \(j\)-minors is at most \(A^{2j}/j!\), by expanding
the \(j\)-th power of the sum of its squared singular
values. Therefore, with \(c_J\) the original conjugation
of the coefficient field,
\[
       |\sigma E_j|\le { A_\sigma^{\,j} A_{\sigma\circ c_J}^{\,j}
                             \over j!}.
\]
This is Cauchy--Schwarz on the two transported lists of
minors, so it also applies to embeddings that do not
commute with conjugation. Apply the product formula to
the nonzero \(q^j E_j\), use
\(\sum_{p\notin S}J_p\log p=B^0-\Delta\), and divide
by \(2j\). This yields the asserted inequality.
\end{proof}

\subsection{A signed integral over sets of positive local measure}\label{nav:IX-68}
Use the supported local path and all compact middle averages
of Theorems~\ref{thm:splitpath322} and~\ref{thm:splitcompact323}.
We integrate its core position matrices \(X_j\), its top
matrix \(B\), and its last frequency \(z\) as additional
variables. Keep the initial endpoint fixed. Matrix measure
is product additive Haar, with the integral matrix lattice
of each shape of volume one. At the frequency use the dual
measure, so the integral dual lattice also has volume one.
Put \(d=n_{s+1}\), \(D_+=r+(s+2)t\) for the row dimension
of \(B\), and \(P_a=\prod_{i=1}^a(1-q^{-i})\). Use \(P_0=1\).
Retain the oriented integral covector \(l_0\) used in
Theorem~\ref{thm:splitpath322}, chosen as a marked coordinate
functional. Put \(U_0=\ker(\bar l_0|_U)\).

\begin{theorem}[A common integral for the tame twists]\leavevmode
\label{thm:splitmass323}
There is a compact open set of the integral position
matrices, independent of the twist and the final residue
unit \(b\), for which the signed path integral over this
set and over \(\pi^{-s-1}\widetilde b\,l_0+O^{1\times d}\)
equals
\[
 {\cal V}\,\eta\,\chi(b)^{-1}\,
 q^{\,r-1-(s+1)d}\,G(\chi)^{\,1+s(s+1)/2}.
\]
Here \(\widetilde b\in O^\times\) is any fixed lift of \(b\);
the position set is the same for every lift. Its exact
measure is
\[
 \begin{split}
 {\cal V}={}&
 \frac{P_r}{q-1}\,q^{-r(D_+-r)}
   P_{a_s}\,q^{-s-(d-a_s)a_s}\\
 &{}\times\prod_{j=0}^{s-1}
  P_{a_j}\,q^{-j n_{j+1}-(n_{j+1}-a_{j+1})a_j}.
 \end{split}
\]
In particular a single determinant condition suffices for
all the marked factors together. At any tuple of distinct
split twist places on the common field the product of these
sets has the product measure and the integral of the product
of the local paths has the product value.
\end{theorem}
\begin{proof}
For \(j<s\) impose the following conditions. At \(j>0\)
require \(X_j z_j\equiv z_{j+1}\pmod{\pi^j}\) and use
\(v_j=(X_jz_j-z_{j+1})/\pi^j\). At \(j=0\) take
\(v_0=-z_1\). In the residue target core the marked
columns of \(X_j\) together with \(v_j\) must be a basis of
the marked plane \(A_{j+1}\). All unmarked columns may
be arbitrary integral columns. Let \(u_j\) be the
determinant of this basis with its original ordering,
including the fixed sign from the flag coordinates.

At \(j>0\) the initial congruence is prescribed on one
column, so has measure \(q^{-j n_{j+1}}\).
Conditionally \(v_j\) is uniform in the integral target
core and independent of the other columns. In the marked
residue determinant there is already one fixed vector
\(z_{j+1}\). The other \(a_j\) random columns are \(v_j\)
and the \(a_j-1\) other marked source columns. At \(j=0\)
the fixed vector is \(v_0\) and the \(a_0\) random columns
are the marked source columns themselves. In either case
their unmarked entries must vanish modulo \(\pi\), with
probability \(q^{-(n_{j+1}-a_{j+1})a_j}\). In the marked
target modulo the line of the fixed vector their images
are independent uniform vectors in \(K^{a_j}\).
The probability that they form a basis is \(P_{a_j}\),
by choosing the vectors successively. These are precisely
the factors in the product in the second line of the
volume formula. On this set the calculation of the
interior entries in Theorem~\ref{thm:splitpath322}
applies pointwise with the unit \(u_j\).

At \(j=s\) require the \(a_s=r-1\) marked columns of
\(X_s\) to reduce to a basis of \(U_0\). This has
probability \(q^{-(d-a_s)a_s}P_{a_s}\). Impose also
\(l_0X_sz_s\in\pi^{s+1}O\). For \(s\ge1\) the tested
entry is the transverse coordinate of one marked column.
Its residue is already zero, and its next \(s\) digits are
unrestricted and uniform. This contributes \(q^{-s}\).
For \(s=0\) the condition is automatic because \(z_0=0\).
Denote by \(u_s\in K^\times\) the scalar for which the
last-column determinant form is \(u_s\,\bar l_0|_U\).
It is well defined also when \(r=1\). Substitution of
\(y=-X_sz_s+\pi^s v\) gives additive phase one on the
entire frequency coset in the statement. Indeed its center
pairs integrally with \(X_sz_s\) by the new condition, and
the additional integral frequency does so by integrality.
On \(v\) the residual phase has parameter \(b/u_s\) in
the determinant coordinate. The Gauss sum therefore
contributes \(\chi(u_s)\chi(b)^{-1}G(\chi)\), with all
volume powers as before.

For \(B\) require the \(r\) marked columns to give an
isomorphism of the marked residue planes. This has measure
\(q^{-r(D_+-r)}P_r\). Its determinant \(u_+\) is uniform
in \(K^\times\) on this event: rescaling one marked column
permutes the possible determinants transitively, preserving
the measure. Impose the single final condition
\[
                u_+\prod_{j=0}^s u_j=1 .
\]
Given all the core matrices, exactly \(1/(q-1)\) of
the admissible top matrices meet it. The product of the
residue character values at the positions is then one
for every \(\chi\). All these tests are congruences in
compact integral matrix lattices, so their joint locus
is compact open. Multiplication of the calculated measures
gives \({\cal V}>0\).

The hypotheses used in the compact averaging proof are
valid on the full set. The interior integrals have the
same marked determinant ranks, and the outgoing covector
and marked columns satisfy the last-column calculation.
Every contributing \(y\) is still primitive with nonzero
reduction in \(U\). Thus the normalized terminal orbit
sum and the other compact middle integrals give the
pointwise signed product above throughout this set.
The frequency coset has volume one. Integration now
gives the displayed value, with the unramified scalar
\(\eta\) retained. Finite products over places follow
by product Haar and Fubini on compact sets.
\end{proof}

This integral isolates specified local positions in the
split auxiliary factors. In the designated global period
the other interior paths, the rational class sum and the
global cuspidal projection still enter with their signs.
Controlling their combined contribution on the same
interval is necessary to deduce the claimed GRH estimate.

\subsection{The phase of the added positive line}\label{nav:V-55}
At one real place of a transverse pair with \(\mathcal G>0\)
put
\[
 L_\tau=|\tau D_d|-T|\tau D_s|>0,\qquad
 \Delta_+=D_d+T D_s,\qquad \kappa=k-d,\quad E_0=\kappa+2ds+s^2 .
\]
Use the extended monomial frame and \(\delta_p\) of
Proposition~\ref{prop:extend300}. For each endpoint \(J_i<d\)
put \(a_i=D_{J_i}/Q^{J_i}\). At the last endpoint put
\(a_t=\Delta_+/Q^d\). These are nonzero algebraic integers by
Proposition~\ref{prop:shift203}. The index \(t\) here counts
blocks. Write \(u_p=d-J_{i(p)}\) as in Theorem~\ref{thm:realtrans300}.

\begin{proposition}[Decay in the difference of the two terminal moduli]\leavevmode
\label{prop:phase301}
For nonnegative endpoint powers \(\eta_i\) with
\(\ell=\sum_i J_i\eta_i\le r^2\), the original real integral
is at most
\[
 \frac{e^{Ckd+C d\ell\log b}}{(k!)^s T^{s(k+d)}}
 \left(\frac{2\,\tau\mathcal G}{L_\tau^2}\right)^{-\kappa}
 \left(\frac{L_\tau}{(2T)^d}\right)^{-E_0}
 \prod_{p=1}^d\delta_p^{-2u_p}
 \prod_i\max(1,|\tau a_i|)^{-\eta_i}.
\]
The constant is uniform in the phase, in the words and in
the real place.
\end{proposition}
\begin{proof}
In Proposition~\ref{prop:extend300} multiply the positive
component of the second terminal vector by \(\zeta\) of
modulus one. Equivalently extend the relative isometry \(w\)
by multiplication by \(\zeta\) on the extra line, and extend
each member of the original Bessel group by the identity.
The new cross determinant is
\(\Delta_\zeta=\tau D_d+\zeta T\tau D_s\). Choose \(\zeta\)
so that \(|\Delta_\zeta|=L_\tau\), taking it to be one if
\(D_s=0\). If the last block is a singleton only its pivot
changes. Otherwise the column operation in
Proposition~\ref{prop:extend300}, with \(T\) replaced by
\(\zeta T\), restores the original monomial matrix.
Every factor \(\delta_p\) with \(u_p>0\) is consequently
unchanged.

The two pairings of the positive vector against the physical
frames are now \(\bar\varpi{\bf e}_d\) and
\(\bar\zeta\bar\varpi{\bf e}_d\) at this embedding.
The projection calculation of Proposition~\ref{prop:extend300}
gives perpendicular norm
\[
    1-2T\operatorname{Re}(\zeta\tau D_s/\Delta_\zeta)
                    =\tau\mathcal G/|\Delta_\zeta|^2 .
\]
In Proposition~\ref{prop:slice300} the equation for the
primal block of the fixed vector is now
\(C^*a=\bar\zeta{\bf e}_d/\sqrt2\). It is still supported
on the first coordinate of the last block. The stabilizer,
the Haar computation, and the Gram integral therefore apply,
with sphere radius squared \(2\tau\mathcal G/L_\tau^2\).
For an original group element, a phase on the extra line leaves
the norm of the extended terminal exterior vector unchanged.
Indeed the two terms from its \(e\) and extra-line components
are orthogonal. It also preserves the extended operator norm.
This justifies the same substitution for the reciprocal
squared integral, where the two frames interchange and the
phase is conjugated. The change of frequency on the original
group is the one in Theorem~\ref{thm:realtrans300}.

It remains to justify the endpoint tests and their larger powers.
On the shell of that proof set \(M=\nu(g)H\). The physical
minors in the unit-frequency frames with the original positive
line phase satisfy
\[
        \left(\frac{Q}{2T}\right)^{J_i}|\tau a_i|\le M^{J_i}
                  \qquad (1\le i\le t).
\]
For the final one use the unrotated exterior pairing of the
two extended frames. Its value is unchanged by any left or
right change in the original Bessel group. After choosing the
minimal representatives of \(g\) and \(y\) in the shell, their
product with one inverted represents the same double coset
and has norm at most \(M\). Extending it by the identity on
the extra line preserves that bound. Pair the unit exterior
frames to obtain the assertion, including on the domain of
integration used for the rotated estimate. The proper endpoint
tests follow by the same argument with the shorter wedges.

For each \(|\tau a_i|\ge1\) insert the resulting inverse
power and its compensating factor \(M^{J_i\eta_i}\).
The additional constant is at most
\(\max(1,2T/Q)^\ell\le\exp(Cd\ell\log b)\).
In the proof of Theorem~\ref{thm:realtrans300} the shell
exponent becomes at most \(K_0+\ell\), where \(K_0=O(d^3)\).
The normalized diagonal moment after shell integration has
order at most \(2(K_0+\ell)+2\le4r^2\) for large \(x\).
Theorem~\ref{thm:moment286} bounds its local logarithm by
\(C s(K_0+\ell+1)\log b\).
The part \(O(d^4\log b)\) is \(O(kd)\) on the common
intervals. This proves the assertion.
\end{proof}

\subsection{The determinant of the difference of two pairings}\label{nav:V-56}
Let the two frames \(F_i\) of a transverse cell be column
matrices with \(d=s+1\) columns each. Write \(Y_i\) for their
spans, \(W=Y_0\oplus Y_1\), \(C=W^\perp\) and
\(Z=F_0^*F_1\), where the star uses the row Hermitian form.
We take it conjugate linear in the first variable. Write
\(D_d=\det Z\), \(D_s\) for the leading \(s\)-minor and
\(\mathcal G=D_d\bar D_d-T^2D_s\bar D_s>0\) at the real
places. Let \(\pi_i:Y_i^\perp\to P_i=Y_i^\perp/X_i\)
be the quotient map, with \(p_i=\pi_i|_C\).
Use the unitary identification \(\beta:P_0\to P_1\) of
Theorem~\ref{thm:rationalrank302}. Set
\[
 C_0=\ker(p_1-\beta p_0),\qquad C_1=C\ominus C_0,\qquad
 h=\dim_E C_1,\qquad W'=W\oplus C_1 .
\]
Define a pairing on \(Y_0^\perp\times Y_1^\perp\) by
\[
 {\cal H}(v,w)=\langle v,w\rangle-
                        \langle\beta\pi_0 v,\pi_1 w\rangle .
\]
Either argument in \(C_0\) gives zero. Indeed pairing a row
vector orthogonal to \(Y_i\) against \(c\in C\) is the same
as pairing its column image against \(p_i c\). Thus we may
restrict \({\cal H}\) to \(S_i=W'\cap Y_i^\perp\) on the two sides.

\begin{lemma}[Archimedean determinant of the comparison pairing]\leavevmode
\label{lem:archdet303}
At a real place \(\tau\) give \(S_0\) the coordinates
\((F_1^*v,\operatorname{pr}_{C_1}v)\) and \(S_1\) the coordinates
\((F_0^*w,\operatorname{pr}_{C_1}w)\). Use orthonormal
coordinates for \(C_1\). These identify each space with
\(\CC^d\oplus\CC^h\). The matrix \(N_\tau\) of \({\cal H}\)
in these coordinates satisfies
\[
     0<|\det N_\tau|\le
          \frac{2^h}{|\tau D_d|-T|\tau D_s|}.
\]
\end{lemma}
\begin{proof}
Suppress \(\tau\) in the computation. The frame Gram matrix on
\(W\) is \(M=\bigl(\begin{smallmatrix}J&Z\\ Z^*&J\end{smallmatrix}\bigr)\),
where \(J=-T e_de_d^*\). Its absolute determinant is \(\mathcal G\),
as follows by expanding in the two rank-one diagonal blocks.
Both \(M\) and \(Z\) are invertible. In dual frame coordinates
\((F_0^*v,F_1^*v)\) its pairing matrix is \(M^{-1}\), proving
the coordinate assertion. Write
\[
            M^{-1}=\begin{pmatrix}B_1&A^*\\ A&B_0\end{pmatrix}.
\]
The complementary minor identity gives
\(|\det A|=|\det Z|/\mathcal G\). Equivalently, reorder the
block rows of \(M\) to put \(Z^*\) in its leading corner
and take the inverse using its Schur complement.
The matrices \(B_i\) are the forms on \(W\cap Y_i^\perp\).
Each is the square norm of the projection to the positive
column line orthogonal to \(p_i C\), with radical \(X_i\).
Choose unit coordinates on those lines and write
\(B_i=\xi_i\xi_i^*\), so the projected coordinate of a
vector with dual frame coordinate \(q\) is \(\xi_i^*q\).
The block equations \(MM^{-1}=1\) give
\[
      B_1 A^{-1}=-(Z^*)^{-1}J,\qquad
      B_0(A^*)^{-1}=-Z^{-1}J .
\]
It follows by taking the trace of their product that
\[
      |\xi_1^* A^{-1}\xi_0|^2
       =T^2|e_d^* Z^{-1} e_d|^2=T^2|D_s/D_d|^2 .
\]

The comparison identifies the copies of \(C_0\) pointwise
and maps their orthogonal complements in the two columns
isometrically to each other. On these complements use the
one line coordinate followed by the \(h\) coordinates of
\(p_i C_1\). Let \(U\) be the matrix of the inverse
comparison from the side-one coordinates to the side-zero
coordinates. With \(J_i=\operatorname{diag}(\xi_i^*,I_h)\)
the pairing matrix is
\[
        N_\tau=\operatorname{diag}(A,I_h)-J_0^*U J_1 .
\]
Use \(\det(1-AB)=\det(1-BA)\) and put \(z=\xi_1^*A^{-1}\xi_0\).
Then
\[
      \det N_\tau=\det A\,\det(1-U\operatorname{diag}(z,I_h)).
\]
Here \(|z|<1\). A fixed vector of
\(U\operatorname{diag}(z,I_h)\) would have zero first
coordinate by equality of norms. Its remaining coordinates
would then give an element of \(C_1\) on which the two
column maps agree. Formation of the kernel commutes with
extension to the complex embedding, so that element is zero.
For the upper bound the first column of
\(1-U\operatorname{diag}(z,I_h)\) has length at most
\(1+|z|\), and each other column has length at most two.
Hadamard's inequality gives the assertion because
\(|\det A|(1+|z|)=1/(|D_d|-T|D_s|)\).
\end{proof}

\subsection{Bounds for projected periods from Hecke eigenvalues}\label{nav:III-75}
Let \(S\) be the terminal table over the coefficient ring
of integers from Theorem~\ref{thm:classtrace326}. Take a
finite list of commuting integral Hecke operators \(H_a\)
which act semisimply over a common splitting field \(K\).
Extend the integral lattice to its ring of integers.
Index the distinct simultaneous eigenvalue lists by \(i\).
Write \(V_i\) for the full associated space, of dimension
\(d_i\), and \(\alpha_{ia}\) for the eigenvalues. For
\(i\ne j\) define the nonzero integral ideal
\[
             \mathfrak d_{ij}=(\,\alpha_{ia}-\alpha_{ja}:a\,).
\]
For a subset \(I\) of the lists, take \(P\) to be the
projection onto \(\bigoplus_{i\in I}V_i\) along the sum
of the remaining spaces. It is defined over \(K\) by the
given direct decomposition. For commuting normal operators
this is the orthogonal projection.

\begin{theorem}[Simultaneous exterior bounds from Hecke congruences]\leavevmode
\label{thm:heckeproj326}
For every prime ideal \(\mathfrak p\) of \(K\), let
\(c_1\ge\cdots\ge c_d\) be the local exponents of
Theorem~\ref{thm:projclass326}. Then
\[
   \sum_{\ell=1}^k c_\ell\ \le\
   \sum_{\substack{i\in I\\ j\notin I}}
          \min(k,d_i,d_j)\,v_{\mathfrak p}(\mathfrak d_{ij})
               \qquad (1\le k\le d).
\]
In particular the normalized projected tests are perfect
at each prime where the two sets of eigenvalue lists remain
disjoint on reduction. Globally the corresponding sum of
losses weighted by \(\log\mathrm N\mathfrak p\) is bounded
by the same sum with \(\log\mathrm N\mathfrak d_{ij}\) in
place of \(v_{\mathfrak p}(\mathfrak d_{ij})\).
\end{theorem}
\begin{proof}
Work over the local discrete valuation ring and set
\(Q=1-P\). The finite module \(C=S/(S_P+S_Q)\) is a
quotient of each of \(P S\) and \(Q S\). For the latter
assertion send \(Qs\) to the class of \(s\) in \(C\).
Changing \(s\) by \(S_P\) preserves the class. The former
assertion follows with \(P,Q\) exchanged. Both maps are
surjective and Hecke equivariant.

Order the \(i\in I\). Filter \(P S\) by intersecting it
with the successive partial sums of the \(V_i\). Each
successive quotient is free of rank \(d_i\), since projection
onto \(V_i\) gives a full lattice there with kernel the
preceding intersection. The action of \(H_a\) there is the scalar
\(\alpha_{ia}\). The images in \(C\) give a filtration
\(A_i\) whose factors are generated by at most \(d_i\)
elements, with the same scalar action. Construct \(B_j\)
from \(Q S\) using the complementary list. Denote the
preceding terms by \(A_{i^-}\) and \(B_{j^-}\), starting
with zero.

In \(A_i/A_{i^-}\) take the filtration induced by
\(A_i\cap B_j\). Its successive factors are
\[
 \frac{A_i\cap B_j}
      {(A_{i^-}\cap B_j)+(A_i\cap B_{j^-})}.
\]
Indeed, an element in the kernel differs from a member
of \(A_i\cap B_{j^-}\) by an element of \(A_{i^-}\)
which is still in \(B_j\). The formula is symmetric in
the two filtrations. Each displayed module is therefore a
subquotient of factors on both sides. It has at most
\(\min(d_i,d_j)\) generators over the discrete valuation
ring. Every \(\alpha_{ia}-\alpha_{ja}\) kills it, since
\(H_a\) acts by each of the two scalars.

Take the sum of the \(k\) largest cyclic factors in any
elementary divisor decomposition of \(C\) and intersect it
with the refined filtration just obtained. Its factors are
subquotients of the displayed modules and of a \(k\)-generated
module. Over a discrete valuation ring a subquotient of a
module on \(k\) generators has at most \(k\) generators,
by the elementary divisor theorem. The length at \((i,j)\)
is thus at most \(\min(k,d_i,d_j)v_{\mathfrak p}(\mathfrak d_{ij})\).
Adding proves the local inequality. Factoring the norms
of the ideals gives the global one.
\end{proof}

The right side has an archimedean bound for actual class
operators. For each pair choose any \(a\) distinguishing
it. Then \(\mathfrak d_{ij}\) contains
\((\alpha_{ia}-\alpha_{ja})\), so
\(\log\mathrm N\mathfrak d_{ij}\) is bounded by the
sum of \(\log|\sigma(\alpha_{ia}-\alpha_{ja})|\) over all
embeddings. For an unnormalized finite double coset with
unitary character and \(b_a=|JgJ/J|\), its convolution
has norm at most \(b_a\), by the triangle inequality for
unitary translations with level volume divided out. If
multiplication by a positive integer \(m_a\) was used to
preserve the integral lattice, the bound for this operator
is \(m_a b_a\). Such an integer exists for a rational
class operator by clearing its denominators on the finite
lattice. When all conjugates have the stated compact class
realizations, this gives the explicit bound
\([K:\QQ]\log(2m_a b_a)\) for that pair. The estimate
separates whole simultaneous eigenspaces. A further
projection inside an eigenspace still has its own quotient
from Theorem~\ref{thm:projclass326}.

\subsection{Maximal minors at prescribed auxiliary frequencies}\label{sec:periodvolume331}
At a place of \(Q(\xi)\) fix \(l\) and the other data of
Theorem~\ref{thm:signedfibers331}. We compare the complex norm and
the integral image of the same map on the position coordinates
\((b,y)\). Take a nonempty set \({\cal X}\) of distinct nontrivial
characters of \(\mathbb F_p^\times\) and a nonempty set
\({\cal B}\subset(\ZZ/p^s\ZZ)^\times\) of frequencies. Singleton
sets allow either datum to be prescribed. For a weight on the full
positions and frequency variables, formula~\eqref{eq:signedfibers331}
uses a separate fiber mean at each \(\nu\). The linear map considered
here uses a single array on the position fibers as input for every
row.

Choose a generator \(g\) of \(\mathbb F_p^\times\). Put
\(\alpha_\chi=\chi(g)\), \(\zeta=\psi(p^{-s})\), \(z_\nu=\zeta^\nu\),
\(k=|{\cal X}|\), \(a=|{\cal B}|\), \(R=ka\), and \(M=(p-1)p^s\).
Let \(J\) be a cyclotomic number field containing the \(p^s\)-th
and \((p-1)\)-st roots of unity. Divide the actual integral row
by \({\cal V}\eta_\chi p^{r-1-(s+1)d}\). The resulting matrix
over the integers of \(J\) is
\[
 K_{(\chi,\nu),(i,j)}
 =G(\chi)^h\chi(\bar\nu)^{-1}\alpha_\chi^i z_\nu^j,
 \qquad 0\le i<p-1,\quad 0\le j<p^s .
\]
Write \(\mathfrak d_R(K)\) for the ideal generated by all minors
of order \(R\). For fixed orderings set
\[
 \Delta_{\cal X}=\prod_{\chi<\chi'}(\alpha_{\chi'}-\alpha_\chi),
 \qquad
 \Delta_{\cal B}=\prod_{\nu<\nu'}(z_{\nu'}-z_\nu).
\]

\begin{theorem}[Complex volume and integral image of the signed tests]\leavevmode
\label{thm:periodvolume331}
At every embedding \(\sigma:J\hookrightarrow\CC\) the rows of
\(\sigma K\) are orthogonal and have squared norm \(Mp^h\).
The maximal-minor ideal is principal, with generator up to a unit
\[
  \Delta_{\cal X}^{\,a}\Delta_{\cal B}^{\,k}
                 \prod_{\chi\in{\cal X}} G(\chi)^{ha}.
\]
Use the absolute ideal norm and define the difference per row by
\[
 {\cal D}(K)=\frac{1}{R[J:\QQ]}
 \left\{\sum_{\sigma}\log\sqrt{\det\big((\sigma K)(\sigma K)^*\big)}
             -\log N\mathfrak d_R(K)\right\}.
\]
Then
\[
 {\cal D}(K)=\tfrac12\log M
   -\frac{\log|N_{J/\QQ}\Delta_{\cal X}|}{k[J:\QQ]}
   -\frac{\log|N_{J/\QQ}\Delta_{\cal B}|}{a[J:\QQ]} .
\]
On all nontrivial characters and all primitive frequencies this is
\[
 {\cal D}_{\rm all}(p)=
      \tfrac12\left(\frac{\log(p-1)}{p-2}
                         +\frac{\log p}{p-1}\right).
\]
On all nontrivial characters and a single prescribed primitive
frequency its value is
\[
 {\cal D}_{\rm fixed}(p,s)=
      \tfrac12\left(\frac{\log(p-1)}{p-2}+s\log p\right).
\]
\end{theorem}
\begin{proof}
At each embedding, summing \(\alpha_\chi^i\overline{\alpha_{\chi'}^i}\)
and \(z_\nu^j\overline{z_{\nu'}^j}\) independently gives zero
for distinct rows, and gives \(M\) for a repeated row. Automorphisms
of a root-of-unity group preserve primitivity and distinctness.
Every conjugate Gauss sum for a nontrivial \(\chi\) has magnitude
\(\sqrt p\) by Lemma~\ref{lem:gauss89}. This proves the norm formula.

Remove the row scalars temporarily. The remaining matrix is the
tensor product of two rectangular power matrices. For the first,
divide each power of degree at least \(k\) by the monic polynomial
\(\prod_{\chi\in{\cal X}}(T-\alpha_\chi)\). Subtract the evaluations
of the remainder, expressed in the first \(k\) columns. These
integral column operations give the square Vandermonde matrix
followed by zero columns. The same division by
\(\prod_{\nu\in{\cal B}}(T-z_\nu)\) works in the second matrix.
The tensor product of the two column operations is integrally
invertible. After permuting columns, only one square block of size
\(R\) remains. Its determinant is \(\Delta_{\cal X}^{\,a}
\Delta_{\cal B}^{\,k}\) up to a sign, by the determinant formula for
a tensor product. Integral invertible column operations preserve the
ideal of maximal minors by expansion of minors, applied in both
directions. Restoring the row scalars multiplies the ideal by their
product. The factors \(\chi(\bar\nu)^{-1}\) are units, giving
the indicated generator. In particular all coefficient primes,
including those over \(p\), enter this ideal.

The sum of the logarithms of the absolute values of any nonzero
algebraic row scalar at all embeddings is the logarithm of the norm
of its principal fractional ideal. Hence multiplying a row by this
scalar leaves \({\cal D}\) unchanged. Applying this to the Gauss
sums proves the formula for \({\cal D}\). It also permits the
restoration of the rational measure and modulus factors. The
unramified \(\eta_\chi\) remains explicit in the original signed
integral. The same invariance applies to it whenever it is taken
as a nonzero algebraic scalar with all of its conjugates.

For the full character set the \(\alpha_\chi\) are the \(p-2\)
roots of \(A(T)=(T^{p-1}-1)/(T-1)\). At a root \(\alpha\),
\(|A'(\alpha)|=(p-1)/|\alpha-1|\). Multiplying over the roots
and using \(A(1)=p-1\) gives
\(|\Delta_{\cal X}|^2=(p-1)^{p-3}\). For the full frequency set,
\(a=\varphi(p^s)\) and the roots are those of
\(\Phi_{p^s}(T)=(T^{p^s}-1)/(T^{p^{s-1}}-1)\).
At a root \(z\) the absolute value of the derivative is
\(p^s/|z^{p^{s-1}}-1|\). Each nontrivial \(p\)-th root occurs
\(p^{s-1}\) times as \(z^{p^{s-1}}\). Their factors \(|1-z^{p^{s-1}}|\)
therefore have product \(p^{p^{s-1}}\), since \(\Phi_p(1)=p\).
Multiplication of the derivatives gives
\(|\Delta_{\cal B}|^2=p^{a(s-1/(p-1))}\). These absolute values
are the same at each embedding. Substitution proves the full formula.
For a single frequency \(\Delta_{\cal B}=1\), giving the second
formula.
\end{proof}

\begin{corollary}[The full auxiliary factors on the common field]\leavevmode
\label{cor:tuplevolume331}
For the fixed \(x,\xi\) let \(S=Q(\xi)\),
\(f=[E_x^+:\QQ]\), and \(N_S=\prod_{v\in S}p_v\), counting
places separately. Take the tensor product of the full local maps
of Theorem~\ref{thm:periodvolume331} on their joint position
array. Its difference per row satisfies, with an absolute constant,
\[
   \sum_{v\in S}{\cal D}_{\rm all}(p_v)
         \ll f\bigl(1+\log(2+f^{-1}\log N_S)\bigr).
\]
For a single prescribed primitive frequency at each place and all
local nontrivial characters the difference per row is exactly
\[
   \frac12\sum_{v\in S}
     \left(\frac{\log(p_v-1)}{p_v-2}+s\log p_v\right).
\]
Both formulas apply at every \(x\) and every admissible tuple on
the intervals of Section~\ref{sec:missing}. The empty tensor
product has difference zero.
\end{corollary}
\begin{proof}
Take one common cyclotomic field for the finitely many matrices.
Extension of the field leaves normalized logarithmic norms and
archimedean means unchanged. Each map has an integral change of
source basis reducing it to a nonsingular square block followed
by zero columns, as in the proof of the theorem, also after
restoring the row factors. For a tensor product of ranks \(R_1,R_2\)
the determinant of this block is the product of their determinants
raised to powers \(R_2,R_1\), respectively. Its Gram matrix before
the column operations is the tensor product of the two original Gram
matrices. The same determinant rule for these matrices gives
additivity of \({\cal D}\) after division by rank. Iteration proves
both exact sums.

For odd \(p\) the full contribution is \(O(\log p/p)\).
There are at most \(f\) places above each \(p\). For \(Z\ge2\)
the terms with \(p\le Z\) sum to \(O(f\log(2Z))\). Indeed for
\(n=\lfloor Z\rfloor\) one has
\(\sum_{p\le n}\log p\lfloor n/p\rfloor\le \log(n!)\)
and \(\lfloor n/p\rfloor\ge n/(2p)\). The terms with \(p>Z\)
sum to \(O(\log N_S/Z)\). Choosing \(Z=2+f^{-1}\log N_S\)
gives the claimed estimate.
\end{proof}

An input supported on \((1,0)\) has a nonzero value at each
selected row. An input constant in the last digit of \(y\)
at one place vanishes on every primitive row there, by
Theorem~\ref{thm:signedfibers331}. Both facts hold for the same
fields, matrices and norms. For a specified global class weight,
its contribution at a given tuple is the corresponding signed
coefficient of its joint fiber sums. The global cuspidal projection
and Bessel pairing in Section~\ref{sec:missing} require an estimate
for those particular sums, with their actual multiplicities, signs
and splitting factors.


\subsection{Integral values of a projected row under split matching}\label{sec:rowcontent335}

Fix \(x\) in an interval of Section~\ref{sec:missing}. Use one square over
\(E_x/E_x^+\), the same \(T=S_0\cup Q(\xi)\) and the same levels
\(e_v\ge c_v\) as in Corollary~\ref{cor:splitproduct334}. Fix every
test outside these places. Regard the large-side projected class sum
as a linear functional \(\mathcal L\) in the remaining tensor of tests.
For example, on the definite class tables of
Theorem~\ref{thm:classtrace326}, a trace-dual test of class \(i\) gives
the value \(\operatorname{vol}(J)\ell_i(y_i)\). A prescribed rational
projection \(P\) gives the row \(y\mapsto\operatorname{vol}(J)\ell_i((Py)_i)\).
Signed sums of such rows have exactly the same interpretation. Finite
stabilizer weights, projections and the common Haar factor are included
in \(\mathcal L\).

For \(W\in\mathcal G_v\), let \(f_W\) be the covariant large-side frame
function supported on the complementary summand and taking value one
at \(X_W\) in standard coordinates. Define
\(J_{A,v}f(X)\) by first placing \(f\) on the block lattice, as
\(X\mapsto f(XD_v^{-1})\), and then applying the inverse Fourier
transform of Theorem~\ref{thm:mixedfourier334} on \(P_v^c\).
For \(\mathbf W=(W_v)_{v\in T}\) form the finite row
\[
      \lambda_{\mathbf W}=
          \mathcal L\!\left(\bigotimes_{v\in T}J_{A,v} f_{W_v}\right).
\]
This definition retains correlations between different places. Let
\(\mathcal L_{\rm in}\) denote the pullback of the same \(\mathcal L\)
to the original small-side level tests by the mixed transfers.
Its entries are its values on the individual residue-class indicators
on the product of the small matrix spaces. For any finite table \(b\)
over a number field and \(w(\ell)=1\), put
\(m_w(b)=\min_j w(b_j)\), with \(m_w(0)=+\infty\).
It is the valuation of the fractional ideal generated by the entries.

\Needspace{5\baselineskip}
\begin{theorem}[The fractional ideal of tested values]\label{thm:rowcontent335}
Write \(H_v=|\mathrm{GL}_{b_v}(O_v/\pi_v^{e_v})|\),
\(d_v=|P_v|\), \(\sigma_v=(\nu_v(\pi_v)q_v^{-z_v})^{e_v(a_v-d_v)}\),
and
\[
                  B=\left(\prod_{v\in T}H_v/\sigma_v\right)
                             \mathcal L_{\rm in}.
\]
The row \(B\) vanishes exactly when all the \(\lambda_{\mathbf W}\) vanish. For algebraic entries
\(\lambda_{\mathbf W}\), work in a number field containing them and
the local factors and phases. At coefficient primes outside
the residue characteristics of \(T\), one has
\[
                 m_w(B)=m_w((\lambda_{\mathbf W})_{\mathbf W}).
\]
In particular the ideal of values of
\((\prod_{v\in T}H_v)\mathcal L_{\rm in}\) on integral tests equals
the ideal of the restricted row after inverting those residue
characteristics. At a coefficient prime over \(\ell\), if the row is
nonzero, write \(m=m_w((\lambda_{\mathbf W})_{\mathbf W})\) and
\(f_v=\log_\ell q_v\) for \(v\mid\ell\). Then
\[
 m-\sum_{\substack{v\in T\\v\mid\ell}}
 \left\{b_v(e_v-c_v)
       \bigl(\delta_v f_v-w(\nu_v(\pi_v))\bigr)_+
                  +e_v b_v d_v f_v\right\}
        \ \le\ m_w(B)\ \le\ m ,
\]
where \(s_+=\max(s,0)\). Each valuation of a splitting character
in this estimate is given by Proposition~\ref{prop:cmplaces335}.
\end{theorem}
\begin{proof}
First use standard coordinates on the small side. For a tuple of
residue matrices \(\mathbf Z\), apply Theorem~\ref{thm:signedrow335}
in each variable. After multiplication by \(\prod H_v\), the coefficient
is zero if any largest Smith exponent exceeds \(e_v-c_v\).
Otherwise it equals the product of
\((\nu_v(\pi_v)q_v^{-\delta_v})^{s(Z_v)}\) times a sum of the
\(\lambda_{\mathbf W}\) with
\(W_v\supseteq\operatorname{row}(Z_v)\) at every \(v\). In this sum
each coefficient is the product of the corresponding unit phases.
Taking every \(Z_v=Y_{W_v}\) singles out any specified
\(\lambda_{\mathbf W}\) with a unit coefficient. This proves the
upper bound on the minimum in standard coordinates, and it detects
every nonzero restricted row. For the lower bound use the valuation
inequality for a sum and \(0\le s(Z_v)\le b_v(e_v-c_v)\).
Proposition~\ref{prop:cmplaces335} makes each determinant factor a
unit at primes of different residue characteristic. At \(v\mid\ell\)
its possible negative valuation is bounded in magnitude by the
first term in braces.

It remains to pass to the original small tests. In each block the
change of coordinates \(U_{B,v}\) is a Fourier matrix on \(R_{e_v}\)
in \(b_vd_v\) entries, divided by \(q_v^{e_vb_vd_v}\), while the entries
in the other columns are unchanged. Multiplying a row by \(U_{B,v}\)
can lower its minimum valuation by at most
\(e_vb_vd_v w(q_v)\). Its inverse has algebraic integer
entries, by finite character orthogonality. Thus the Fourier change
cannot raise the minimum. These matrices are invertible, which also
proves the nonvanishing assertion in the original coordinates.
The large-side Fourier changes have already been made in
\(J_{A,v}f_{W_v}\). The remaining scalars \(\sigma_v\) are exactly those
in Theorem~\ref{thm:mixedfourier334}. This gives the stated inequalities
for \(B\). At primes away from \(T\) all the \(\sigma_v\) are units.
Since the integral tests are freely generated by the residue-class
indicators, equality of the local minimum valuations there proves the
asserted equality of fractional ideals.
\end{proof}

This determines the ideal obtained when the small test ranges over
its full prescribed level. On one specified small test the value is
the contraction of the row with that particular test. Section~\ref{sec:missing}
still requires the size and divisibility of this individual signed
value and its identification with the designated Bessel pairing.

\Needspace{5\baselineskip}
\begin{corollary}[One prescribed test and its compact translates]\label{cor:compactglobal335}
On the same fields and levels, let \(f\) be the actual matched output
of a fixed small test at all marked places and restrict \(\mathcal L\)
to that image. Give the image its additive Hilbert norm in the original
charts. Conjugate the compact coordinate actions of
Theorem~\ref{thm:compactrow335} by the partial Fourier changes of variables.
For \(G=\prod_{v\in T}G_{e_v}\), let \(P\) average the product of the
congruence kernels down to \(c_v\). Then
\[
   |G|^{-1}\sum_{g\in G}|\mathcal L(g f)|^2
       \ \ge\
   \frac{\|\mathcal L P\|_*^2\,\|P f\|^2}{\prod_{v\in T}R_v}.
\]
Here \(R_v\) is the dimension in Theorem~\ref{thm:compactrow335}.
There is equality if \(e_v=c_v\) throughout \(T\), and the same
lower bound holds for the maximum of the squared values over these
translates of this one test.
\end{corollary}
\begin{proof}
The Fourier changes identify the image with the tensor product of
the frame spaces up to a scalar change of Hilbert norm at each place.
A constant change of norm cancels with its dual in the displayed
product. For the usual action on large covariant frames the common
character twist is \(\nu_v\) on determinants, which is trivial on
the congruence kernel at level \(c_v\). Thus \(P\) projects onto
the tensor product of the conductor-level subspaces. Average first
inside the cosets of its kernel group as in Theorem~\ref{thm:compactrow335}.
On the projected space the operator average for the product group
is scalar. Indeed the commuting algebra of the first factor consists
of the identity on its irreducible space tensored with arbitrary
endomorphisms of the others, by expanding in matrix units of the other
factors. Applying this to each factor leaves only scalars. Taking
the trace gives the reciprocal of \(\prod_v R_v\) as in the local
proof. This argument permits arbitrary dependence between the local
entries of both \(f\) and \(\mathcal L\). The maximum is at least the
mean. If a functional on a minimal-level factor has already been projected
by the full stabilizer and character of a specified \(W\), the same
theorem gives its value on that factor directly as \(\lambda_W f(X_W)\)
in linear frame coordinates. In general the translate \(g=1\) is one
summand of the mean. Its individual value still requires the global
comparison.
\end{proof}


\subsection{Elementary divisors of the terminal column transform}\label{sec:fouriersmith338}
The compact part \(B\) of Theorem~\ref{thm:incidence338} has a
finite matrix on each unrefined terminal class. With the compatible
twisted trivializations, its input coordinates are indexed by
nonzero residue covectors and its output coordinates by nonzero
vectors. The entry at \((v,\eta)\) is \(\psi_0(\eta(v))\), by the
reindexing identity in that theorem. The omitted zero coordinate
has zero value on a nontrivial scalar character. We compute the
elementary divisors before moving from the vector level back to
the covector level.

\begin{theorem}[Elementary divisors in each scalar character]\label{thm:fouriersmith338}
Let \(p=T\), \(d=p-1\), and write \(\zeta=\psi_0(1)\).
Work over \(R=\mathbb Z_p[\zeta]\) with uniformizer
\(\lambda=\zeta-1\), and let \(\omega\) be the Teichm\"uller
lift on \(\mathbb F_p^\times\). On input basis covectors use
\([x]\mapsto[ux]\) for the scalar \(u\), and on output functions
use \(f(y)\mapsto f(uy)\). For \(1\le j<d\), on the character
\(u\mapsto\omega(u)^j\) the elementary divisors of the unscaled
Fourier matrix are
\[
      \lambda^{|m|},\qquad 0\le m_i\le d\ (1\le i\le a),
       \qquad |m|=\sum_{i=1}^a m_i\equiv j\pmod d .
\]
Each tuple contributes one factor. In particular their number is
\(N_a=(p^a-1)/(p-1)\). With \(v_p(p)=1\), their total valuation is
\[
                    a\left(\frac{N_a-1}{2}+\frac jd\right).
\]
If the input scalar action is written with \(u^{-1}\), the
corresponding index is \(-j\bmod d\).
\end{theorem}
\begin{proof}
First include zero on both sides and identify the input with the
group ring \(R[(\mathbb F_p)^a]\) after choosing dual bases.
Let \(Z_i=[e_i]-1\). The elements \(b_m=\prod_i Z_i^{m_i}\)
with \(0\le m_i<p\) form an \(R\)-basis. Indeed in one coordinate
the binomial expansion is unit triangular in
\(1,[e_i],\ldots,[(p-1)e_i]\), and tensoring these bases covers
all coordinates. The Fourier image of \(b_m\) at \(y\) is
\(\prod_i(\zeta^{y_i}-1)^{m_i}\). After division by
\(\lambda^{|m|}\) let its value function be \(c_m\). The reduction
of \(c_m\) is \(\prod_i y_i^{m_i}\), because
\(((1+\lambda)^t-1)/\lambda\equiv t\pmod\lambda\) for
\(0\le t<p\). These monomials form a basis of the functions on
\(\mathbb F_p^a\). In one coordinate a polynomial of degree
less than \(p\) vanishing at all \(p\) arguments is zero, and the
multivariable evaluation matrix is a tensor product. Hence the
\(c_m\) form an \(R\)-basis, since their determinant is a unit.

The order \(d\) of the scalar group is a unit in \(R\). Its
character projectors \(e_j=d^{-1}\sum_u\omega(u)^{-j}[u]\)
therefore act integrally, with \([u]\) denoting the scalar action.
They commute with Fourier transformation under the two specified
actions. The reduction of \(c_m\) already has eigencharacter
\(u^{|m|}\). On the input side the group ring modulo \(\lambda\)
is \(\mathbb F_p[Z_1,\ldots,Z_a]/(Z_1^p,\ldots,Z_a^p)\).
The action of \(u\) takes \(Z_i\) to \((1+Z_i)^u-1\), whose
linear term is \(uZ_i\). Therefore it sends \(b_m\) to
\(u^{|m|}b_m\) plus terms of strictly higher total degree
modulo \(\lambda\). With \(j\equiv |m|\), the union of the
\(e_j b_m\) is a basis by unit triangularity after reduction.
The corresponding \(e_j c_m\) form a basis by the monomial
calculation. Separating the bases by their exact characters now
gives the displayed diagonal factors. For \(j\ne0\) every vector
in the character space has zero coefficient at the fixed origin,
and the same holds on the output. Deleting zero has no effect.

For a residue \(l\bmod d\), the number of tuples of length \(b\)
with total degree \(l\) is \(N_b+\mathbf1_{l=0}\), with \(N_0=0\).
Indeed one coordinate in \(0,\ldots,d\) represents every residue
once and zero once more. Induction gives the formula since
\(N_{b+1}=pN_b+1\). By symmetry the sum of the degrees in the
\(j\)-component equals
\[
   a\sum_{t=0}^d t\big(N_{a-1}+\mathbf1_{t\equiv j\ (d)}\big)
                   =a\big(N_{a-1}dp/2+j\big).
\]
Finally \(p\) is a unit times \(\lambda^d\), by
\(p=\prod_{u=1}^{p-1}(1-\zeta^u)\). Division by \(d\) gives the stated
valuation. Extending the coefficient field leaves these normalized
valuations unchanged.
\end{proof}

For the \(f\) places on the interval the compact transform is a
tensor product on each class of Corollary~\ref{cor:rank111}.
The stabilizers there are trivial, so the finite matrix acts on each
coefficient coordinate. For fixed nontrivial character indices
\(j_v\), its elementary-divisor valuations are
\(\sum_v |m_v|/(T-1)\) with each \(|m_v|\equiv j_v\bmod(T-1)\).
Its mean valuation per dimension differs from \(af/2\) by
\(a\sum_v(j_v/(T-1)-1/2)/N_a\). Away from \(T\) finite Fourier
inversion on the separately zero-sum lattice has a unit denominator
even at primes dividing \(T-1\), by the integral sums in
Theorem~\ref{thm:incidence338}. At those primes this last assertion
uses the full subgroup with zero fiber sums and requires no division
into character eigenspaces.

\begin{proposition}[The change of coefficient lattice for one adjacent class]\label{prop:weightstep338}
At a split coefficient place use the integral flag-minor realization
of weight \(\mu_1\ge\cdots\ge\mu_a\), with the determinant power
included. For a full column lattice \(M\) write \(W(M)\) for the
resulting coefficient lattice. For a sublattice \(M'\subset M\)
of residue index \(p\) one has
\[
          \pi^{\mu_1}W(M)\ \subseteq\ W(M')                         \ \subseteq\ \pi^{\mu_a}W(M).
\]
In a basis adapted to \(M'\), its relative factors are
\(\pi^t\), where \(t\) is the last-coordinate weight with its
multiplicity in the representation. For a tensor product of split
coefficient factors, the exponents add over the affected factors.
\end{proposition}
\begin{proof}
Choose a basis of \(M\) in which only the last basis vector must
be multiplied by \(\pi\) to obtain \(M'\). Factor the representation
as the indicated power \(\det^{\mu_a}\) times the polynomial weight
with entries \(\mu_i-\mu_a\). The flag-minor generators for that
weight use one minor per column of the Young diagram, in total
\(\mu_1-\mu_a\) minors. The last basis vector occurs at most once
in each minor. Thus each generator has last-coordinate degree
\(t\) between \(\mu_a\) and \(\mu_1\), including the determinant
power. The lattice is the direct sum of the spans of generators of
each weight, since distinct weight spaces are independent over the
fraction field. These spans are free over the local valuation ring.
The basis change multiplies every vector in the degree-\(t\) summand
by \(\pi^t\). This gives both inclusions, also for negative weights.
Tensoring the weight-space bases proves the last assertion.
\end{proof}

For the column tables formed with these coefficient lattices the
factor \(\rho(h^{-1})\) in \(D\) evaluates the input at the adjacent
index-\(p\) sublattice of the output lattice. In local frames this
is the inclusion just computed. The determinant twist of
Section~\ref{sec:graphtransport337} takes unit values and sends the
uniformizer to one. The compact sums in \(B,C\) preserve the
coefficient lattices within each unrefined class. Thus, if \(S^\circ\)
denotes the separately zero-sum table at this level, a single split
coefficient factor satisfies
\[
             \pi^{a+\mu_1} S^\circ\ \subseteq\ D S^\circ
                                      \ \subseteq\ \pi^{\mu_a}S^\circ .
\]
For the left inclusion use \(p^aD^{-1}=C\rho(h)\) and the reverse
change along the same inclusion of lattices. For several affected
factors replace \(\mu_1,\mu_a\) by their respective sums. For a
dual testing lattice use its dual highest weight. These inclusions
retain the algebraic weight cost and the Fourier cost separately.


\subsection{The unmarked translation in an integral period}\label{sec:splittranslation343}
Fix one of the fields and the signed rows and columns of
Theorem~\ref{thm:rationalperiod342}. At unmarked finite places use
the prescribed compact levels and theta tests. Outside a finite set
these levels are hyperspecial and the tests are unit lattice
indicators, as in Lemma~\ref{lem:finite167}. The rational flag has an
integral basis compatible with these lattices outside finitely many
places. Indeed the transition matrices between rational bases and
their inverses are integral outside the primes in their denominators,
and adelic lattices agree with a fixed rational lattice almost
everywhere. Include all exceptional unmarked places for the data
under consideration in a finite set. At each underlying rational
prime, membership in the open stabilizers of the smooth tests and
vectors follows if all $c_j$ are one to sufficiently high congruence
precision. This follows by continuity of the map from the flag torus
to the local groups. The same applies to any fixed finite list of
data by intersecting the open stabilizers. Take $M$ divisible by
these prime powers and by $T^b$ in Theorem~\ref{thm:splitflag343}.

Let $\mathcal U$ be the set of unmarked places of $F$ at which any
of the flag multipliers has nonzero valuation. Each is now good
and lies over a rational prime splitting completely in $E$.
At the chosen half $\mathfrak p$ above $v\in\mathcal U$, the
initial list $\lambda_v$ of valuations on the $r$ row coordinates
consists of the entries
\[
      v_{\mathfrak p}(c_j),\quad -v_{\bar{\mathfrak p}}(c_j)
                          \quad (1\le j\le s)
\]
and zeros on the remaining coordinates. Set $B_v=\sum_t|\lambda_{v,t}|$
and let $q_v$ be its residue cardinality. Complete splitting gives
\[
                   \sum_{v\in\mathcal U}B_v\log q_v
                                  \le \sum_{j=1}^s h_f(c_j).
\]
At an unmarked place outside $\mathcal U$ the translation fixes the
chosen data, either by the congruence at an exceptional place or by
integrality in the hyperspecial lattice.

\begin{proposition}[The unit lattice tests at all compensating places]\leavevmode\label{prop:unmarkedscore343}
On a growing string of $m$ theta edges let $r_k$ be the rank at
vertex $k$, with $r_0=r$. At every $v\in\mathcal U$ extend the
diagonal by the identity in the extra coordinates of each vertex.
Write $\lambda_{k,v}$ for $\lambda_v$ extended by these zeros.
Let $\Phi^0$ be the tensor of the unmarked unit matrix lattice
indicators at the places in $\mathcal U$, on all the edges.
The normalized absolute inner product $C$ of this tensor with its
simultaneous Weil translate is positive and satisfies
\[
 \begin{aligned}
 -\log C
 &=\tfrac12\sum_{v\in\mathcal U}\sum_{k=0}^{m-1}
       \sum_{\alpha=1}^{r_{k+1}}\sum_{\beta=1}^{r_k}
       |\lambda_{k+1,v,\alpha}-\lambda_{k,v,\beta}|\log q_v\\
 &\le \tfrac12\sum_{k=0}^{m-1}(r_k+r_{k+1})
                              \sum_{j=1}^s h_f(c_j).
 \end{aligned}
\]
The choice of $\mathcal U$ and the formula use the same rational
translation as the Bessel identity of Theorem~\ref{thm:rationalperiod342}.
\end{proposition}
\begin{proof}
In split coordinates the two supports on edge $k$ are the lattice
$W=M_{r_{k+1},r_k}(O_v)$ and $W'=d_{k+1,v}W d_{k,v}^{-1}$.
With Haar measure giving $W$ volume one, the normalized absolute
inner product of their indicators equals
$([W:W\cap W'][W':W\cap W'])^{-1/2}$. The scalar in the Weil
action cancels on dividing by the product of the two norms. In
matrix entry $(\alpha,\beta)$ the exponent of the second lattice
relative to the first is
$\lambda_{k+1,v,\alpha}-\lambda_{k,v,\beta}$. The product of the
two indices on this entry is $q_v$ to the absolute value of this
exponent. Multiplying over entries, edges and places gives the identity.
For each edge the triangle inequality bounds the sum of absolute
differences by $(r_k+r_{k+1})B_v$. The preceding height estimate
gives the asserted upper bound for $-\log C$.
\end{proof}

\begin{theorem}[A common denominator for returning the unmarked row to its level]\leavevmode\label{thm:unmarkedtrace343}
At $v\in\mathcal U$ write $K_v=\mathrm{GL}_r(O_v)$ in the chosen
split frame and $t_v=d_{*,v}$. Put $K=\prod_{v\in\mathcal U} K_v$
and $t=\prod_{v\in\mathcal U}t_v$, and normalize the average $P_K$
on the automorphic row space by Haar probability. There is a positive
integer $I$ and a fixed list of $I$ finite translates such that for
every $K$-fixed row $A$ one has
\[
   I\,P_K R(t) A=\sum_{\nu=1}^{I} R(\gamma_\nu t) A,
          \qquad \log I\le r^2\sum_{j=1}^s h_f(c_j).
\]
The representatives $\gamma_\nu\in K$ and $I$ are independent of
the row and column. The formula remains an exact identity on applying
the global Bessel functional, the fixed signed projections at the
marked places, or a linear period map. In particular, if the resulting
values of the summands belong to a fractional ideal $\mathfrak a$
in a common coefficient field, the projected value belongs to
$I^{-1}\mathfrak a$. Both displayed logarithmic bounds in this section
are polynomial in $x$ when $\log M$ and $\log(1/\delta)$ are
polynomially bounded on the common intervals.
\end{theorem}
\begin{proof}
Set $J_v=K_v\cap t_vK_vt_v^{-1}$ and
$D_v=\max_t\lambda_{v,t}-\min_t\lambda_{v,t}$. If $D_v>0$ then
$1+\pi_v^{D_v}M_r(O_v)\subset J_v$. Indeed conjugation by $t_v^{-1}$
and $t_v$ scales the entries of the difference from the identity by
factors of valuation at least $-D_v$. The resulting matrix is
integral and still has unit determinant. Reduction modulo
$\pi_v^{D_v}$ now gives
$[K_v:J_v]\le q_v^{r^2D_v}$. If $D_v=0$, $t_v$ is a scalar times
a unit diagonal, so $J_v=K_v$. In all cases $D_v\le B_v$.

Put $J=\prod_v J_v$, choose representatives of the cosets
$\gamma_\nu J$ in $K$ and let $I=[K:J]$. The function
$k\mapsto R(kt)A$ is constant on each such coset since
$t^{-1}Jt\subset K$. Every coset has measure $1/I$, proving the
identity. Multiplying the index bounds proves the assertion on
$\log I$. Linearity gives the assertion for functionals and fractional
ideals. For the polynomial bounds use Theorem~\ref{thm:splitflag343}
and the polynomial bounds for $m,r_k,s$ on the indicated fields.
\end{proof}

This computation determines the unmarked lattice overlap and the
cost of the additional level average $P_K$ on the actual row.
Passing from the period of $R(t)A$ in
Theorem~\ref{thm:rationalperiod342} to the prescribed integral
theta period still requires identifying the endpoint functional.
In particular the Bessel functional need not be invariant under $K$
on the span of these translates, and the signed automorphic values
in the finite sum may cancel.

\begin{corollary}[Avoidance through the cohomological prime range]\leavevmode\label{cor:avoidrange343}
Take the same fields, points $x\in[V,V^{3/2}]$ and cutoff
$Y=x^{-3}\exp(x^\varepsilon/4)$ as in
Corollary~\ref{cor:allsample341}. Fix $\delta=\exp(-x^B)$ for
any fixed $B>0$, and take $b=1$. There is a matching rational
diagonal all of whose unmarked nonzero valuations lie over completely
split rational primes greater than $Y$. Its coordinates satisfy
\[
              \sum_{j=1}^s h_f(c_j)
                     \le \exp\bigl((1/4+o(1))x^\varepsilon\bigr).
\]
For rows fixed by the full stabilizer of the prescribed signed
coordinate lattice at unmarked places, the integer returning the
row to those levels satisfies the same bound on $\log I$.
For the two estimates on the automorphic data above, let $M_0$ be
an integer imposing the additional unmarked stabilizer congruences.
If $\log M_0\le\exp((1/4+o(1))x^\varepsilon)$ uniformly on the
intervals, the same choice can impose those levels and gives
\[
           -\log C+\log I=o\bigl(\exp(x^\varepsilon/2)\bigr).
\]
\end{corollary}
\begin{proof}
In Theorem~\ref{thm:splitflag343} take $M$ to be $T$ times the product
of all rational primes up to $Y$ except $T,g_1,g_2$. This sets all
unmarked valuations through $Y$ equal to zero. The added primes
increase $\log M$ by at most $Y\log(2Y)$, by counting the factors
and their individual sizes. All the remaining terms in $L$ are
polynomially bounded, as are the number of coordinates and the
prefactor in the height bound. Since $\log x=o(x^\varepsilon)$,
any fixed polynomial factor times $Y\log(2Y)$ is bounded by
$\exp((1/4+o(1))x^\varepsilon)$. This proves the coordinate assertion.

For the full row lattice level use the rational frame of
Section~\ref{sec:char94}. Its two vectors on each signed plane are
$e_++e_-$ and $(e_+-e_-)/2$. They are an integral basis at odd
primes. When both diagonal entries $c,\bar c^{-1}$ are units, the
two operators on the plane given by the diagonal and its inverse
therefore preserve the lattice. At dyadic primes we have
$c,\bar c^{-1}\equiv1\bmod 2$ by the choice of $M$. Their sum and
difference divided by two are integral, as are the corresponding
entries for the inverse. These are precisely the entries in the
original basis $e_+,e_-$, by the calculation of
Theorem~\ref{thm:neighbour93}. The other row coordinates are fixed.
Thus all unmarked places with zero valuations already fix a row of
this level. The remaining places are odd and split completely,
where the index proof in Theorem~\ref{thm:unmarkedtrace343} applies
to this lattice with the displayed exponents. Its bound for $\log I$
has just the additional polynomial factor $r^2$.

For the data with additional stabilizers take the least common
multiple with $M_0$, which is coprime to $g_1g_2$. Under the stated
hypothesis the logarithm of this modulus has the same bound.
Apply Proposition~\ref{prop:unmarkedscore343} and
Theorem~\ref{thm:unmarkedtrace343}. Their polynomial rank factors
are absorbed in the $o(1)$ in the exponent. Since $1/4<1/2$, this
gives the last display. The integer $I$ may have prime factors at
most $Y$ even when each local translation in $\mathcal U$ is above
$Y$. Their total weighted valuation is bounded by $\log I$ by
factorization of this single integer.
\end{proof}

\section{Arithmetic theta columns on marked classes}\label{nav:chapter-marked-theta-columns}
\subsection{The coefficient formula on the prescribed terminal classes}\label{sec:terminalformula347}

Fix one of the fields \(E_x\), a translate, a selected word and its
terminal prime \(T=\pi\bar\pi\) as in Theorem~\ref{thm:common111}.
Choose a coefficient field containing the images of the CM
embeddings \(\tau\). At \(\tau\) use the actual large and small
terminal highest weights \(\lambda_\tau,\mu_\tau\). Write
\(l_\tau,m_\tau,b_{\tau i},c_\tau,\kappa_\tau\) for the
partitions, last-row exponents and constants in
Section~\ref{sec:lowerlattice347}. Put
\[
 \delta_\tau=\sum_i\lambda_{\tau i}-\sum_i\mu_{\tau i},\qquad
 Z=\prod_\tau \tau(\bar\pi)^{-\delta_\tau},\qquad
 U=\prod_\tau\frac{c_\tau\kappa_\tau}{H_{m_\tau}}.
\]
Use the finite class pairing with coefficient \(1/|\Gamma_j|\)
on each decorated terminal class \(j\). In adelic Haar measure
this is the pairing divided by the volume of its right level, as in
Theorem~\ref{thm:classtrace326}. Let \(P_j\) be the resulting
large polynomial obtained by evaluating the entire signed column at
the lift of the chosen small class representative in the fixed
polynomial frame. Let \(F_j\) be the polynomial of the testing
form at the original small representative, with the dual action in
the bilinear pairing. This convention includes the transport from
the large class representative and all determinant characters in
\(P_j\). It fixes the same representative for both evaluations.
The last hyperplane is the first \(n\) coordinates in rank
\(N=n+1\). Its normal coordinate uses \(\bar\pi\) in the
algebraic restriction. In fact an element acting as \(z\) on
the normal coordinate acts on the \(\mu_\tau\)-summand by
\(z^{\delta_\tau}\), by the two central weights. Pullback by
its inverse multiplies by \(z^{-\delta_\tau}\). The construction
takes \(z=\tau(\bar\pi)\) at that embedding. The representative,
the map in this frame and the class pairing are specified in
\srcloc{Formalism/LocalToGlobal/Columns/FiniteSumNonzeroLargeDepth}{197}
and lines 203--218 there.

\begin{theorem}[The signed period as a sum of specified coefficients]\leavevmode
\label{thm:terminalformula347}
For a joint monomial index \(\alpha=(\alpha_\tau)_\tau\) on the
smaller matrices let \(\widehat\alpha\) extend it by zero to
the last row and column at all embeddings. Let \(B\) have entries
\(b_{\tau i}\) at \((\tau,N,i)\) and zero elsewhere. For the
prescribed class-sum normalization the exact scalar is
\[
 {\cal P}(P,F)=ZU\sum_j\frac{1}{|\Gamma_j|}
          \sum_\alpha \alpha!\,(P_j)_{\widehat\alpha+B}(F_j)_\alpha ,
 \qquad \alpha!=\prod_{\tau,i,i'}\alpha_{\tau ii'}! .
\]
The formula holds at every complex embedding and over the coefficient
field whenever the entries are defined there. At every finite place
over \(p>\max_\tau(\lambda_{\tau1}-\lambda_{\tau N}+n)\), the
factor \(U\) is a unit and all nonzero monomial weights \(\alpha!\)
occurring in this pairing are units. For the word parameters
of \eqref{eq:letters111} this covers \(g_1,g_2,T\) simultaneously
at every sufficiently distant \(x\) on the intervals of
Theorem~\ref{thm:criticalmarks339}.
\end{theorem}
\begin{proof}
At a single representative, multiply the large polynomial by the
last restriction operator in Theorem~\ref{thm:lowerlattice347},
including the separate normal-coordinate factor. In its pairing
with \(F_j\), the orthogonal projection can be moved to the second
slot since it is self-adjoint. It fixes \(F_j\), which belongs to
the smaller flag space. The normalized bilinear Fischer form now
contributes \(\alpha!/H_m\) for the matching smaller monomial.
Coefficient extraction reads exactly the larger monomial indicated
in the statement. Tensoring gives the inner sum. Integration of a
class function counts each representative with weight
\(1/|\Gamma_j|\). Linearity keeps all the signs, class transports
and earlier theta operations inside the same \(P_j\). Every
operation is defined by rational coefficients and the displayed
field elements, which proves the embedding assertion.

Lemma~\ref{lem:highestline347} gives the assertion for \(c_\tau\)
and \(\kappa_\tau\). The hook bound in its proof applies to
\(H_{m_\tau}\). Each exponent of a smaller monomial is at most
the degree of its column, hence at most \(\mu_{\tau1}-\mu_{\tau n}\).
This proves the factorial assertion. For the actual weights, all
the ranks and the spacings \(\Delta,H_0\) are polynomial in \(x\).
The large parameters consist of the initial \(u_i\) with their
rank shifts and the repeated \(-e_j\) with rank shifts. In
\eqref{eq:letters111} and its following formulas all \(u_i,e_j\)
are \(O(xb+x^C)\) for a fixed \(C\), with \(b=\exp(x^\eps)\).
Thus every cutoff in the theorem is \(O(xb+x^{C'})\).
The chosen \(g_i\) exceed \((kb)^8\), and \(\log T\gg m x^\eps\)
with \(m\to\infty\), by Theorems~\ref{thm:levels88}
and~\ref{thm:markedstrict337}. These inequalities hold throughout
the same intervals. They also put all three primes above
\(1+n[E_x:\mathbb Q]\). A class stabilizer acts faithfully on the
underlying integral lattice of that rank. If its order were divisible
by \(p\), an element of order \(p\) would have a primitive
\(p\)-th root as an eigenvalue over \(\mathbb Q\), of degree
\(p-1\). This proves also that the displayed mass denominators
are units at these three primes.
\end{proof}

\begin{corollary}[The normal coordinate on a common terminal prime]\leavevmode
\label{cor:normaldegree347}
Write \(d=t_{\rm r}+1\), and use the \(u_i,v_i,b_j,e_j\) at one
embedding in the lists of \eqref{eq:letters111}. Then
\[
 \delta_\tau=\sum_{i=1}^k u_i-\sum_{i=1}^{k-1}v_i
                       +\sum_{j=1}^m (b_j-e_j)
             =\sum_i u_i-\sum_i v_i-mH_0.
\]
For a fixed translate this integer at each embedding, and the element
\(Z\) itself, are independent of the word in
Theorem~\ref{thm:common111}, with the same generator \(\pi\).
Normalize \(v_w(T)=1\) for \(w\mid T\). Then
\[
 v_w(Z)=-\!\sum_{\tau:\,w\mid\tau(\bar\pi)}\delta_\tau ,
 \qquad |\sigma Z|=T^{-\sum_\tau\delta_\tau/2}
 \quad(\sigma:J\hookrightarrow\mathbb C).
\]
At all finite places outside \(T\) the valuation of \(Z\) is zero.
\end{corollary}
\begin{proof}
Write \(s_0=md/2\) and \(\rho_j(i)=(j+1)/2-i\), with one-based
indices. In the large list the initial entries are
\(u_i-\rho_k(i)-s_0\). The following \(m\) blocks each have
length \(d\), with constant value
\[
       L_j=-e_j+k/2+(2j-1)d/2-s_0.
\]
The small list begins with \(v_i-\rho_{k-1}(i)-s_0\) for
\(i<k\). Its \(j\)-th block consists of one entry
\(A_j=-b_j+k/2+(j-1)d-s_0\) followed by \(d-1\) entries
\(L_j\). The shifted highest weights are obtained by subtracting the
half-sum of positive roots from the row parameters. Within the
string following \(-b_j\), each parameter and its root shift
decrease by one, giving the repeated value \(L_j\).
The first entries interlace by \(u_i>v_i>u_{i+1}\) and
half-integrality. At the first junction use
\(u_k+b_1\ge1/2\), and at later junctions use
\(b_j-e_{j-1}\ge(d+2)/2\).
Also \(A_j-L_j=e_j-b_j-d/2\ge0\). All remaining comparisons
are equalities. In particular both final entries equal \(L_m\),
as required in Section~\ref{sec:lowerlattice347}.

The root shifts sum to zero in each initial list. Subtracting
their sums and those of the blocks gives
\(\sum u_i-\sum v_i-s_0+\sum_j(L_j-A_j)\), which is
the formula since \(L_j-A_j=b_j-e_j+d/2\).
The only word-dependent contribution is
\(-\Delta\sum_{i=2}^{k-1}c_i\). Its sum is fixed in the selected
family. Finally \(\pi\) belongs to the quadratic field used for
the split generator, \(\pi\bar\pi=T\), and their ideals are
coprime. At \(w\mid T\) each conjugate of \(\bar\pi\) consequently
has valuation zero or one. At every archimedean embedding its
absolute value is \(\sqrt T\). Taking products proves the claims.
\end{proof}

The formula places the known normalization and the cohomological
vanishing at the marked primes on the same fields. The coefficients
of the earlier signed theta transforms enter as the values
\((P_j)_{\widehat\alpha+B}\). Their cancellation against the tests
and the comparison with the incoming Bessel row still enter the
estimate in Section~\ref{nav:I-28}.


\subsection{The change of determinant on the image of the arithmetic period}
\label{sec:imagedet355}

Take the period fields, common intervals and coded weights of
Theorems~\ref{thm:region198}, \ref{thm:common111}
and~\ref{thm:prime210}. For the moment fix one terminal prime
\(T\). Take the direct sum over any selected words with this
prime and their ordinary initial classes. Denote the sum of the
actual initial lattices by \(S\), the corresponding sum of the
unrefined vector lattices in Section~\ref{sec:fourierslopes354}
by \(V\), and their common rank by \(d>0\).
Write \(\mathcal F\) for the direct sum of the unnormalized
transforms. Work over a common coefficient field \(J\).

Let \(B\) be any \(J\)-linear map on \(V\otimes J\) of rank
\(j>0\). In its image compare the two full lattices
\[
               M_0=B(V),\qquad M_1=B(\mathcal F S).
\]
In particular \(B\) can be the signed arithmetic period on the
vector coordinates. Its blocks are the centered maps of
Theorem~\ref{thm:rankone354} divided by \(Q=T^{fk}\), with the
specified terminal lowering and projections. In that case
\(M_1\) is the image of the actual ordinary input under the
signed column. The reference image \(M_0\) uses exactly the
same \(B\), with the unrefined integral vector input. These
lattices measure the input contribution to the period
determinant even if \(B\) has a large kernel.

At a prime over \(T\) use \(v(T)=1\). Write \(\ell_v\) for the
valuation of a local generator of \(\det M_1\) in a local
basis of \(\det M_0\). Set \(a=\lceil xb\rceil\),
\(b_*=\lceil xb/\exp(x^{\eps/3})\rceil\). Choose a positive
polynomial bound \(C(x)\) for \(\lambda_{\tau,i}+k-i\) for
all \(i\ge2\), and also for \(i=1\) away from the high places.
Such a bound uniform in the words follows directly from
\eqref{eq:letters111}. At a high place the first entry before
adding \(k-1\) is \(a\).

\begin{theorem}[Determinant bound for the actual projected image]\leavevmode
\label{thm:imagedet355}
At every prime over this \(T\),
\[
 -f b_* \ \le\ \frac{\ell_v}{j}\
 \le\ fk+f C(x)
        +\frac{2(a+k)}{k-1}\bigl(f\log 2+\log(d/j)\bigr).
\]
The determinant valuation is zero away from \(T\). Uniformly in
every such map, on the same sufficiently distant intervals,
\[
 \begin{aligned}
   \frac{|\ell_v|}{d}\log T
       &\ll_{\psi,\eps} b\,x^{-1-17\eps/16},\\
   \frac{|\ell_v|}{j}\log T
       &\ll_{\psi,\eps} b\,x^{-1-17\eps/16}
              &&\text{if }\log(d/j)\le x^{2+8\eps}.
 \end{aligned}
\]
In the quotient Hilbert norms induced from the ordinary class
norm on \(S\) and the unrefined class norm on \(V\), the norm
of every image vector for \(B\mathcal F\) is \(Q^{-1/2}\)
times its norm for \(B\).
\end{theorem}
\begin{proof}
Use a uniform random index in the list of identity exponents
\(\alpha_i\) in Theorem~\ref{thm:fourierslopes354}. Within each
word and class it has the distribution of a sum of independent
single-coordinate weights at the affected embeddings. Indeed,
the weight lists are tensored with their multiplicities. The
other embeddings and the orbit multiplicities only repeat
the distribution. At an affected embedding every exponent is
at least \(\lambda_{\tau,k}\ge-b_*\), by the coordinate bounds
of Lemma~\ref{lem:openflags353}. There are at most \(f\)
affected embeddings.

Put \(t=(k-1)/(2(a+k))\). The initial lists satisfy the
hypotheses of Lemma~\ref{lem:oneweight355} with \(A=a+k\)
for large \(x\), since \(a\) dominates every fixed polynomial.
Independence gives
\[
         \frac{1}{d}\sum_i e^{t\alpha_i}
                   \le e^{tf C(x)}2^f.
\]
Each summand space has this bound and the same lower endpoint
\(-f b_*\). Taking the mixture with their dimension weights
proves both assertions for the direct sum. For its Fourier
exponents \(\beta_i\), the same exponential moment gains at
most \(e^{tfk}\), by Theorem~\ref{thm:fourierslopes354} on
each summand. Its least exponent is still at least \(-fb_*\).

Use Smith bases for \(\mathcal F S\) and \(V\). In these bases
the former is obtained from the latter by multiplying coordinate
\(i\) by a scalar of valuation \(\beta_i\). In a local basis
of \(M_0\) the matrix of \(B\) on \(V\) is integral and has a
unit \(j\)-minor, since its columns generate \(M_0\). For a
set \(I\) of \(j\) columns, the valuation of the corresponding
minor on \(\mathcal F S\) is its old valuation plus
\(\sum_{i\in I}\beta_i\), with valuation \(+\infty\) for zero.
The minimum of these values is \(\ell_v\). It is at least the
sum of the smallest \(j\) exponents and at most the sum of the
largest \(j\), by using a unit minor for the latter bound.
The average of \(e^{t\beta_i}\) on the largest \(j\) terms
is at most \(d/j\) times its full average. Jensen's inequality
for the exponential now gives the stated upper bound for
\(\ell_v/j\). The lower endpoint gives the opposite bound.
Away from \(T\) the two source lattices agree after applying
the transform, so their images agree.

By Theorems~\ref{thm:rangenew50}, \ref{thm:common111}
and~\ref{thm:prime210}, the ranks and terminal scale satisfy
\[
 f\ll_\psi x^{2+8\eps},\quad k\asymp x^{9/2+10\eps},\quad
 m\asymp x^{1/2-\eps/16},\quad \log T\ll m x^\eps .
\]
In particular \(xf\log T/k\ll_{\psi,\eps}x^{-1-17\eps/16}\).
Every fixed polynomial in \(x\) and the quantity \(f b_*\log T\)
are \(o(b/x^A)\) for any fixed \(A>0\). If
\(\log(d/j)\le x^{2+8\eps}\), substitution gives the estimate
per image dimension. For all \(j\) multiply the first display
of the theorem by \(j/d\) and use
\((j/d)\log(d/j)\le1/e\). This gives the estimate per source
dimension. Finally \(\mathcal F/\sqrt Q\) is unitary by
Theorem~\ref{thm:vectorchain353}. Taking the infimum of the
norms of lifts of the same image vector under the two maps
gives the asserted quotient norm.
\end{proof}

\begin{corollary}[A joint period map with several terminal primes]\leavevmode
\label{cor:jointimage355}
At a fixed \(x\) take any nonempty selection of the words and
translates on this interval, allowing their selected terminal
primes to differ. Suppose their signed arithmetic columns take
values in one coefficient space, with their specified tests
and normalizations. On the direct sum of their initial spaces
use the joint map, adding their outputs. Let \(d\) be the
total source dimension and \(j\) its image dimension. Define
\(M_0,M_1\) as above using each summand's own transform.
Normalize \(v(p)=1\) at a prime of \(J\) above \(p\). Then
\[
  \frac1d\sum_{p}\sum_{v\mid p}
     \frac{[J_v:\mathbb Q_p]}{[J:\mathbb Q]}\,
               |\ell_v|\log p
          \ll_{\psi,\eps} b\,x^{-1-17\eps/16}.
\]
In the quotient metrics the absolute change of the logarithm
of the determinant norm at any embedding, divided by \(d\),
is at most \(\tfrac{j}{2d}fk\max_T\log T\) when the ranks
\(f,k\) are shared. With varying ranks use
\(\tfrac{j}{2d}\max_T (f_T k_T\log T)\), taking the maximum
over summands with their actual ranks. These bounds are uniform
in the number of summands. They hold also in rank zero, using
the unit determinant in degree zero.
\end{corollary}
\begin{proof}
At a given rational terminal prime \(T\), let \(d_T\) be the
sum of the dimensions of all summands with that prime. Other
summands contribute only zero elementary-divisor exponents,
by Theorem~\ref{thm:fourierslopes354}. By the minor argument
above, \(|\ell_v|\) is bounded by the sum of the positive
exponents plus the absolute sum of the negative ones in the
active summands. In each active block of dimension \(d'\) the
absolute sum of the negative ones is at most \(d' f b_*\).
The sum of the positive ones is the full signed sum plus that
absolute negative sum. The preceding theorem in full rank,
applied to the identity map on the unrefined vector lattice
of this block, bounds the absolute full signed sum times
\(\log T\) by \(O_{\psi,\eps}(d'b x^{-1-17\eps/16})\).
The negative contribution satisfies the same bound. Summing
the blocks gives this estimate with \(d_T\) in place of \(d'\).
The local degree weights above one \(T\) sum to one. Finally
\(\sum_T d_T=d\); all other finite primes contribute zero.

At infinity the direct sum of transforms is a unitary change
of coordinates followed by positive scalars \(\sqrt{Q}\) on
the individual summands. The norm of a preimage in the new
coordinates is between the reciprocal of the largest of these
scalars and the reciprocal of the smallest, times the old
norm. Taking infima over all preimages gives the same bounds
on the image. On a \(j\)-dimensional determinant line their
logarithms multiply by \(j\), proving the stated estimate.
\end{proof}

If a further finite-index sublattice \(S'\subseteq S\) is
imposed before applying the signed column, the quotient
\(B(\mathcal F S)/B(\mathcal F S')\) is an image of \(S/S'\).
Thus its extra local length is bounded by that of the imposed
restriction. The reference image \(B(V)\) in this comparison
still depends on the signed theta coefficients. Estimating
its determinant in the prescribed Bessel test lattice and
bounding the residual orders in Proposition~\ref{prop:res23}
are the unresolved steps in Section~\ref{sec:missing}.


\subsection{Forgetting auxiliary frames in the terminal period}
\label{sec:rowfibers380}

Fix a field \(E=E_x\), a word and a translate from the interval
in Section~\ref{sec:missing}. Put \(n=r+mt\) for the small terminal
rank and \(N=n+1\). The primes above the two auxiliary rational
primes at which frames are retained form a set \(S_g\). The row
can have additional frames at a set \(Q\) disjoint from \(S_g\).
For a residue place \(w\) write \(k_w=\mathcal O_E/w\).
At the terminal prime \(T=\pi\bar\pi\), let \(S_T\) consist
of the primes dividing \((\pi)\). The prime \(T\) splits completely
in \(E\), and \((\pi)\) and \((\bar\pi)\) are relatively prime.

A small object consists of a positive definite unimodular Hermitian
lattice \(M\) in the specified space and its residue frames.
Each frame is an isomorphism \(M/wM\simeq k_w^n\).
The large object used for evaluation has lattice
\(M^+=M\perp\mathcal O_E e\), with \(e\) its last unit vector.
Its covector at \(P\in S_T\) and frame at \(w\in S_g\) are
\[
 f_P(v+a e)=\bar\pi a\pmod P,\qquad
 \beta_w^+(v+a e)=(\beta_w(v),\pi^{-1}a\bmod w).
\]
Here \(\pi\) is a unit at every \(w\in S_g\). No frame at \(Q\)
occurs in either formula. These are the decorations in
\srcloc{Formalism/LocalToGlobal/Columns/StdLatUnimodular}{168}.
Their use in the finite period, with a specified representative of
each row class, is described in
\srcloc{Formalism/LocalToGlobal/Columns/FiniteSumNonzeroLargeDepth}{197}.

Let \(C_w\subseteq\mathrm{GL}_N(k_w)\) be the column level, with
its character. For each retained frame define
\[
 J_w=\{h\in\mathrm{GL}_n(k_w):\operatorname{diag}(h,1)\in C_w\}.
\]
Give \(J_w\) the restriction of the column character along this
block embedding. Denote the class set of small lattices with
these levels and only the retained frames by \(\mathcal C\).
Morphisms preserve the lattice and the stated orbits of frames.
On the large side they also preserve every specified covector.
Write \(i(c)\) for the class obtained by adjoining \(e\).

\begin{proposition}[Recovery of the small lattice with its induced level]
\leavevmode\label{prop:liftclasses380}
For arbitrary groups \(C_w\), extension by the identity on \(e\)
is a bijection on the sets of morphisms between any two lifted
objects and their small objects at the induced levels.
Consequently \(i\) is injective and it preserves stabilizer orders.
\end{proposition}
\begin{proof}
Write the image of \(e\) under a large isometry as \(v+a e\).
Lattice preservation gives \(a\in\mathcal O_E\).
Evaluating each covector at \(e\) gives \(a\equiv1\pmod P\)
for \(P\in S_T\), since \(\bar\pi\) is a unit there. At every
embedding \(\tau\), definiteness and orthogonality give
\(|\tau a|\le1\). If \(a-1\ne0\), its absolute norm is at
most \(2^{[E:\mathbb Q]}\) and is divisible by
\(\prod_{P\in S_T}|k_P|=T^{[E:\mathbb Q]/2}\).
The latter is larger because \(T>4\). Thus \(a=1\) and
\(v=0\) by the length identity. The complement of \(e\) is
preserved and the isometry is \(\gamma_0\oplus1\).
The frame relation on \(M/wM\) and on \(e\) forces its entire
residue change of frame to be \(\operatorname{diag}(h_w,1)\).
Indeed the images of \(M/wM\) fill the first \(n\) coordinates
and the image of \(e\) spans the last. Thus \(h_w\in J_w\).
Conversely every such small isometry extends by the identity
and respects the displayed formulas. Surjectivity of the frames
also proves uniqueness of all the frame changes. This proves
the assertions, including the assertion about automorphisms.
\end{proof}

Use the coefficient representations and the normalized Fischer
forms of Section~\ref{sec:lowerlattice347}. Put
\(L=\bigotimes_\tau L_\tau\), \(s=\prod_\tau c_\tau^2\kappa_\tau\)
and take \(Z\) from Corollary~\ref{cor:normaldegree347}.
Theorem~\ref{thm:lowerlattice347} gives \(LL^*=sI\).
A coefficient table assigns to each class a value fixed by
its stabilizer acting on the coefficient space with the level
character. In its algebraic dual table one uses the dual
actions and inverse characters. Give the complex tables the
squared norms \(\sum_c\|Y_c\|^2/|\Gamma_c|\), where
\(\Gamma_c\) is the stabilizer in the relevant class set.
Using one representative per class is equivalent to using
equivariant functions on all decorated lattices.

Write \(\mathcal H_N\) for the full large coefficient table
and \(\mathcal H\) for the small one at the induced levels.
The map \(R:\mathcal H_N\to\mathcal H\) takes the value
\(ZL P(i(c))\) in the matching rational frames at \(c\).
It is well defined by the small block equivariance of \(L\)
and Proposition~\ref{prop:liftclasses380}. Let \(\mathcal T\)
be the small table with dual coefficient action. It pairs
bilinearly with \(\mathcal H\), using the same class masses.

Now choose the representatives \(j\) of the actual row classes
used in the period. Remove their frames at \(Q\), take the
\(J_w\)-orbits of their retained frames, and write \(c(j)\) for
the resulting class. This procedure is defined on the chosen
representatives for every row type. Choose a small transport from
the representative of \(c(j)\) to this object at the induced
levels. Its action on small output coefficients, including the
character, is denoted by \(U_j\). Thus the output evaluated
at the lift of \(j\) is \(U_j(RP)_{c(j)}\). For a row test
\(F\), view \(F_j\) as a linear functional in that frame, with
the normalized bilinear form used to take its values.
Write \(U_j^\vee F_j=F_j\circ U_j\), and let \(A_c\) be
averaging over \(\Gamma_c\) on dual coefficients. Define
\[
  (D F)_c=A_c\!\sum_{j:\,c(j)=c}
           \frac{|\Gamma_c|}{|\Gamma_j|}\,U_j^\vee F_j.
\]
The averaging sends the vector to the invariant dual table
and makes the expression independent of the choice of transports.

\begin{theorem}[Transfer of the prescribed terminal class sum]
\leavevmode\label{thm:transfer380}
For every column \(P\) in the large table and every test \(F\),
the finite period with the specified representatives satisfies
\[
 {\cal P}(P,F)=\langle RP,DF\rangle_{\mathcal C},\qquad
 RR^\dagger=|Z|^2sI.
\]
At a complex embedding, \(R^\dagger R/(|Z|^2s)\) is zero off
the classes \(i(c)\) and acts there as \(L^*L/s\) in the
matching frames. At the primes above \(g_1,g_2,T\), this
projection preserves the integral coefficient tables. The map
\(R\) identifies its projected full lattice with the full
integral small table, using the lattice frames of
Proposition~\ref{prop:normalframe349}. At these primes \(D\)
carries integral dual row values to integral dual values in the
induced-level table.
\end{theorem}
\begin{proof}
Pair an invariant value \((RP)_c\) against the displayed sum.
It pairs equally with a vector and its stabilizer average.
The factor \(|\Gamma_c|\) cancels the mass of \(c\), leaving
\(F_j(U_j(RP)_c)/|\Gamma_j|\) at each specified representative.
Their sum is the period of Theorem~\ref{thm:terminalformula347}.
Two choices of transport differ on the source by an element
of \(\Gamma_c\). This also proves the claimed independence after
averaging. Passage from this representative sum to a Haar integral
requires the contracted integrand to be invariant under the row level.

The large class \(i(c)\) has the same mass as \(c\), and the
different lifted classes are disjoint. In the Hermitian norms,
\(R^\dagger\) therefore inserts \(\bar Z L^*\) at \(i(c)\)
and zero elsewhere. The inserted values are fixed by the large
stabilizer by Proposition~\ref{prop:liftclasses380} and
equivariance. The two operator formulas follow. In rational
frames a right inverse inserts \(Z^{-1}L^*/s\). In lattice frames
the small frame is extended on the large dual lattice with
normal factor \(\tau(\bar\pi)^{-1}\) at \(\tau\).
The normal action on the branching component has exponent
\(\delta_\tau=|\lambda_\tau|-|\mu_\tau|\).
Consequently the frame change cancels the factor of \(Z\) in
restriction and the factor of \(Z^{-1}\) in insertion, just as in
Proposition~\ref{prop:normalframe349}. The matrices in these
paired integral frames are \(L\) and \(L^*/s\). Both preserve
their lattices by Theorem~\ref{thm:lowerlattice347}, since the
three primes exceed all the hook cutoffs. Their composition on
small values is the identity, and in the opposite order it is
the asserted projection.

For \(D\), the small lattice frame is attached to the underlying
lattice, and removing residue frames does not change it. Each
transport carries it to the frame lattice at the target since
its rational isometry carries the underlying lattice onto that
lattice. The finite character factors are roots of unity.
The stabilizers act faithfully on integral lattices of rank
\(n[E:\mathbb Q]\), by surjectivity of the residue frames.
If a prime \(p\) divides such a group order, an element of order
\(p\) has an eigenvalue of degree \(p-1\) over \(\mathbb Q\).
The degree bound of Theorem~\ref{thm:terminalformula347}
therefore makes every small mass denominator a unit at the
three stated primes. The formula for \(D\) proves integrality.
\end{proof}

There is a more precise answer when the retained row groups
\(H_w\) are contained in \(J_w\) with the matching dual
characters. This includes the ordinary row, its further
fixed-column subgroup, and the row with trivial first block,
paired with the ordinary terminal column. Indeed its first
block is upper unitriangular and its other diagonal blocks
carry the specified determinant characters. Restricting its
first-block character to each of the indicated subgroups gives
the row character in the dual pairing. At every extra place
\(v\in Q\) the row stabilizes an \(r\)-dimensional subspace
of \(k_v^n\), and its character is the specified character of
the determinant on the quotient, with the dual convention.
These groups and characters are given explicitly in
\srcloc{Formalism/LocalToGlobal/Columns/RowForms}{203}
and lines 217--233 there.

Fix \(c\). The possible refinements over its representative form
a finite set \(\Omega_c\), acted on by \(\Gamma_c\). At a retained
place one chooses a refinement of its frame orbit from \(J_w\)
to \(H_w\). At an extra place one chooses an \(r\)-subspace.
These choices are independent, since a residue frame is arbitrary
at each of these places. In particular
\[
 |\Omega_c|=\prod_{w\in S_g}[J_w:H_w]\,
                \prod_{v\in Q}|\operatorname{Gr}(r,n)(k_v)|.
\]
The fine classes over \(c\) are exactly the orbits in this set.
For \(y\in\Omega_c\), write \(\Gamma_y\) for its stabilizer.
On this group the transported fine dual action differs from
the coarse dual action by a finite scalar character \(\theta_y\).
It is the product of the extra determinant characters with
their indicated signs. This follows by cancelling the common
rational coefficient action and the matching retained
characters. Set
\[
  b_c=\#\{y\in\Omega_c:\ \theta_y|_{\Gamma_y}=1\}.
\]
In particular this is a count of individual refinements.

\begin{theorem}[Multiplicity and integral image of the row transfer]
\leavevmode\label{thm:fiberlength380}
On the dual coefficient space of each coarse class \(c\),
\(DD^\dagger=b_cI\). At every prime above \(g_1,g_2,T\)
the image of the full integral fine test table under \(D\)
is precisely the direct sum of the full integral coarse test
spaces at the classes for which \(b_c>0\). These assertions
use the representatives of the given class sum even when
a forgotten character is nontrivial. In particular the
archimedean multiplicities and the integral image are
independent of that choice. If the orders of all extra
characters are relatively prime to \(|\Gamma_c|\), then
\(b_c=|\Omega_c|\).
\end{theorem}
\begin{proof}
Choose one representative of each refinement orbit. Via
\(U_j^\vee\), its testing space is the fixed space for the
fine dual action of \(\Gamma_y\) in the coarse coefficient
space. The transport is unitary in the complex metrics. The
weight of the orbit in \(D\) is
\(m_y=|\Gamma_c|/|\Gamma_y|\). A fine invariant transforms
under the coarse action of \(\Gamma_y\) by \(\theta_y^{-1}\).
If this character is nontrivial, averaging over the subgroup
is zero. Averaging over \(\Gamma_c\) first averages over
\(\Gamma_y\), so this orbit contributes zero to \(D\).
If the character is trivial, its fine invariant space contains
the entire coarse invariant space. The operator \(A_c\) on
that fine space is orthogonal projection onto the coarse
invariants, since all these actions are unitary.

The ratio \(m_y\) cancels between the two class measures on
taking the adjoint. At a contributing orbit the adjoint hence
inserts the coarse value in the fine space, by inverse
transport. It vanishes at every other orbit. Applying \(D\)
once more sums the numbers \(m_y\) over the contributing
orbits, giving \(b_c\). Changing representatives transports
the stabilizer and its character isomorphically and makes
unitary changes in the coefficient spaces. The same computation
therefore applies to the prescribed section.

Locally at any of the three primes all \(m_y\) are units.
Given an integral coarse value at \(c\) with \(b_c>0\),
choose one contributing orbit, insert its inverse transport
divided by \(m_y\), and take zero at all other orbits.
This is an integral fine test and has the desired transfer.
Theorem~\ref{thm:transfer380} gives the reverse inclusion,
and the averaging calculation gives zero at all other classes.
Finally, the order of the character on any \(\Gamma_y\) divides
both its original order and \(|\Gamma_c|\). The stated
coprimality thus makes every orbit contribute.
\end{proof}

\begin{proposition}[A uniform estimate from the extra residue fields]
\leavevmode\label{prop:freeflags380}
Use the ordinary induced level with its upper unitriangular
first block, assume \(1\le r<n\), and put \(d=n[E:\mathbb Q]\),
\(M_d=|\mathrm{GL}_d(\mathbb F_3)|\). Let \(Q'\) consist of
the extra places of residue characteristic greater than \(d+1\).
Uniformly in the lattice class and in the extra characters,
\[
 \frac{b_c}{|\Omega_c|}\ \ge\
  1-(M_d-1)\prod_{v\in Q'}
                \min\!\left(1,\frac{2(r+1)}{|k_v|}\right).
\]
The bound also holds with common upper bounds for \(d,r\) on
the interval, taking \(Q'\) with the larger degree cutoff.
\end{proposition}
\begin{proof}
The faithful integral action of \(\Gamma_c\) embeds it in
\(\mathrm{GL}_d(\mathbb F_3)\) after reduction, as in the proof
of Theorem~\ref{thm:weylperiod351}. For completeness, a torsion
element in the kernel would give an element of prime order
\(p\) of the form \(1+3^aB\), \(a\ge1\), with \(B\) nonzero
modulo three. In its \(p\)-th power the leading term is
\(p3^aB\), which cannot be cancelled by the later binomial
terms, also when \(p=3\). Thus \(|\Gamma_c|\le M_d\).

A scalar in \(\Gamma_c\) is a root of unity, since the group is
finite. The first block at one retained auxiliary place forces
its reduction to be one. Its order is prime to that residue
characteristic by the eigenvalue degree bound in the preceding
theorem. For a root \(u\) of such order \(h\), the sum
\(1+u+\cdots+u^{h-1}\) reduces to the unit \(h\).
The identity \((u-1)(1+\cdots+u^{h-1})=0\) gives \(u=1\).
In fact every nonidentity \(\gamma\) in \(\Gamma_c\) is
nonscalar modulo every \(v\in Q'\). Otherwise conjugation by
\(\gamma\) on the endomorphisms of the localized lattice is a
matrix \(A\) congruent to the identity. Its order divides the
order \(h\) of \(\gamma\), a unit at \(v\) by the same degree
bound. The identity \((A-I)(I+A+\cdots+A^{h-1})=0\)
implies \(A=I\), since the second factor is invertible modulo
\(v\). Hence \(\gamma\) commutes with all endomorphisms over
the fraction field and is already a global scalar.

Over a field of \(q\) elements, let \(g\) be a nonscalar
invertible matrix of size \(n\). Some \(v\) and \(gv\) are
independent, since a map preserving every line is scalar by
application to the basis vectors and their pairwise sums.
Choose an \(r\)-space containing \(v\) and excluding \(gv\).
If \(X\) is its basis matrix, \([X\ \ gX]\) has rank at
least \(r+1\). Thus one of its minors is a nonzero polynomial
of degree \(r+1\) in the entries of a variable \(X\).
An invariant \(r\)-space forces this minor to vanish on every
one of its basis matrices. A nonzero polynomial of total
degree \(a\) vanishes on a proportion at most \(a/q\) of
the full affine array. To see this, write its degree in the
last variable as \(b\). By induction the leading coefficient
vanishes on a proportion at most \((a-b)/q\) of the earlier
variables. On the others there are at most \(b\) roots in the
last variable. The one-variable root bound starts the induction.
The probability that a uniformly chosen \(n\)-by-\(r\)
matrix has dependent columns is at most
\(\sum_{i=0}^{r-1}q^{i-n}\le1/(q(q-1))\). Conditioning on
independence therefore costs at most a factor of two, and the
resulting column spaces are uniform. The fraction of
\(g\)-invariant \(r\)-spaces is at most \(2(r+1)/q\).

For each nonidentity \(\gamma\in\Gamma_c\), independence of
the choices at \(Q'\) bounds its fixed proportion in
\(\Omega_c\) by the product in the statement. Summing those
proportions bounds the fraction having a nontrivial stabilizer.
Every refinement with trivial stabilizer contributes to \(b_c\)
for every choice of characters. This proves the inequality.
Increasing the degree and rank bounds weakens it uniformly
and proves the interval assertion.
\end{proof}

Here is a Haar integral for the same representative sum. Work at
the compatible retained levels of Theorem~\ref{thm:fiberlength380}
and let \(H\) denote the original compact open row level. Use the
double-coset measure of Theorem~\ref{thm:classtrace326}. The
ratio between the fine and the pulled-back coarse actions on
dual coefficients is a character \(\theta\) of the frame-change
group, trivial on rational isometries. At a fine stabilizer it
is \(\theta_y\) in transported coordinates. On output
coefficients the ratio is \(\theta^{-1}\).

\begin{corollary}[The character in a Haar realization of the finite sum]
\leavevmode\label{cor:rowhaar380}
There is a scalar function \(a\) on the finely framed objects,
with values zero or roots of unity in the coefficient field,
such that for every \(P,F\) as above,
\[
 \frac{1}{\operatorname{vol}(H)}
      \int \! F(h)\bigl(a(h)(RP)(h)\bigr)\,dh
                    ={\cal P}(P,F).
\]
Here the integral is over the definite small unitary quotient
with the weight at infinity contracted, and \(RP\) is pulled
back by forgetting the fine frames. If \(\theta\) is nontrivial
on the right compact level, the integral with \(a\) replaced
by the constant one is zero for every \(P,F\).

More generally suppose a candidate scalar kernel \(a'\) has
the same transformation law and let \(a'_j\) be its values
at the prescribed representatives. For any coarse output table
\(Y\), the fine table norms satisfy
\[
 \|(a'-a)Y\|^2
   =\sum_{c\in\mathcal C}\frac{\|Y_c\|^2}{|\Gamma_c|}
      \sum_{\substack{j:\,c(j)=c\\ \theta_j|_{\Gamma_j}=1}}
          \frac{|\Gamma_c|}{|\Gamma_j|}\,|a'_j-1|^2 .
\]
The first identity applies in particular to any earlier signed
column with values in \(\mathcal H_N\).
\end{corollary}
\begin{proof}
On a contributing orbit give the selected representative value
one and set \(a(gj)=\theta(g)^{-1}\) at a change of decorated
object by \(g\). This is well defined since \(\theta\) is one
on the stabilizer. On every other orbit set \(a=0\).
The resulting function is invariant under rational isometries
and cancels the character of the paired expression on right
level translates. It defines a scalar weight on the quotient
with the required transformation law. A noncontributing
orbit had zero contraction in the original representative
sum. Indeed on its stabilizer the coarse output and fine test
pairing transforms by a nontrivial scalar, while both class
values at the stabilized point are fixed in their respective
actions. Its contraction must therefore be zero. At each
contributing representative the inserted factor is one.
The contracted product is now right \(H\)-invariant. A
double coset at \(j\) has mass
\(\operatorname{vol}(H)/|\Gamma_j|\), by counting the
finite stabilizer in the compact quotient. This proves the
integral identity on summing the cosets.
Without the factor, right translation by any \(h_0\in H\)
multiplies the contracted integrand by its discrepancy
character (or its inverse with the opposite convention for
right translation). Haar invariance forces its integral to
vanish when that character takes a value different from one.

Any other scalar function with the same law must also vanish
on an orbit with nontrivial stabilizer character, by evaluating
at a point fixed by that stabilizer. On the other orbits, the
squared norm at the representative is
\(|a'_j-1|^2\|Y_{c(j)}\|^2\), since coefficient transports are
unitary. Summing with the fine masses proves the second identity.
The construction is integral at every finite place on coefficients,
since each of its nonzero scalar values is a root of unity.
\end{proof}

For any earlier signed column map \(\Theta\) with values in
\(\mathcal H_N\), the tested functional is exactly
\(\langle R\Theta(\cdot),DF\rangle_{\mathcal C}\).
Thus on unrestricted fine tests the transfer contributes the
explicit multiplicities above. For a specified projected test
space or lattice \(S\) its contribution is \(D(S)\).
To deduce the estimate of Section~\ref{sec:missing} one still
has to control the earlier signed image on the lifted classes
against these transported tests and identify its Gram form with
that of the designated analytic row.


\subsection{The signed terminal column on ramified residue flags}
\label{sec:ramifiedplanes383}

Fix the field \(E_x\), word, translate and auxiliary character of
Section~\ref{sec:missing}. Write \(Q\) for its extra marked places,
\(n=r+mt\) for the small terminal rank and \(a\) for the preceding
column rank. At \(v\in Q\) the local group is split, with residue
field \(k_v\) of order \(q_v\), and the character on residue units
is a nontrivial character \(\xi_v:k_v^\times\to J^\times\).
Here \(J\) is the coefficient field, containing the character
values. We use the character itself in the determinant below.
Write \(A=k_v^r\subset k_v^n\) for the first coordinate plane,
\(G_v=\GL_n(k_v)\) and \(P_v\) for its stabilizer. At a complex
embedding all the character values have absolute value one.

The preceding section treated restriction of a column whose extra
frames had already been forgotten. We now evaluate the terminal
kernel of the signed arithmetic chain at those frames. A basis
of the small target at \(v\) means an isomorphism from the
coordinate space to the residue space of its lattice. Changing
this basis by \(h\in G_v\) means composition on the right by
\(h\). Thus \(hA\) is the new marked plane in the original
coordinates. Set
\[
       \lambda_v(p)=\xi_v\!\left(\det(p\mid k_v^n/A)\right)
       \qquad(p\in P_v).
\]

\begin{proposition}[The flag and phase of a terminal matrix]
\leavevmode\label{prop:termplane383}
Fix a summand at the last edge, including its source frame and
coupling vector coset. Let \(B:k_v^a\to k_v^n\) be the reduction
of the small projection of its matrix in the reference bases,
when this projection is integral. If it is nonintegral, or if
\(B\) has rank less than \(r\) on the first \(r\) coordinates,
the summand vanishes for every target basis. Otherwise put
\(I=B(k_v^r)\). Apart from one factor independent of \(h\),
its value as a function of the target basis is
\begin{equation}
 b_B(h)=\mathbf 1_{hA=I}\,\xi_v(\det h)\,
          \xi_v\!\left(\det((h^{-1}B)|_{k_v^r}:k_v^r\to A)\right).
 \label{eq:flagphase383}
\end{equation}
The determinant on the right is evaluated only on the support
of the indicator. For \(p\in P_v\) one has
\(b_B(hp)=\lambda_v(p)b_B(h)\).
\end{proposition}
\begin{proof}
Let \(x\) be the matrix from the preceding lattice to the large
terminal lattice \(M\perp\mathcal O_E e\). Its small projection
is projection along \(e\) followed by expression in the basis
of \(M\). At this step the index is \(m-1\), so the number
\(m-1-(m-1)\) of auxiliary columns in the source is zero.
The target already has zero auxiliary columns. Consequently
the reordered entries in the terminal position test are the
ordinary entries of this projection. The test requires them all
to be integral. It requires the last \(n-r\) entries in each
of the first \(r\) columns to vanish modulo \(v\), and weighs
the first \(r\)-by-\(r\) minor by \(\xi_v\), extended by zero
on nonunits. These statements follow by reading the defining
position test in
\srcloc{Formalism/LocalToGlobal/Columns/FiniteSumNonzeroLargeDepth}{711}
and its terminal arguments on lines 1141--1170. An invertible
integral basis change preserves integrality in both directions.
On the reduction it changes the matrix to \(h^{-1}B\).
The two remaining requirements say exactly that the first
\(r\) columns give a basis of \(hA\). This proves the indicator
and the minor factor.

For the other factor, the coefficient of a matrix coset in the
theta distribution is its derivative operator multiplied by a
source determinant character and the inverse target determinant
character. The derivative operator uses the two lattices and
the matrix coset. Its arguments contain no residue bases.
The target determinant factor before inversion, on multiplying
the target basis by an integral unit matrix \(u\), changes by
\(\xi_v^{-1}(\det\bar u)\). Indeed its local character is
defined by the uniformizer power times the inverse residue
character of the unit part. Multiplication by a unit preserves
the power and multiplies the residue unit by its reduction.
The extension of the small basis to the large target sends
the last basis vector to \(\pi e\) in residue coordinates,
as follows from the frame formula in
Proposition~\ref{prop:liftclasses380}. Its basis change has
residue \(\operatorname{diag}(h,1)\), hence determinant
\(\det h\). Any two integral lifts of that residue change
have determinant ratio one modulo \(v\). Inverting the target
factor thus supplies precisely \(\xi_v(\det h)\).
The determinant and derivative formulas are in
\srcloc{Formalism/LocalToGlobal/KernelLattices/ThDistC}{123}
and lines 137--145 there. The local character is specified in
\srcloc{Formalism/LocalToGlobal/Columns/FiniteSumNonzeroLargeDepth}{935}.

The last covector, the additive Fourier test on the coupling
coset, and the finite masses at the preceding vertices use
the original lattices and source bases. None uses the
small target residue basis at \(v\). This accounts for every
remaining factor in the summand. Finally on replacing \(h\)
by \(hp\), the minor acquires the inverse determinant of
\(p|_A\). Multiplication by \(\det p\) leaves the determinant
on the quotient. This proves the transformation law.
\end{proof}

Fix representatives \(c\) of small lattice classes retaining the
actual row levels at the original two auxiliary primes and forgetting
only the extra frames at \(Q\). Let \(\Gamma_c\) be the stabilizer.
Choose one reference basis at each extra place in each class.
Put
\[
 \Omega=\prod_{v\in Q}G_v/P_v,\qquad
 d_Q=|\Omega|,\qquad G_Q=\prod_{v\in Q}G_v.
\]
For each tuple of planes \(I\in\Omega\), choose
\(h_I=(h_{I,v})\) carrying the standard planes to \(I\).
For a fixed arithmetic input \(\phi\), expand the signed
theta chain at \(c\). A term \(\gamma\) indexes its preceding
classes, matrices, finite types and coupling coset. Discard
terms failing the integrality and rank tests above. Denote the
tuple of surviving images by \(I_\gamma\). Let \(C_{\gamma,c}\)
be the coefficient vector with the factors in
\eqref{eq:flagphase383} removed. It includes the earlier signed
tests, all preceding masses, the Fourier factors, the derivative
operators, and the last branching map. This map is independent
of the extra bases. Thus the finite array
\begin{equation}
 Y_{\phi,c}(I)=
    \sum_{\gamma:I_\gamma=I} C_{\gamma,c}
                       \prod_{v\in Q} b_{B_{\gamma,v}}(h_{I,v})
 \label{eq:flagarray383}
\end{equation}
uses the actual coefficient in its chosen basis. In particular
a term can occur in just one entry of this array. If \(Q\) is
empty the array has one entry.

\begin{theorem}[The Haar pairing and the square norm on the fine row]
\leavevmode\label{thm:planehaar383}
Let \(F\) be a test in the prescribed dual row space, including
any imposed projection on that space. Write \(F_c(h_I)\) for
its value, viewed as a functional on the small coefficient
space. The signed arithmetic period is
\begin{equation}
 {\cal P}_{\rm chain}(\phi,F)
  =\sum_c\frac{1}{|\Gamma_c|}
       \sum_{I\in\Omega} F_c(h_I)\big(Y_{\phi,c}(I)\big)
  =\frac{1}{\operatorname{vol}(H)}
       \int F(h)\big(Y_{\phi}(h)\big)\,dh .
 \label{eq:flaghaar383}
\end{equation}
Here \(H\) is the original compact open row level, infinity
is contracted in the same coefficient pairing, and \(Y_\phi\)
is the signed arithmetic column after branching. Its squared
norm in the mass metric of this fine level is
\[
 \|Y_\phi\|^2
     =\sum_c\frac1{|\Gamma_c|}
            \sum_{I\in\Omega}\|Y_{\phi,c}(I)\|^2 .
\]
With probability measure on the extra full residue bases, their
mean square in any fixed class is
\(d_Q^{-1}\sum_I\|Y_{\phi,c}(I)\|^2\).
\end{theorem}
\begin{proof}
The class sum for the signed chain pairs its evaluated terminal
column with the original test. The coefficients at all preceding
vertices are independent of the last row basis. Expanding those
finite sums and applying Proposition~\ref{prop:termplane383}
at each extra place gives \eqref{eq:flagarray383}. Finiteness
here follows from the finite lattice quotients for each matrix
and coupling vector, and from finiteness of definite classes
at the indicated compact levels. For the former one may take
local bases and reduce modulo the common product of the prescribed
prime powers, as in Section~\ref{sec:covectorpaths352}. The
coupling cosets are quotients of two local lattices of the same
rank. For the class number assertion, the underlying integral trace
lattice is positive definite of fixed determinant, so it has
finitely many isometry classes by Minkowski reduction. On
any one of them multiplication by each element \(a_0\) of a
fixed integral basis of \(E_x\) has operator norm bounded by
\(\max_\tau|\tau(a_0)|\), by the orthogonal decomposition at
infinity. Only finitely many integral matrices satisfy this
bound. The trace pairing on the field recovers the Hermitian
form from its traces after these multiplications. Residue
frames add only finitely many choices.

Each test has change-of-basis character \(\lambda_v^{-1}\) in
this basis convention. To see the sign, the frame in the row
definition is the map \emph{to} the coordinate space. Right
multiplication of its inverse basis by \(p\) acts on this map by
left multiplication by \(p^{-1}\). The character acting on tests
for that left operation is the determinant character on the
quotient, since the row uses the inverse of the level character
specified on line 231 of
\srcloc{Formalism/LocalToGlobal/Columns/RowForms}{230}
in the coefficient action on line 269. Evaluating at
\(p^{-1}\) gives \(\lambda_v(p)^{-1}\). Proposition~\ref{prop:termplane383}
therefore makes the contraction independent of the choices of
\(h_I\). A rational isometry transports both coefficients and
preserves the branching contraction. Equivalently one can
reindex the finite theta sum by transporting each matrix with
its two endpoint lattices. It preserves the matrix cosets, the
coupling cosets and all their masses. The integrand thus has
one value on each \(\Gamma_c\)-orbit in \(\Omega\). The fine
classes over \(c\) are these orbits. An orbit with stabilizer
\(\Gamma_{c,I}\) has \(|\Gamma_c|/|\Gamma_{c,I}|\) elements.
This gives the first equality with the correct fine masses.
The Haar mass of its double coset is
\(\operatorname{vol}(H)/|\Gamma_{c,I}|\) in the normalization
of Theorem~\ref{thm:classtrace326}, proving the second equality.
The same count applies to the squared norm, since all level
characters and transports are unitary. Each plane tuple has
\(\prod_v |P_v|\) bases, giving the asserted probability mean.
\end{proof}

The frame calculation also gives an integral projection in the
global row space with full residue frames at the places of \(Q\).
Let \(R(g)Y(h)=Y(hg)\) be right translation. The finite groups
at different places commute. The mass metric, equivalently the
Haar metric with any common scaling, is invariant under this
action. Take a \(J\)-subspace \(V\) of the actual column space
with the standard character \(R(p)=\prod_v\lambda_v(p_v)\)
on \(V\). For example one can take the span of any of the signed
columns in this row with fixed extra characters. Put
\[
       W=\sum_{I\in\Omega}R(h_I)V,\qquad b=\min(r,n-r).
\]
These translates describe the span inside the given fully
framed space. The choice of representative of \(I\) does
not change its summand.

\begin{theorem}[Orthogonal flag summands and an integral projector]
\leavevmode\label{thm:flagproject383}
Assume \(0<r<n\). The sum defining \(W\) is orthogonal at
each complex embedding. Its Gram matrix on the translates
of any fixed list in \(V\) consists of \(d_Q\) copies of the
Gram matrix of that list. For each \(v\in Q\) choose a
unit \(t_v\in k_v^\times\) such that \(z_v=\xi_v(t_v)\)
generates the image of \(\xi_v\). Define
\[
    D_Q=\prod_{v\in Q}
              q_v^{\binom n2}(z_v-1)^{b}\in\mathcal O_J
              \setminus\{0\}.
\]
On \(W\), multiplication of the projection onto any one
summand by \(D_Q\) is the action of an element of
\(\mathcal O_J[G_Q]\). The assertions still apply after
a projection on the full row space that commutes with \(G_Q\),
using the resulting image of \(V\).
\end{theorem}
\begin{proof}
First work at one place and omit its subscript. Write
\(V_I=R(h_I)V\). An element fixing \(I\) acts on \(V_I\)
by \(\xi\) of its determinant on \(k^n/I\), by conjugation.
For distinct planes \(I,I'\), choose a line in \(I\) outside
\(I'\). There is a hyperplane containing \(I'\) complementary
to that line, whose intersection with \(I\) complements it in
\(I\). For instance take the kernel of a functional vanishing
on \(I'\) and nonzero on the line. Multiplication by \(t\) on
the line and by one on this hyperplane preserves both planes.
Its determinants on the quotients by \(I,I'\) are respectively
one and \(t\). The two spaces thus have different eigenvalues
for a unitary operator and are orthogonal. Right translation
preserves the Gram matrix within a single space. On a product
of places two different tuples differ at some place, giving
the same orthogonality. This proves the first assertions.
The argument uses \(V\) as a multiplicity space and imposes
no independence on its chosen list.

For the projector at one place take \(U\) to be the upper
unitriangular subgroup in the reference basis. Its order is
\(q^{\binom n2}\). Its orbits on the Grassmannian consist of
the subspaces with a fixed subset \(J_0\) of \(r\) pivot positions,
using the last nonzero entry as pivot in each echelon column.
Indeed the coordinate plane with these positions can be sent
to any plane in the cell by using its pivot-normalized columns
as the corresponding columns of an upper unitriangular matrix.
The other columns can be chosen to be the original coordinate
vectors. The dimensions of intersection with the initial
coordinate spaces are invariant under \(U\), so different
pivot sets lie in different orbits.

Let \(E_U=|U|^{-1}\sum_{u\in U}R(u)\).
In the direct sum over a cell the \(U\)-invariant vectors are
determined by their component at its coordinate plane. The
stabilizer there acts trivially on that component, since its
quotient action is unipotent and has determinant one.
Every value at the coordinate plane therefore extends uniquely
to a \(U\)-invariant vector over the cell, by translation.
A diagonal matrix normalizes \(U\) and fixes that plane.
On these invariant vectors the diagonal matrix \(t_i\)
having entry \(t\) only in position \(i\) and ones elsewhere
acts by \(z\) if \(i\notin J_0\), and by one if \(i\in J_0\).
This follows by evaluating the translate at the coordinate
summand and using its quotient determinant.

Let the standard pivot set be \(A_0=\{1,\ldots,r\}\).
When \(r\le n-r\), apply to the image of \(E_U\) the operator
\[
              \prod_{i\in A_0}\frac{z-R(t_i)}{z-1}.
\]
It retains the cell with every \(i\in A_0\) a pivot. There
is only one such pivot set, and its cell consists of the
standard plane alone. When \(n-r<r\), use instead
\[
              \prod_{i\notin A_0}\frac{R(t_i)-1}{z-1}.
\]
This requires every index outside \(A_0\) to be a nonpivot,
with the same consequence. On the standard summand \(E_U\)
and the selected product both act as the identity. Every other
cell is killed. Thus their composite is the required projection.
Clearing the displayed denominators and \(|U|\) produces an
integral combination of group elements. In the product case
one applies this one-place operator with all other places fixed,
and multiplies the commuting projections. Conjugation by
\(h_I\) gives the projection onto the desired translate.
A projection commuting with \(G_Q\) carries a standard
character vector to a standard character vector, so the entire
proof applies to its image.
\end{proof}

\begin{corollary}[Cost for the integral row lattice]
\leavevmode\label{cor:flagcost383}
Let \(\Lambda\) be a \(G_Q\)-stable full \(\mathcal O_J\)-lattice
in the full-frame row space. Inside \(W\) put
\[
 \Lambda_{\rm sat}=\Lambda\cap W,\qquad
 \Lambda_{\rm sep}=\bigoplus_{I\in\Omega}R(h_I)(\Lambda\cap V).
\]
Then \(D_Q\Lambda_{\rm sat}\subseteq\Lambda_{\rm sep}
\subseteq\Lambda_{\rm sat}\). For \(d=\dim_J W\) their
index satisfies
\begin{align*}
 \frac{\log[\Lambda_{\rm sat}:\Lambda_{\rm sep}]}{[J:\QQ]}
 &\le d\sum_{v\in Q}\left\{\binom n2\log q_v+
       b\frac{\log\Phi_{e_v}(1)}{\varphi(e_v)}\right\} \\
 &\le d\left(\binom n2+b\right)\log N_{E_x/\QQ}\mathfrak f_\xi .
\end{align*}
Here \(e_v=\operatorname{ord}(\xi_v)>1\), \(\Phi_j\) is the
cyclotomic polynomial, and \(\mathfrak f_\xi\) is the finite
conductor of the auxiliary character on \(E_x\). At a coefficient place
whose residue characteristic divides none of the \(q_v\), the two lattices coincide unless
at least one \(e_v\) is a power of its residue characteristic.
These conclusions apply to the full-frame integral test tables
and their duals, and to intersections with row projections at
the retained primes.
\end{corollary}
\begin{proof}
For each vector of \(\Lambda\cap W\), its component in the
\(I\)-summand multiplied by \(D_Q\) lies in \(\Lambda\) by the
group-ring formula. The intersection of \(\Lambda\) with that
summand is exactly \(R(h_I)(\Lambda\cap V)\), since the
translation and its inverse preserve \(\Lambda\). Summing the
components gives the inclusions. The quotient is locally a
quotient of \(d\) copies of \(\mathcal O_J/(D_Q)\), since a
full lattice is locally free over each discrete valuation ring.
Multiplying the local index bounds gives
\([\Lambda_{\rm sat}:\Lambda_{\rm sep}]\le|N_{J/\QQ}D_Q|^d\).

The conjugates of \(z_v\) in its cyclotomic field are the
primitive \(e_v\)-th roots. This gives the first logarithmic
bound by taking norms. For \(e>1\), the identity
\(\prod_{j\mid e,\,j>1}\Phi_j(1)=e\) follows by factoring
\((X^e-1)/(X-1)\) and setting \(X=1\). Induction on \(e\)
gives \(\Phi_e(1)=p\) when \(e\) is a power of a prime \(p\),
and one otherwise. Also
\(\log\Phi_e(1)/\varphi(e)\le\log2\) since every conjugate
of \(1-z_v\) has absolute value at most two. Each extra
prime divides the conductor since its character on units is
nontrivial. Its contribution to the norm of the conductor
is at least \(q_v\). These observations give the second
inequality and the local assertion.

For the integral coefficient tables the lattice attached to a
small object depends on its Hermitian lattice. After adjoining
the last vector its coefficient lattice depends on the lattice
and the terminal covectors. A change of the extra residue bases
fixes these data, permutes the entries and preserves the full
integral table. The definitions occur in
\srcloc{Formalism/LocalToGlobal/Columns/CoefCol}{89}
and lines 93--103 there. Dualizing preserves this stability.
Frame changes or averages at the two retained primes commute
with \(G_Q\), since they act on separate local factors. Their
rational images intersected with the integral table are
therefore stable as required. This use of intersections allows
a projector with a nonunit denominator.
\end{proof}

In the fixed standard row of \eqref{eq:flaghaar383}, cancellation
can still occur among the paths with the same tuple of image
planes. When that row is paired with a specified analytic
functional, the values of the functional on the surviving entries
also enter. The flag projection computes the integral loss from
separating translates. Section~\ref{sec:missing} still needs
the Gram estimate for the specified words and the comparison to
the analytic Bessel periods on the same interval.


\subsection{The common vector at the two marked primes}
\label{sec:coupledgauss386}

Keep a field, word, translate and choice of sign from
Section~\ref{sec:ramifiedplanes383}. Let \(S_g\) be the marked
places over \(g_1,g_2\). These primes split completely.
Fix \(w\in S_g\), put \(O=\mathcal O_{E_w}\), and write
\(k=O/qO\), \(q=|k|\). Thus the rational integer \(q\) is
a uniformizer here. Fix all the class representatives, matrix
cosets and type indices on one path. Leave unsummed the vector
shared by the last two edges. We compute its contribution
before the signed paths are added.

At the penultimate edge put \(n_0=N(m-2)\), \(d=N(m-1)\)
for the two core dimensions. There is one extra source coordinate
and no extra target coordinate. In the chosen source and target
frames let \(X\) be the \(d\)-by-\(n_0\) core matrix and \(y\)
the extra column. In the terminal edge write \(z\) for the row
covector obtained by composing the last matrix with the last
coordinate divided by \(\pi_T\), in the same intermediate frame.
Here \(\pi_T\) is the specified generator at the terminal prime.
Its local Fourier factor is zero unless \(z\in O^{1\times d}\).
When \(z\) is integral, its factor is \(\psi(-zy)\), where
\(\psi\) has conductor \(O\). The shared vector ranges over
\(q^{-1}O^d/O^d\). Write \(\theta(t)=\psi(\tilde t/q)\)
for the nontrivial additive character on \(k\), using any
integral lift. Then \(v=\overline{qy}\in k^d\) and this factor
is \(\theta(-\bar z v)\). Additive Haar measure with
\(\operatorname{vol}(O^d)=1\) gives each vector coset mass one.
These are the two marked factors in
\srcloc{Formalism/LocalToGlobal/Columns/FiniteSumNonzeroLargeDepth}{1131}
and lines 1149--1170 there.

Use indices starting at zero. Let \(h\) be the penultimate
source representative in its full \(n_0+1\) coordinates.
It preserves \(S O^{n_0+1}\), where
\(S=\operatorname{diag}(1,\ldots,1,q)\).
This is the specified marked-frame condition, expressed on the
lattice before the columns are reordered. Write \(P\) for the
permutation matrix inserting the last coordinate at its marked
position. Then \(U=S^{-1}hSP\) belongs to
\(\mathrm{GL}_{n_0+1}(O)\). With \(C=[X\mid 0]U\) and \(b\)
the last row of \(U\), the tested matrix is
\[
     [X\mid y]hSP=[X\mid qy]U=C+(qy)b .
\]
The determinant multiplier from the source action on the
Schr\"odinger test is denoted by \(c_h\). It is the original
splitting character evaluated on \(\det h\), since \(h\)
preserves a full lattice and hence \(\det h\) has valuation
zero. All values of this character on units are roots of unity.

Put \(b_0=m-2\), \(a=r+b_0t\), and, for \(0\le i<b_0\), let
\(B_i=\{r+it,\ldots,r+(i+1)t-1\}\) be the diagonal blocks
after the initial \(r\) coordinates. The rank formulas are
\(N(j)=k_0+j(t+1)\) and \(r=k_0+m-1\), so
\[
       a=N(m-2)+1=n_0+1,\qquad d=N(m-1)=a+t .
\]
Thus all source columns enter the penultimate test.
Write \(A_j\subset k^d\) for the first \(j\) coordinate
positions, and \(s_i=r+(i+1)t\).
The ranks are specified in
\srcloc{Formalism/LocalRingCardinality/Params}{32}.
Write \(\chi_i\) for the
residue determinant characters and
\(\nu_i=\chi_i\chi_{b_0}^{-1}\). Extend every multiplicative
character, including the trivial one, by zero at zero.

\begin{proposition}[The support and the varying determinant block]
\label{prop:markedcell386}
If any entry in the first \(a\) columns of \(C\) is nonintegral,
the local summand vanishes for every \(v\). Otherwise apply
the entry congruences of the marked test to the residue matrix
\(\bar C+v\bar b\). Their solution set \(D\subset k^d\) is
empty or is given by fixing a subset of the individual
coordinates. Choose \(v^0\in D\) when it is nonempty, write
\(D=v^0+W\), and put \(f=\dim_k W\).

On \(D\) the penultimate test, with \(c_h\) removed, equals
its prescribed additive character in the superdiagonal
entries times
\[
                 \prod_{i=0}^{b_0-1}\nu_i(d_i+L_i v'),
                 \qquad v'=v-v^0\in W .
\]
Here \(d_i\) is the determinant of the \(B_i\) diagonal block
at \(v^0\), and \(L_i\) is a linear form on its free row
coordinates, extended by zero on the other free coordinates.
At most one \(L_i\) is nonzero. All constants and linear
forms in this description are computed from \(C,b,v^0\).

Suppose that the test is nonzero at some position. Let \(H\)
be the image of the first \(a\) rows of \(\bar X\) in \(A_a\).
It is a hyperplane. Let \(\lambda_X\) be a nonzero form on
\(A_a\) with kernel \(H\). If the first nonzero entry of
\(\bar b\) occurs among the initial \(r\) columns, all
determinants are constant on \(D\). If it occurs in \(B_i\),
then \(W=A_{s_i}\), the \(i\)-th determinant varies, and
\(L_i\) is a nonzero multiple of \(\lambda_X|_{A_{s_i}}\).
On this space it is supported on the \(B_i\) coordinates.
\end{proposition}
\begin{proof}
The equality of tested matrices follows by multiplying out
the source dilation and permutation. It uses the identity
target representative in the penultimate edge. Since \(U\)
is integral, adding \((qy)b\) on the prescribed window cannot
change failure of integrality of \(C\). Once integrality
holds, every required residue entry is of the form
\(\bar C_{ij}+v_i\bar b_j\).

We detail the congruences used in the source's two tests.
Rows with index at least \(a\) vanish on the first \(a\)
columns. The leading \(r\)-by-\(r\) square is upper
unitriangular in the plain test. In the symmetrized test
it is upper unitriangular after ordering these indices by
the permutation exchanging the two initial consecutive groups
of sizes \(k_0,r_S\). Its additional zero entries in column
\(k_0+r_S\) at rows \(i<k_0\) are imposed whenever that column lies in the
initial \(r\) columns. In both tests the rows in a later
block vanish on the first \(r\) columns and on columns of
every earlier block. These are entrywise zero or one equations.
Each constrains just one variable \(v_i\), either fixing its
value, giving no restriction, or being inconsistent. Intersecting
all of them gives the assertion about \(D\). The lists, including
the additive sums in the two orders, are in
\srcloc{Formalism/LocalToGlobal/Columns/FiniteSumNonzeroLargeDepth}{771},
lines 780--810 there.

In the leading \(a\)-by-\(a\) matrix the block below the
initial square is zero, and the submatrix on the remaining
blocks is upper block triangular. The initial square has determinant
one, also after its simultaneous row and column permutation.
It follows that the full determinant is the product of the
\(B_i\) determinants. Multiplying its inverse character
\(\chi_{b_0}^{-1}\) by the prescribed individual characters
gives the displayed \(\nu_i\). If a determinant vanishes,
both versions of the original expression vanish, since the
unit-character factor of a nonunit is zero.

In block \(i\) the row changes at \(v^0\) are scalar multiples
of the same row \(\bar b|_{B_i}\). By multilinearity of the
determinant, every term taking changes from two rows vanishes.
The coefficient of \(v'_j\) for a free row \(j\in B_i\) is
the determinant obtained by replacing row \(j\) of the base
block by \(\bar b|_{B_i}\). This computes \(L_i\).
The last row of the invertible matrix \(\bar U\) has a
nonzero entry. Take the earliest of the initial
group and subsequent blocks in which its restriction is
nonzero. All earlier blocks have zero replacement rows.
In each later block every row is constrained to vanish in
the nonzero column of that earliest group. Hence the
variable in that row is fixed throughout \(D\). There are
no free rows in these later blocks. Only the earliest group
can supply a varying determinant, and the initial group
has constant determinant one.

If the earliest group is \(B_i\), none of the entry equations
on a row of index less than \(s_i\) uses a nonzero coefficient
of \(v\). Indeed the equations on the first \(r\) columns have
\(\bar b_j=0\), and a block column can constrain only rows
in strictly later blocks or below the whole square. Columns
before \(B_i\) have zero coefficient as well. Every row of
index at least \(s_i\) is fixed by a column of \(B_i\).
This gives exactly \(W=A_{s_i}\). At a position where the
test is nonzero, the \(B_i\) block is an invertible matrix
\(M_i\). The linear part of its determinant in the free block
rows has coefficient row \(\det(M_i)\bar b|_{B_i}M_i^{-1}\),
which is nonzero. This follows either from the row replacements
above or by multiplying the rank-one update on the right
by \(M_i^{-1}\).

Since all columns are tested and \(U\) is invertible over
\(O\), integrality of \(C\) implies integrality of \(X\).
Write \(\bar X^{(a)}\) for its first \(a\) rows after reduction.
The top determinant of the tested matrix is
\[
       \det(\bar U)\det[\,\bar X^{(a)}\mid v_{<a}\,].
\]
At a position passing all tests this is nonzero. Thus
\(\bar X^{(a)}\) has rank \(a-1\), and the second determinant
as a linear function of its last column has kernel exactly
\(H\). Also all rows of \(\bar X\) below row index \(a-1\)
vanish, since the tested matrix vanishes there modulo \(q\)
and \(\bar U\) is invertible. On \(D\) this same top
determinant is the product of the diagonal determinants.
For the earliest group \(B_i\), the other factors are
nonzero constants on \(D\). Taking linear parts on
\(A_{s_i}\) proves the assertion about \(\lambda_X\).
\end{proof}

\begin{theorem}[The marked Fourier coefficient of the coupled edges]
\label{thm:markedgauss386}
Let \(\mathcal F_w(z)\) be the sum over the common vector of
the penultimate local test times its terminal Fourier factor.
The top matrix test at the last edge is independent of this
vector and is kept as a separate multiplier. Write \(\Phi_X(y)\)
for the penultimate local factor, including \(c_h\). Then, with
the measure fixed above, one has the equality of local integrals
\[
       \mathcal F_w(z)=\int_{E_w^d}\Phi_X(y)\psi(-zy)\,dy.
\]
Set
\(\mathcal F_w(z)=0\) for a failed integrality test, an empty
\(D\), or a singular constant determinant block. In the other
cases there is a root of unity \(\zeta\) and a form
\(\omega\in W^\vee\) prescribed by the data as follows.
Write the product of the additive superdiagonal character and
\(\theta(-\bar z v)\) on \(v^0+W\) as its value at \(v^0\)
times \(\theta(\omega(v'))\). Multiply this value by \(c_h\)
and the \(\nu_i(d_i)\) of the constant blocks to obtain \(\zeta\).

If all determinants are constant, \(\mathcal F_w(z)\) is
\(\zeta q^f\) for \(\omega=0\) and zero otherwise. If a
determinant varies, write \(d_*,L,\nu\) for its constant,
nonzero linear form and character. Then \(\mathcal F_w(z)=0\)
unless \(\omega=sL\) for \(s\in k\). For this frequency
one has the exact formula
\[
 \mathcal F_w(z)=
     \zeta\,q^{f-1}\theta(-s d_*)G_s(\nu),\qquad
 G_s(\nu)=\sum_{a'\in k^\times}\nu(a')\theta(sa').
\]
For nontrivial \(\nu\) this is nonzero precisely when \(s\ne0\)
and has absolute value \(q^{f-1/2}\). For trivial \(\nu\) the
Gauss factors are \(q-1\) at \(s=0\) and \(-1\) at \(s\ne0\).
These absolute values hold at every complex embedding.
\end{theorem}
\begin{proof}
Because every source column is tested, integrality of the
tested matrix \([X\mid qy]U\) is equivalent to integrality
of both \(X\) and \(qy\). Hence the integrand is supported in
\(q^{-1}O^d\). The penultimate factor is invariant under
translation by \(O^d\), since such translations preserve
every tested residue. For integral \(z\) the Fourier factor
is constant on each of the same cosets, of mass one. For
nonintegral \(z\) its integral on each coset vanishes by
translation, since \(\psi\) has conductor \(O\). These give
respectively the finite sum and the zero specified above.

The additive exponent for the plain test is the sum of
entries \(M_{i,i+1}\), \(0\le i<r-1\). In the symmetrized
test it is the sum within each of the two specified strings
\(0,\ldots,k_0-1\) and \(k_0,\ldots,r-1\).
Each entry is affine linear in \(v\). Any additive character
of \(k\) equals \(u\mapsto\theta(cu)\) for a unique \(c\in k\).
Indeed these characters are distinct and there are \(|k|\)
of them. This proves the prescription of \(\omega\), also for
the two different choices of superdiagonal character.

When all determinants are constant, averaging a nontrivial
additive character over the vector space \(W\) gives zero by
translation, and the trivial character sums to \(q^f\).
Otherwise sum on fibers of \(L\). The sum vanishes when the
character is nontrivial on \(\ker L\). Triviality there is
equivalent to \(\omega=sL\), since each nonzero \(k\)-linear
map to \(k\) is surjective. The fiber \(d_*+Lv'=a'\)
then contributes \(q^{f-1}\nu(a')\theta(s(a'-d_*))\).
Summing proves the formula with the stated sign.

For nontrivial \(\nu\), summation over the multiplicative
group gives \(G_0=0\). For \(s\ne0\), substitution gives
\(G_s=\nu(s)^{-1}G_1\). Expanding the squares and summing
first over \(s\) gives \(\sum_s|G_s|^2=q(q-1)\) by additive
orthogonality. Hence \(|G_1|=\sqrt q\). For the trivial
character, direct summation of the nonzero additive phases
gives the remaining two values. Applying a complex embedding
permutes the roots of unity, preserves triviality of a
character and preserves the same orthogonality calculation.
\end{proof}

\begin{corollary}[Application to the signed column and its local Fourier tables]
\label{cor:coupledglobal386}
For each path in the flag array of Theorem~\ref{thm:planehaar383},
fix all choices other than the common vector at the places of
\(S_g\). The sum over these components of the vector is obtained
by inserting \(\prod_{w\in S_g}\mathcal F_w(z_w)\) in the
original signed coefficient in place of their two
vector-dependent factors. This formula includes the base
tuple \(Q=\varnothing\). For a path on which the cells are
nonempty, write
\[
     D_g=\prod_{w\in S_g}q_w,\qquad
     K_g=\prod_{w\in S_g}q_w^{-f_w}\mathcal F_w(z_w).
\]
Then \(D_g K_g\) is an algebraic integer and at every complex
embedding \(\sigma\) a nonzero \(K_g\) satisfies
\[
                 D_g^{-1}\le |\sigma K_g|\le1,\qquad
                 \log D_g\le [E_x:\QQ]\log(g_1g_2).
\]
The factor \(\prod_w q_w^{f_w}\) is retained individually
when adding the paths.
For the marked characters prescribed on the common intervals,
let \(D_{\rm var}\) be the product of \(q_w\) over the places
with a varying determinant. If \(K_g\ne0\), then
\[
       |\sigma K_g|=D_{\rm var}^{-1/2}
       \quad\text{at every embedding},\qquad
       D_{\rm var}K_g\ \text{is an algebraic integer}.
\]

For the full local frequency table at one marked place, assuming
the cell is nonempty and the constant determinants are units, one has
\[
   \sum_{\bar z\in(k^d)^\vee}|\mathcal F_w(z)|^2
       =q^{d+f}(1-q^{-1})^\delta ,
\]
where \(\delta=1\) if a determinant varies and is zero otherwise.
For a nonzero table with varying determinant in \(B_i\) and
nontrivial \(\nu_i\), its frequency support is exactly
\[
       \bar z|_{A_{s_i}}\ \in\ k^\times
                    \bigl(\lambda_X|_{A_{s_i}}\bigr).
\]
Thus on the first \(s_i\) coordinates the covector vanishes
on \(H\cap A_{s_i}\) and is nonzero on \(A_{s_i}\). There
is no restriction on its later coordinates. For two such
tables in one intermediate frame the supports are disjoint
if their varying blocks differ or, for the same block, if the
two hyperplanes \(H\cap A_{s_i}\) differ. Their inner product
then vanishes for every nonnegative weight on the residue
covectors, including weights supported on prescribed
terminal matrices.
\end{corollary}
\begin{proof}
At the intermediate lattice the global common vector ranges
over the quotient prescribed by the position and translation
lattices at \(S_g\cup Q\). The map to the product of these
local quotients is bijective. Injectivity is exactly their
defining congruence. For surjectivity choose a common integral
denominator, clear it, and prescribe the resulting residues
modulo powers of the finitely many different prime ideals.
The Chinese remainder theorem applies to the intermediate
lattice, a finite projective module over \(\mathcal O_{E_x}\).
Equivalently one may prescribe these residues over its
semilocalization, with the same finite quotients after clearing
the denominator. Passing to the chosen local frames identifies
each marked quotient with \(q_w^{-1}O_w^d/O_w^d\).

In the two matrix sums the derivative operators depend on the
matrix cosets, coefficient data and frames. On fixing those
choices, the vector occurs only in the penultimate position
test and the terminal Fourier character. The terminal indicator
at \(T\), the projected small-matrix test and the determinant
frame factors of the terminal matrix are independent of the
vector. The pair uses one common vector in the composition.
Its class mass is taken once, outside the sum over that
vector. These are the explicit composition and two matrix
formulas in
\srcloc{Formalism/LocalToGlobal/Columns/FiniteSumNonzeroLargeDepth}{298}
and lines 1269--1277 there. By the product decomposition of
the quotient we may sum first over all the marked components
while fixing the components at \(Q\). This gives precisely the
product of the marked Fourier coefficients. In particular it
leaves the image flag of each terminal matrix unchanged.
The formula then commutes with summation over paths, branching
on the last coefficient space, and pairing against any of the
projected rows, by linearity of those operations.

Theorem~\ref{thm:markedgauss386} expresses each normalized
factor as zero, a root of unity, or a root of unity times
\(G_s(\nu)/q_w\). The Gauss sum is an algebraic integer
because it is a sum of products of roots of unity. Each of
its nonzero absolute values at a complex embedding lies between
one and \(q_w\). This proves both bounds for the product.
Each marked place above a completely split prime has that
prime as residue size. In fact exactly half the primes above
each \(g_j\) are marked. The ideals of its generator and its
conjugate are coprime, have product \(g_j\mathcal O_{E_x}\),
and conjugation interchanges them. Hence \(\log D_g\) equals
\(\tfrac12[E_x:\QQ]\log(g_1g_2)\).

At either marked rational prime, the prescribed characters on the
residue units are \(\chi_j=\chi^{2(j+1)}\) for a faithful character
\(\chi\) of \(k^\times\). These are the explicit local choices in
\srcloc{Formalism/LocalToGlobal/Columns/AbelianHostFieldConditions}{301}.
The cyclic group \(k^\times\) has order \(q-1>2m\) on the
sufficiently distant intervals, by the prime choices in
Proposition~\ref{prop:ray50}. For \(0\le i<b_0\) the exponent
\(2(i-b_0)\) is nonzero and has absolute value less than
\(q-1\). Thus every varying character \(\nu_i\) in the
arithmetic column is nontrivial. The normalized local factor
has absolute value \(q_w^{-1/2}\) when it varies and is
nonzero. At each other contributing place it is a root of
unity. The formula for \(D_{\rm var}\) follows.

Finally the local table is the unnormalized Fourier transform
of a function of absolute value one on those positions passing
all tests, and zero elsewhere. There are \(q^f\) positions in
the cell. If \(L\ne0\), exactly \(q^{f-1}\) make its determinant
zero. Expanding the squared Fourier norm and summing first
over \(\bar z\) kills all unequal pairs of positions and gives
a factor \(q^d\) on the diagonal. This proves the norm identity.

If the earliest group is \(B_i\), then \(\bar b\) vanishes
in the initial \(r\) columns. Both choices of superdiagonal
character use only those columns, so each character is constant
as \(v\) varies. Thus \(\omega=-\bar z|_{A_{s_i}}\). For
nontrivial \(\nu_i\), Theorem~\ref{thm:markedgauss386}
gives precisely the nonzero multiples of \(L_i\), with
absolute value \(q^{s_i-1/2}\) at every such frequency.
Proposition~\ref{prop:markedcell386} identifies their
common kernel with \(H\cap A_{s_i}\).
For two different blocks, the earlier block must carry a
nonzero restriction for its own table and a zero restriction
for the later table. For equal blocks with different
hyperplanes, a nonzero covector cannot have both kernels.
In either case the product of one table entry with the
conjugate of the other vanishes at each covector separately.
Multiplication by any weights and summation preserves zero.
\end{proof}

The condition on \(H\) couples the reduction of a specified
penultimate matrix to the last row of a specified terminal
matrix. Several paths may satisfy this incidence for the same
covector or contribute at different covectors to the same
output flag. Their signed sums still enter the projected
Gram estimate of Section~\ref{sec:missing}.


\subsection{Inner products and mean squares of the marked Fourier tables}
\label{sec:markedtablegram387}

Fix \(x\), an intermediate class, a marked place \(w\) and its frame
as in Section~\ref{sec:coupledgauss386}. Write \(V=k^d\). For
\(B\subset V\) a linear subspace, let \(B^\circ\subset V^\vee\)
denote its annihilator. We use the Fourier transform
\(\widehat h(z)=\sum_{v\in V}h(v)\theta(-zv)\).
All residue characters and the splitting characters on units take
values in roots of unity. We work in a common number field \(J\)
containing them, and use the ordinary complex inner product after
choosing an embedding. We consider paths whose local test is
nonzero as a function of position.

\begin{lemma}[The singular positions of a marked matrix]
\label{lem:markedsingular387}
For a path with a varying determinant, let \(W=A_{s_i}\) and
\(K=H\cap W\) be the two subspaces of
Corollary~\ref{cor:coupledglobal386}. Choose once and for all,
for this pair, a form \(\lambda:W\to k\) with kernel \(K\).
The positions in its cell where the determinant vanishes are a
single coset \(u+K\) of \(K\) in \(V\). The entire cell is \(u+W\).
There is a root of unity \(\gamma\) such that the path's full
Fourier integral is
\[
  \mathcal F(z)=
   \gamma q^{s_i-1}G_{-1}(\nu_i)\nu_i(t)^{-1}\theta(-zu),
   \qquad z|_W=t\lambda,\quad t\in k^\times ,
\]
and it is zero elsewhere. The class of \(u\) in \(V/K\) is
intrinsic to the tested matrix.

For a path with constant determinants there is a nonempty cell
\(u+W_c\) with \(W_c\subset A_r\), a form
\(\rho_c\in W_c^\vee\), and a root of unity \(\gamma_c\) such that
\[
  \mathcal F_c(z)=
    \gamma_c q^{\dim W_c}\theta(-zu)
               \mathbf1_{z|_{W_c}=\rho_c}.
\]
\end{lemma}
\begin{proof}
In the varying block, the determinant on \(v^0+W\) is
\(d_i+c\lambda(v-v^0)\) for some \(c\in k^\times\).
This assumes every value, so its zeros form a coset \(u+K\).
Taking any \(u\) on that coset writes the determinant as
\(c\lambda(v-u)\). The superdiagonal character and the
other determinants are constant on the cell. Put their values,
the unit splitting factor and \(\nu_i(c)\) into \(\gamma\).
The local function at \(u+y\), \(y\in W\), then equals
\(\gamma\nu_i(\lambda(y))\), with value zero for \(y\in K\).
At \(z|_W=t\lambda\), the fiber above \(a\in k\)
contributes \(q^{s_i-1}\gamma\theta(-zu)\nu_i(a)\theta(-ta)\).
Substituting \(a'=ta\) gives the formula when \(t\ne0\).
At \(t=0\) multiplicative orthogonality gives zero, since
\(\nu_i\ne1\) for these marked characters. At frequencies
outside \(K^\circ\), summing along \(K\) gives zero. Changing
\(u\) within \(u+K\) preserves every value on \(K^\circ\).

For a nonzero constant test the first nonzero column of the
coefficient row \(\bar b\) lies in the initial group, by
Proposition~\ref{prop:markedcell386}. That column fixes all rows
beyond the first \(r\) by the entry congruences. On \(u+W_c\)
the remaining position function equals
\(\gamma_c\theta(\rho_c(v-u))\), since its multiplicative
factors are constant and its superdiagonal exponent is affine.
Summing this character gives the second formula.
\end{proof}

\begin{theorem}[The Gram entries at a marked place]
\label{thm:markedtablegram387}
For two nonzero tables \(F,F'\) of the lemma define
\[
   R(F,F')=\frac{|\sum_{z\in V^\vee}F(z)\overline{F'(z)}|^2}
                    {(\sum_z|F(z)|^2)(\sum_z|F'(z)|^2)} .
\]
The following values hold at every embedding of \(J\).
\begin{enumerate}
\item For two varying tables, \(R=0\) for distinct pairs \((W,K)\).
For the same \((W,K)\), let \(\delta=u-u'\) be the difference of
representatives of the singular cosets. Then
\[
  R=
  \begin{cases}
    1,&\delta\in K,\\
    (q-1)^{-2},&\delta\in W\smallsetminus K,\\
    0,&\delta\notin W.
  \end{cases}
\]
\item For one constant and one varying table put
\(f_c=\dim W_c,\ s=\dim W\), and \(\delta=u-u_c\).
If \(\rho_c=0\) and \(\delta\in W\smallsetminus K\), then
\(R=q^{1+f_c-s}/(q-1)\). In every other case \(R=0\).
\item For two constant tables put \(B=W_c+W'_c\),
\(h=\dim B,\ f_c=\dim W_c,\ f'_c=\dim W'_c\).
If the two forms agree on \(W_c\cap W'_c\) and
\(u_c-u'_c\in B\), then \(R=q^{f_c+f'_c-2h}\).
In every other case \(R=0\).
\end{enumerate}
\end{theorem}
\begin{proof}
Within the support of each table its absolute value is constant.
Thus the constants in front of the characters in the lemma
cancel from the quotient for \(R\). For two varying tables
with different block index or different hyperplane in the same
block, the supports are disjoint by
Corollary~\ref{cor:coupledglobal386}. For the same pair, the
character \(\nu_i\) and the form \(\lambda\) are common to
both paths. Their support is \(K^\circ\smallsetminus W^\circ\)
of size \((q-1)q^{d-s}\). The frequency phases from \(\nu_i\)
cancel. The sum of the remaining additive phases is
\[
       q^{d-s+1}\mathbf1_{\delta\in K}
           -q^{d-s}\mathbf1_{\delta\in W}.
\]
Indeed the sum on either annihilator is its cardinality when
the tested difference belongs to the original space and zero
otherwise, by additive orthogonality. Dividing by the support
size and squaring gives the three values.

For the mixed pair, \(W_c\subset A_r\subset K\) because
the form of the varying block is supported in \(B_i\).
Every frequency on the varying support restricts to zero on
\(W_c\). Consequently the supports are disjoint for
\(\rho_c\ne0\), and the constant table contains the varying
support for \(\rho_c=0\). For a fixed \(t\ne0\) sum first on
the \(q^{d-s}\) extensions of \(t\lambda\) to \(V\).
Their phases sum to zero unless \(\delta\in W\).
In that case the remaining sum, up to a constant of absolute
value one, is
\[
       q^{d-s}\sum_{t\in k^\times}
                \nu_i(t)^{-1}\theta(-t\lambda(\delta)).
\]
It vanishes for \(\delta\in K\) and otherwise has absolute
value \(q^{d-s+1/2}\) by Theorem~\ref{thm:markedgauss386}.
The two support sizes are \((q-1)q^{d-s}\) and \(q^{d-f_c}\).
Dividing the square of the last magnitude by their product
gives the stated value of \(R\).

Two prescribed restrictions on \(W_c,W'_c\) extend to \(V\)
precisely when they agree on their intersection. If they do,
the common frequency set is a translate of \(B^\circ\).
The phases there have sum of absolute value \(q^{d-h}\)
for \(u_c-u'_c\in B\), and sum zero otherwise. Dividing the
square by \(q^{d-f_c}q^{d-f'_c}\) gives the last assertion.
These calculations use orthogonality and the absolute values
of nontrivial Gauss sums, which are preserved at each embedding.
\end{proof}

We apply the singular cosets to groups of signed paths in the
arithmetic column. Fix an intermediate class with its frames,
an input, and a terminal matrix coset. Take one scalar
coefficient of the output, applying the terminal derivative
and the chosen coefficient functional by linearity. Choose a
nonempty subset \(\Sigma\) of the marked places. At these places
group incoming paths by a common varying block and hyperplane
\((W_w,K_w)\) at each \(w\in\Sigma\). At all other places
retain the prescribed values of the tests. In particular
the two neighboring kernels still use the same vector.
The arithmetic meaning of this grouping is that for the fixed
terminal matrix the last covectors \(z_w\) are specified. We
can evaluate the group's coefficient as a finite function of
these variables and then substitute their specified values.

Put \(s_w=\dim W_w\) and write
\[
   A=\prod_{w\in\Sigma} V_w/K_w,\qquad
   L_w=W_w/K_w\subset A,\qquad D_\Sigma=\prod_{w\in\Sigma}q_w.
\]
Here \(L_w\) is supported in the \(w\)-th factor and has order
\(q_w\). Identify \(A^\vee\) with \(\prod_{w\in\Sigma}K_w^\circ\)
by the additive characters \(\theta_w\).
For each path remove the Fourier factors at \(\Sigma\), retain
its sign and all the other factors, and multiply by
\(\prod_{w\in\Sigma}\gamma_w\) from
Lemma~\ref{lem:markedsingular387}. Add the resulting
coefficients for all paths having the same tuple of singular
cosets \(u\in A\), and call the sum \(c(u)\). It includes
the original class masses. For a finite table \(c:A\to J\)
let
\[
  M_wc(u)=\frac1{q_w}\sum_{h\in L_w}c(u+h),\qquad
  P=\prod_{w\in\Sigma}(1-M_w).
\]
The factors commute since the translations commute.

\begin{theorem}[A sharp lower bound for a nonzero integral group of paths]
\label{thm:markedtableintegral387}
Let \(N(z)\) be the group's signed coefficient as a function
of the marked frequencies with the common factor
\(\prod_{w\in\Sigma}q_w^{s_w}\) divided out. It vanishes off
\(A^\vee\). On \(A^\vee\) it is supported where every
\(z_w|_{W_w}=t_w\lambda_w\) has \(t_w\ne0\). There its
value is exactly
\[
    N(z)=
       \left(\prod_{w\in\Sigma}q_w^{-1}
                       G_{-1}(\nu_{i_w})\nu_{i_w}(t_w)^{-1}\right)
       \sum_{u\in A}c(u)\prod_{w\in\Sigma}\theta_w(-z_wu_w).
\]
In particular \(N\) is nonzero somewhere if and only if \(Pc\ne0\).
At any one embedding,
\[
       \frac1{|A|}\sum_{z\in A^\vee}|N(z)|^2
                   =D_\Sigma^{-1}\sum_{u\in A}|Pc(u)|^2.
\]
If \(c(u)\in\mathcal O_J\) for every \(u\) and \(Pc\ne0\), then
\[
 \frac{1}{[J:\QQ]|A|}\sum_{\sigma:J\hookrightarrow\CC}
                 \sum_{z\in A^\vee}|\sigma N(z)|^2
   \ \ge\ \prod_{w\in\Sigma}\frac{q_w-1}{q_w^2}.
\]
This constant is attained among integral tables \(c\).
For each prescribed \(z\) with \(N(z)\ne0\), the embedding
average of \(|\sigma N(z)|^2\) is at least \(D_\Sigma^{-1}\)
under the same integrality assumption.
\end{theorem}
\begin{proof}
The signed coefficient of the two coupled edges is the sum
over incoming paths with the product of the local factors
from Corollary~\ref{cor:coupledglobal386}. The terminal matrix
coset, its derivative, and the coefficient functional are
fixed during this summation. The position at the terminal prime,
the top matrix test and the mass of the intermediate class
are independent of the varying marked covectors in this
formula. All dependence under consideration occurs in the
partial Fourier factors at \(\Sigma\). At each such place
the lemma expresses it as a common factor times
\(\gamma_w\theta_w(-z_wu_w)\). The latter phase depends only
on the singular coset because \(z_w\in K_w^\circ\) on the
support. Grouping and restoring the prescribed covectors
therefore recovers the actual contribution, including the
signs within each \(c(u)\).

Averaging translates along \(L_w\) multiplies a Fourier
coefficient by one if its character is trivial there and
by zero otherwise. Thus the Fourier transform of \(Pc\)
agrees with that of \(c\) exactly on the retained frequency
set and vanishes on its complement. The squared absolute
value of the product of Gauss factors in the displayed
formula is \(D_\Sigma^{-1}\). Expanding the squared Fourier
transform and using character orthogonality on \(A\) gives
the norm identity and the equivalence of nonvanishing.

For the integral estimate put the norm
\(\|a\|^2=[J:\QQ]^{-1}\sum_\sigma|\sigma a|^2\) on
\(\mathcal O_J\), extending it to arrays by summing over
coordinates. Every nonzero integral scalar has squared norm
at least one by the arithmetic-geometric mean inequality
and its nonzero integral absolute norm. The operator \(P\)
acts separately on each coset of \(\prod_w L_w\) in \(A\).
Choose one on which the result is nonzero. On this product
of lines perform the mean subtractions one coordinate at
a time. After centering in any first set of coordinates,
every nonzero resulting slice has squared norm at least
the product of \(1-q_w^{-1}\) over those coordinates.
For the empty set this is the scalar assertion just proved.

For the induction step let \(v_1,\ldots,v_q\) be the \(q\)
slices in a new coordinate after the preceding mean
subtractions, and put \(\bar v=q^{-1}\sum_jv_j\). In the
Hilbert space of arrays over all embeddings,
\[
       \sum_{j=1}^{q}\|v_j-\bar v\|^2
             =\frac1q\sum_{1\le j<l\le q}\|v_j-v_l\|^2 .
\]
This follows by expanding both sides. For a nonzero centered
array at least two entries differ, giving at least \(q-1\)
unequal unordered pairs. Indeed a value occurring \(h\) times,
\(1\le h<q\), differs from the entries in \(h(q-h)\ge q-1\)
pairs. Each difference is the preceding projection of a
difference of integral input arrays, so every nonzero
difference satisfies the inductive lower bound. Multiplying
that bound by \((q-1)/q\) proves the step. Summing over all
product cosets and using the norm identity proves the inequality.
For \(c\) equal to a single point mass of value one the
centered vector on each line has squared norm \(1-q_w^{-1}\).
Their tensor product gives equality. Finally the Fourier sum
of \(c\) at any prescribed \(z\) is an algebraic integer.
If it is nonzero, its embedding average of squared absolute
values is at least one by the same norm argument. The common
Gauss factor has squared absolute value \(D_\Sigma^{-1}\)
at every embedding, giving the last assertion.
\end{proof}

We next sum over terminal matrix cosets themselves. Let \(N\)
be the intermediate lattice and \(N'=M\perp\mathcal O_{E_x}e\)
the lifted terminal lattice in its large space. Write
\(\mathcal H\) for the rational homomorphisms taking \(N\)
into \(N'\) locally at all the selected places, including
those over \(T\). The terminal matrix index is the finite
module \(c_A^{-1}\mathcal H/c_B\mathcal H\) of
Section~\ref{sec:covectorpaths352}. The two positive
integers \(c_A,c_B\) are the position denominator and
congruence modulus of this sum. Each marked rational prime
divides \(c_B\), by the positive depth used at the marked
places. These parameters and the matrix sum are specified in
\srcloc{Formalism/LocalToGlobal/Columns/FiniteSumNonzeroLargeDepth}{636}
and lines 1107--1113 there.

\begin{proposition}[Which marked frequencies the terminal matrices sample]
\label{prop:markedsampling387}
Fix the intermediate and terminal frames, all components of the
terminal matrix coset outside \(\Sigma\), and its projection
to the small lattice \(M\) at the places in \(\Sigma\).
Suppose the tests determined by these fixed choices pass.
Put \(b_w=v_w(c_B)\ge1\). After the terminal integrality
test on its normal row, every tuple of residue covectors
\(z\in\prod_{w\in\Sigma}V_w^\vee\) occurs, with exactly
\[
                 m_\Sigma=\prod_{w\in\Sigma}q_w^{d(b_w-1)}
\]
terminal matrix cosets for each tuple. In particular the
fraction in the frequency support of a fixed group of
varying paths is exactly
\(\prod_{w\in\Sigma}(q_w-1)q_w^{-s_w}\).

Let \(C(u)\) be the coefficient table of the preceding theorem
with values in the intermediate coefficient space, before applying
the terminal derivative and output functional. Let \(T_\xi\)
be the terminal coefficient operator followed by this functional
for the matrix coset \(\xi\). Include in \(T_\xi\) the fixed
terminal scalar tests away from the varying Fourier factors,
counting each original scalar once. Put
\(J(z)=\sum_{\xi\mapsto z}T_\xi\). With
\(\mathcal Z=\prod_w(K_w^\circ\smallsetminus W_w^\circ)\)
and \(\widehat C\) defined entrywise, the contribution of all
these terminal cosets and incoming paths to the projected
coefficient equals
\[
    \left(\prod_{w\in\Sigma}q_w^{s_w}\right)
        \sum_{z\in\mathcal Z} g(z)\,J(z)\bigl(\widehat C(z)\bigr),
    \qquad
    g(z)=\prod_{w\in\Sigma}q_w^{-1}G_{-1}(\nu_{i_w})
                                      \nu_{i_w}(t_w)^{-1}.
\]
\end{proposition}
\begin{proof}
After multiplication by \(c_A\) the matrix quotient is
\(\mathcal H/(c_Ac_B)\mathcal H\). Over the ring of elements
integral at the selected places, \(\mathcal H\) is the
homomorphism module of the two localized lattices. Their
localizations are finite free at each selected place.
The Chinese remainder theorem modulo \(c_Ac_B\) identifies
this quotient with the product of the corresponding local
matrix quotients. To see surjectivity for projective lattices,
present their homomorphism module as a direct summand of a
finite free module, first lift arbitrary prescribed residues
in the free module, then apply the projector. The same
decomposition identifies the kernel with multiplication by
\(c_Ac_B\).

At \(w\in\Sigma\), division by \(c_A\) restores a matrix
between the local lattices, with entries allowed in
\(c_A^{-1}O_w\) modulo \(c_BO_w\). Since \(N'=M\perp O_we\)
locally, the small projection and last row give independent
coordinates on this quotient in the chosen bases. The
covector used in the Fourier factor is the last row divided
by \(\pi_T\) in the intermediate frame. The generator
\(\pi_T\) is a unit at \(w\), so this change preserves the
integral row lattice and its congruence sublattice. The
allowed integral rows are therefore \(O_w^d/c_BO_w^d\).
Reduction onto \(V_w^\vee\) is surjective and a fiber has
\(q_w^{d(b_w-1)}\) elements. The top position test uses the
small projection only, and the position equation at \(T\)
uses components outside \(\Sigma\). Thus fixing these data
imposes no further equation on the last rows. Multiplication
of the local counts gives \(m_\Sigma\).
The support at \(w\) consists of \(q_w-1\) nonzero multiples
of a form on \(W_w\), each with \(q_w^{d-s_w}\) extensions.
Division by \(q_w^d\) and multiplication over \(w\) gives
the fraction asserted.

Apply Theorem~\ref{thm:markedtableintegral387} to each
coordinate before applying \(T_\xi\). On the fiber above
\(z\), both \(g(z)\) and \(\widehat C(z)\) are fixed.
Summing the terminal operators gives \(J(z)\). Restoring the
common powers \(q_w^{s_w}\) and summing over \(z\) proves
the formula.
\end{proof}

After clearing the incoming denominators, the remaining
estimate in Section~\ref{sec:missing} concerns the operators
\(J(z)\) and their signed sum in the same projected output.


\subsection{The theta sections in the last row of the terminal matrix}
\label{sec:normalrowtheta388}

Fix a word, an intermediate class and a terminal class in
Proposition~\ref{prop:markedsampling387}. Write \(d=N(m-1)\) for the
intermediate rank and \(n+1=N(m)\) for the large terminal rank.
In the terminal frame the target lattice is
\(N'=M\perp\mathcal O_E e\), where \(E=E_x\) and \(M\) has rank
\(n\). Thus the last row of a homomorphism from \(N\) to \(N'\)
belongs to \(N^\vee=\operatorname{Hom}_{\mathcal O_E}(N,\mathcal O_E)\)
when the homomorphism is integral. Put \(f=[E:K]\) and \(D=fd\).
Use the normal Hermitian coordinates of Theorem~\ref{thm:lattice55}.

Choose one or both of the two marked rational primes. Take for
\(\Sigma\) all the selected places over their marked halves.
Write \(\beta=\prod_i\bar\pi_{g_i}\), \(G=\prod_i g_i\), where
both products run over the chosen primes, and put \(Q=G^D\).
The prime factorizations and the orientation of the halves are
given in \srcloc{Formalism/LocalToGlobal/Columns/AbelianHostFieldConditions}{190}.
In particular, \(\beta\bar\beta=G\), each chosen prime splits
completely, and
\(Q=\prod_{w\in\Sigma}q_w^d\). At these places write \(U_w\)
for the lattice of integral normal rows, in the coordinates in
which the Fourier covector was computed. Division by \(\pi_T\)
at a marked place is a unit change of coordinates.
Fix the small part of the terminal matrix everywhere, and the
normal cells at the other selected places. One may also fix a
residue cell at the other marked prime by summing over its higher
digits. All fixed finite tests are assumed to pass. The position
test at \(T\) is applied in its original coordinates.

\begin{proposition}[The sections indexed by the marked normal rows]
\label{prop:normalcells388}
In the complex theta realization, summing the terminal theta
sections over one integral residue tuple
\(z\in\prod_{w\in\Sigma} U_w/q_w U_w\), before differentiation,
gives a product
\(\Theta_{\rm top}\Theta_z\). Here \(\Theta_{\rm top}\) is fixed.
The \(Q\) normal sections \(\Theta_z\) are an orthogonal basis of
sections of one positive line bundle on a torus of dimension \(D\).
Its period lattice and covolume in the normal metric are
\[
 \Gamma=\delta\Phi(\beta\,\mathfrak b^{-1}\mathfrak a N^\vee),
 \qquad \operatorname{vol}(\Gamma)=Q .
\]
Here \(\mathfrak b\) and \(\delta\) are the period data of
Theorem~\ref{thm:lattice55}. The fractional ideal \(\mathfrak a\)
has exponent \(h_v\) at a frozen split half with normal cell of
depth \(h_v\), exponent \(-h_v\) at its conjugate and exponent zero
elsewhere. Put \(s_0=\|\sum_z\Theta_z\|_{\rm av}\), with the
Hermitian metric and probability Haar measure on the torus. Then
\[
     s_0>0,\qquad \|\Theta_z\|_{\rm av}^2=s_0^2/Q
                       \quad\text{for every }z .
\]
\end{proposition}
\begin{proof}
Consider first one split place, with uniformizer \(q\), position
lattice \(U\), dual frequency lattice \(U^\circ\), and an additive
character \(\psi\) of conductor the local integers. For \(a\in U\)
and \(h\ge1\), finite character orthogonality gives the identity
of functions of \(y\)
\[
 {\bf1}_{a+q^hU}(y)
  ={\bf1}_{U}(y)\,q^{-hd}
     \sum_{\eta\in U^\circ/q^hU^\circ}
       \psi\bigl(\langle\eta,y-a\rangle/q^h\bigr).
\]
Indeed, for \(y\in U\) the sum of a nontrivial character on the
quotient is zero, and the trivial character sums to \(q^{hd}\).
The finite theta intertwiner sends each summand to the same
modulation of the vacuum. In the class formula its finite Fourier
coefficient is \(\sqrt{\mathrm{vol}}/|S_A|\), and the class
distribution multiplies it by \(1/\sqrt{\mathrm{vol}}\). Their
product is the average over the frequency representatives.
For the position translate \(a\), conjugating a modulation by
translation supplies the displayed phase. Equivalently one can
translate the characteristic function at zero by \(a\) in the
split position model. These are the definitions in
\srcloc{Formalism/LocalToGlobal/KernelLattices/XiIndCClass}{20}
and \srcloc{Formalism/LocalToGlobal/KernelLattices/ThDistC}{87}.
Thus the fine cell enters with its characteristic function.
Its higher digits partition a residue cell disjointly. Their sum
is exactly \({\bf1}_{\tilde z+qU}\), with \(\tilde z\) any lift of \(z\).
Taking products gives this assertion at every chosen half.

The period module decomposes into the homomorphisms to \(M\) and
the homomorphisms to \(\mathcal O_Ee\). The Hermitian form is the
orthogonal sum of the two trace forms. Hence its quadratic
semicharacter is the product of the two semicharacters. The fixed
finite matrix coset is also a product in these two sets of
coordinates. The theta realization respects this decomposition.
For clarity, the vacuum line of each self-dual summand is
one-dimensional by Riemann--Roch for a principal polarization.
Their tensor product gives the vacuum of the orthogonal sum.
Each split translation or modulation in the finite expansion acts
on the corresponding factor. The displayed characteristic-function
identity and linearity now give \(\Theta_{\rm top}\Theta_z\).
The normalization of the initial vacuum can be included once in
\(\Theta_{\rm top}\). A change of this convention multiplies all
the normal sections by the same nonzero scalar.

We describe the periods on the normal factor more explicitly.
The function of position on a frozen half is
\({\bf1}_{a+q^hU}\). Translations by \(q^hU\) fix it. Modulations
by \(q^{-h}U^\circ\) multiply it by a unit scalar depending on
\(a\). Relative to the unrestricted integral vacuum this replaces
the two local lattices by \(q^hU\) and \(q^{-h}U^\circ\).
Their product is the self-dual change specified by \(\mathfrak a\).
At a varying half, translations by \(qU\) and modulations by
\(U^\circ\) have the same action on every residue cell.
Accordingly the common period lattice is obtained by multiplying
the position component by \(q\) at each varying half, leaving its
conjugate component as it was. This is multiplication of the
global normal module by \(\beta\). These local descriptions
determine the global lattice, since two fractional lattices that
agree at every finite prime agree over \(\mathcal O_E\).
All the \(\Theta_z\) therefore have the same automorphy factors on
\(\Gamma\).

The module \(\mathfrak a N^\vee\) is self-dual for its Hermitian
form because \(\mathfrak a\overline{\mathfrak a}=1\).
The trace polarization gives covolume one before multiplication
by \(\beta\), by Theorem~\ref{thm:lattice55}. Multiplication by
\(\beta\) is a dilation by \(\sqrt G\) followed by a unitary map
in \(D\) complex variables. Its real determinant is \(G^D\).
The resulting positive line bundle has \(Q\) holomorphic sections
counted by dimension. Indeed, the top exterior power of its
polarization divided by \(D!\) integrates to \(Q\). Riemann--Roch
on the complex torus gives this Euler characteristic, and the
higher cohomology vanishes by positivity. Changing the fixed
unitary character affects neither the curvature nor this computation.

Finally \(q^{-1}U^\circ/U^\circ\) acts by normalized unitary
modulations on these sections. Its characters distinguish the
residue cells by the perfect residue pairing. Such operators
preserve the indicated line on the common torus, since their
commutators with the common periods are trivial. Distinct cells
are therefore orthogonal. The position translations by
\(U/qU\) permute the cells up to scalars of absolute value one,
so all norms coincide. The positions and modulations at distinct
places can be prescribed independently. One may realize them
by rational lifts in the corresponding fractional lattice modulo
the common lattice, by the Chinese remainder theorem.
The sum of all the residue cells is the whole integral normal
vacuum at the varying halves, with the frozen test elsewhere.
It is nonzero. Locally a frozen cell is a nonzero multiple of a
unitary translate of the normalized vacuum for its changed
self-dual lattice. The finite isometric intertwiner takes its
tensor product to the vacuum line for that lattice, as in the
split-lattice calculation of Theorem~\ref{thm:one61}.
Orthogonality and equal norms now give the norm formula, and the
dimension calculation proves that the family is a basis.
\end{proof}

\begin{theorem}[Normal derivatives in the projected terminal operator]
\label{thm:normaljet388}
Let \(J(z)=\sum_{\xi\mapsto z}T_\xi\) be the terminal operator in
Proposition~\ref{prop:markedsampling387}, in its complex theta
realization. Use an orthonormal normal frame respecting the
embeddings \(\tau\), and put
\[
 \ell_\alpha(F)=
   \frac{\partial^\alpha F(0)}
        {\pi^{|\alpha|/2}\sqrt{\alpha!}},\qquad
 B_\tau=\sum_{j=1}^n(\lambda_{\tau j}-\mu_{\tau j}).
\]
There are fixed coefficient maps \(A_\alpha\), determined by the
small coordinates and the prescribed branch projection, such that
\[
   J(z)=\sum_{\alpha\in\mathcal I}\ell_\alpha(\Theta_z)A_\alpha,
   \qquad
   \mathcal I=
     \left\{\alpha\in\mathbb N^{f\times d}:
             \sum_{i=1}^d\alpha_{\tau i}=B_\tau
                         \text{ at every }\tau\right\}.
\]
In particular the linear span of these operators has dimension at
most \(\prod_\tau\binom{B_\tau+d-1}{d-1}\).
\end{theorem}
\begin{proof}
Here is a construction of \(A_\alpha\) which also fixes the
relative scalars. Multiply the input flag polynomial by the
prescribed determinant powers \(\prod_\tau\det(X_\tau)^{\nu_\tau}\).
Replace its variable \(X_{\tau i j}\) by
\(\sum_{\gamma=1}^{n+1}Y_{\tau\gamma j}D_{\tau\gamma i}\).
In each term of the expansion a factor \(D_{\tau\gamma i}\)
denotes one differentiation in that matrix coordinate. Include
the assigned differential scale \(c_{D,\tau}\) for each such
factor. This is the ordinary mixed-derivative definition in
\srcloc{Formalism/LocalToGlobal/KernelLattices/ThDistC}{21}
and lines 62--104 there. Apply the derivatives with
\(\gamma\le n\) to \(\Theta_{\rm top}\) at zero. Express the
remaining derivatives in the orthonormal normal coordinates.
For the derivative of index \(\alpha\) put the factor
\(\pi^{|\alpha|/2}\sqrt{\alpha!}\) into its coefficient.
An orthonormal change within the \(d\) coordinates at one
embedding preserves their total degree. In this manner the only
factor involving the varying section is \(\ell_\alpha(\Theta_z)\).

To see which indices survive projection, first observe that
the substitution already lies in the prescribed target flag
space, coefficient by coefficient in the \(D\)'s. At one
embedding, that space is spanned by products of minors of the
first \(j\) columns, with column counts prescribed by the
partition of the weight after subtraction of its last entry.
Each source minor of \(D^{t}Y_{:,\le d}\) expands by
Cauchy--Binet into target minors of the same first \(j\)
columns, with minors of \(D\) as coefficients. Each determinant
factor adds a minor of size \(d\). The target partition for this
theta step is the source partition plus \(\nu_\tau\) in its
first \(d\) entries, extended by zeros. This is the shift of the
weight lists used in \eqref{eq:letters111}, expressed in partition
coordinates in
\srcloc{Formalism/LocalToGlobal/KernelLattices/FirstKernelLawReading}{233}.
Thus the target retraction fixes the resulting polynomial.

By Theorem~\ref{thm:terminalformula347} the terminal pairing uses
precisely the monomials having last-row exponents
\(\lambda_{\tau j}-\mu_{\tau j}\) for \(1\le j\le n\) and
zero in the last column. These exponents are nonnegative by
interlacing. Each occurrence of \(Y_{\tau,n+1,j}\) came from
exactly one \(D_{\tau,n+1,i}\), and every normal derivative arose
this way. Their number at \(\tau\) is \(B_\tau\).
Apply the indicated coefficient extraction, Fischer factors and
small output functional to each coefficient still present.
Include the class factor \(\chi_F^{-1}\), the frame determinants,
the fixed finite scalars and the normal-coordinate factor from
the branch map in the same places as in the prescribed operator.
This defines \(A_\alpha\) independently of \(z\) and proves the
identity. Testing every small coefficient proves the corresponding
map identity before evaluation by a single functional. Counting
the multi-indices separately at each embedding gives the bound.

The computation takes place in a specified complex realization.
For the normalized CM jets, the change of embedding uses the
transported differentials of Theorem~\ref{thm:transport59} and
the theta functional of Theorem~\ref{thm:vac242}. Dividing by
the Fock factors in the definition of \(\ell_\alpha\) and putting
them into \(A_\alpha\) is an identity in that realization. Descent
applies to the original normalized coefficient whenever it
belongs to the designated coefficient field.
\end{proof}

\begin{corollary}[Fourier synthesis of the signed normal rows]
\label{cor:normaltranslate388}
Let \(h(v)\) be a scalar incoming position test at the varying
marked places and write \(\theta(zv)=\prod_w\theta_w(z_wv_w)\).
Use \(\widehat h(z)=\sum_v h(v)\theta(-zv)\). Define
\(\mathcal J_\eta\) by applying the same terminal operation to
the full integral normal row, modulated by the residue character
\(z\mapsto\theta(z\eta)\). Then
\[
      \sum_z \widehat h(z)J(z)=\sum_v h(v)\mathcal J_{-v}.
\]
For an individual marked path the right side is an integral
linear combination over the ring generated by the values of
the residue and splitting characters.
\end{corollary}
\begin{proof}
A full integral row modulated by \(\theta(z\eta)\) is the sum of
its residue cells with those coefficients. In the normal theta
line this modulation is the Weyl operator of the dual residue
class in \(q^{-1}U^\circ/U^\circ\) at each varying half. Thus
\(\mathcal J_\eta=\sum_z\theta(z\eta)J(z)\), including the
original scalar of the full row when \(\eta=0\).
Expanding \(\widehat h\) and interchanging the two finite sums
proves the identity with the indicated minus sign. For a marked
path the original position test takes the values zero and roots
of unity, by Lemma~\ref{lem:markedsingular387}. Hence the
coefficients \(h(v)\) are algebraic integers. The identity
requires no division by the residue cardinalities. Integrality
of the modulated operators themselves and the size of the
resulting signed period require their own estimates.
\end{proof}


\subsection{A uniform mean square for the projected terminal operators}
\label{sec:normalbergman389}

The operators in Theorem~\ref{thm:normaljet388} use derivatives
at one point of many sections of a single normal theta line
bundle. We estimate their Gram matrix before contracting it with
the incoming path coefficients. The dependence on the dimension
and on the maximum order of differentiation matters because both
grow along the intervals in Section~\ref{sec:missing}.

\begin{lemma}[Derivatives of a theta basis on a separated lattice]
\label{lem:thetaframe389}
Let \(\Gamma\subset\mathbb C^D\) be a full lattice of real
covolume \(Q\) for the usual Euclidean metric. Use an integral
positive polarization with Hermitian form the usual inner
product, complex-linear in its first variable. A holomorphic
section in the Hermitian frame satisfies
\[
 F(u+\lambda)=a_\lambda
    e^{\pi u\cdot\bar\lambda+\pi|\lambda|^2/2}F(u)
       \quad(\lambda\in\Gamma),\qquad |a_\lambda|=1 .
\]
Give the space of sections its inner product with weight
\(e^{-\pi|u|^2}\) and probability Haar measure on the torus.
Let \((s_j)_j\) be an orthonormal basis and, for \(|\alpha|\le L\),
put \(B^\circ_{\alpha j}=\ell_\alpha(s_j)\), with the derivative
normalization of Theorem~\ref{thm:normaljet388}. Suppose
\(|\lambda|\ge R>0\) for every nonzero \(\lambda\in\Gamma\) and
\[
       \pi R^2\ge \max\{64L,\ 16D\log 5,\ 8\log 2\}.
\]
Then, in operator norm on the space of derivative indices,
\[
     \left\|Q^{-1}B^\circ(B^\circ)^*-I\right\|
                  \le 2e^{-\pi R^2/8}.
\]
The same bound holds after restricting the matrix to any
common subset of its row and column indices.
\end{lemma}
\begin{proof}
Write \(K(u,w)=\sum_j s_j(u)\overline{s_j(w)}\). With the
convention on volume used here its period expansion is
\[
 \frac{K(u,w)}{Q}
  =e^{\pi u\cdot\bar w}
    \sum_{\lambda\in\Gamma}a_\lambda
      e^{-\pi|\lambda|^2/2+\pi u\cdot\bar\lambda
                                  -\pi\lambda\cdot\bar w}.
\]
We give the normalization argument. For Gaussian measure
\(e^{-\pi|w|^2}\,dw\) on the full space, the squared norm of
\(w^\alpha\) is \(\alpha!/\pi^{|\alpha|}\), by polar integration
in each coordinate. Distinct monomials are orthogonal by angular
integration. Their normalized sum in the reproducing kernel is
\(e^{\pi u\cdot\bar w}\). Completeness follows by expanding an
entire function in its Taylor series on polydiscs, integrating
the angular variables first and then increasing the radii. These
are the Gaussian conventions of
\href{https://doi.org/10.1002/cpa.3160140303}{V.~Bargmann, Part~I}.

For a torus section \(F\), the function
\(|F(w)|e^{-\pi|w|^2/2}\) is bounded by compactness. Apply the
full-space reproducing identity first to \(F(tw)\) for \(0<t<1\).
Let \(t\) increase to one. The integrand at fixed \(u\) is bounded
by a constant times
\(\exp(\pi|u|\,|w|-\pi|w|^2/2)\), which is integrable. Split the
resulting integral into translates of a fundamental domain. On
its \(\lambda\)-translate, substitute the automorphy formula for
\(F(w+\lambda)\) and expand \(|w+\lambda|^2\). This yields the
displayed series as the kernel for ordinary Euclidean integration
over that domain. The series and all of its derivatives converge
normally on compact sets by Gaussian decay. For fixed \(u\) its
complex conjugate satisfies the section law in \(w\), as follows
by reindexing the translated domains, or by substituting
\(w+\mu\) in the series and using the cocycle law for \(a_\lambda\).
It thus represents evaluation in the space of sections. Changing
the integration measure to the probability measure multiplies
the reproducing kernel by \(Q\), as asserted.

In the variables \(v=\sqrt{\pi}\,u\) the normalized Fock
monomials are \(v^\alpha/\sqrt{\alpha!}\), and \(\ell_\alpha\)
reads their coefficient. Take the corresponding derivatives
in \(u,\bar w\) of the kernel at the origin. The term
\(\lambda=0\) gives the identity matrix. A term with
\(\lambda\ne0\), apart from \(a_\lambda\), gives the matrix on
degrees at most \(L\) of the translation on the first variable
\[
 e^{-\pi|\lambda|^2/2}
 \exp\bigl(\sqrt{\pi}\,\bar\lambda\cdot v\bigr)
 \exp\bigl(-\sqrt{\pi}\,\lambda\cdot\partial_v\bigr).
\]
Indeed its action on the full-space kernel yields exactly the
three \(\lambda\)-dependent exponential factors in the series.
If \(P_L\) is orthogonal projection onto degree at most \(L\),
then the norm of \(\lambda\cdot\partial_v\) there is at most
\(|\lambda|\sqrt L\). A unitary rotation reduces this to
differentiation in the first coordinate, which lowers a monomial
with multiplier \(\sqrt{\alpha_1}\). The compression of
multiplication by \(\bar\lambda\cdot v\) is its adjoint and
satisfies the same bound. The derivative exponential preserves the
degree cutoff. Multiplication raises degree at every positive
power, so its exponential, projected back to the cutoff, equals
the exponential of the compressed multiplication operator on
that space. The norm of the compressed translate is consequently
bounded by
\[
        \exp\{-\pi|\lambda|^2/2+2\sqrt{\pi L}\,|\lambda|\}
                    \le e^{-\pi|\lambda|^2/4}.
\]
The last inequality uses \(\pi R^2\ge64L\).

In the shell \(jR\le|\lambda|<(j+1)R\), \(j\ge1\), put an
open ball of radius \(R/2\) at each lattice point. The balls
are disjoint and lie in a ball of radius \((j+3/2)R\).
Comparison of real \(2D\)-dimensional volumes bounds the shell
count by \((2j+3)^{2D}\). Since
\(\log(2j+3)\le j^2\log5\) and \(\pi R^2\ge16D\log5\),
the sum of the norms on this shell is at most
\(e^{-\pi j^2R^2/8}\). The sum over \(j\ge1\) is bounded by
\(2e^{-\pi R^2/8}\), using \(\pi R^2\ge8\log2\).
The triangle inequality proves the estimate. Compression to
the specified indices can only decrease the error norm.
\end{proof}

\begin{theorem}[The normal derivative matrix on the original interval]
\label{thm:normalframe389}
In Proposition~\ref{prop:normalcells388} take
\(R^2=2fG/\operatorname{rd}(E_x)\) and
\(\eta_x=2\exp(-\pi R^2/8)\). For every sufficiently large
\(x\) on the intervals of Section~\ref{sec:missing}, the
matrix \(B_{\alpha z}=\ell_\alpha(\Theta_z)\),
\(\alpha\in\mathcal I\), satisfies
\[
               \left\|s_0^{-2}BB^*-I_{\mathcal I}\right\|
                       \le\eta_x\le e^{-x^{\varepsilon/8}} .
\]
This holds for either choice of one marked prime and for the pair,
uniformly in the word, the frozen normal cells, the class frames
and the transported complex realization with its own metric.
\end{theorem}
\begin{proof}
The squared length of a nonzero element of
\(\delta\Phi(\mathfrak b^{-1}\mathfrak aN^\vee)\) is at least
\(2f/\operatorname{rd}(E_x)\) by the trace and norm calculation
of Theorem~\ref{thm:lattice55} for the self-dual source dual with
the norm-one ideal twist. Multiplication by \(\beta\) scales
every squared length by \(G\). The frozen normal translate
changes the character on the period lattice and leaves this
length calculation intact. For a conjugate CM realization the
trace metric transports as in Theorem~\ref{thm:transport59}.
The conjugate ideal of the elliptic differential is accompanied
by its inverse norm in the Hermitian form. In the trace proof,
the lower bound for the geometric mean of the squared lengths
gains exactly that ideal norm before multiplication by the
metric factor. Thus the same lower bound holds, and every
conjugate of \(\beta\) still has complex absolute value \(\sqrt G\).

Put \(L=\sum_\tau B_\tau\). The ranks and \(f\) are bounded by
fixed powers of \(x\) on these intervals. The weight lists in
Theorem~\ref{thm:terminalformula347} bound each difference in
the interlacing partitions by \(O(xb+x^C)\), where
\(b=\exp(x^\varepsilon)\) and \(C\) is fixed. Hence
\(L\ll x^{C'}\exp(x^\varepsilon)\), and \(D=fd\ll x^{C'}\)
after increasing \(C'\). The abelian field \(E_x\) has conductor
dividing the common modulus \(N\) of Theorem~\ref{thm:conductor232},
for which \(\log(2N)\ll_{\psi,\varepsilon}x^{\varepsilon/4}\).
The conductor--discriminant formula expresses
\(|\operatorname{disc}(E_x)|\) as a product of the conductors
of its characters. Each divides \(N\), so
\(\operatorname{rd}(E_x)\le N\). Each marked rational prime
is at least \((kb)^8\) by Proposition~\ref{prop:ray50}.
Already for one chosen prime it follows that
\[
 R^2\ \ge\ 2f\exp\{8x^\varepsilon-O_{\psi,\varepsilon}
                                    (x^{\varepsilon/4})\}.
\]
This dominates the degree and dimension cutoffs in
Lemma~\ref{lem:thetaframe389} uniformly, and also gives the
stated upper bound on \(\eta_x\).

By Proposition~\ref{prop:normalcells388} an orthonormal basis on
this common torus is \(s_z=\sqrt Q\,\Theta_z/s_0\).
Consequently \(Q^{-1}B^\circ(B^\circ)^*=s_0^{-2}BB^*\)
after restriction to \(\mathcal I\). Applying the lemma proves
the assertion. The other marked prime may stay frozen throughout
this proof when only one prime is varied.
\end{proof}

\begin{corollary}[The projected Gram matrix and the signed input]
\label{cor:weightedkernel389}
Equip the fixed complex input and output coefficient spaces with
their chosen Hermitian norms. For the \(A_\alpha\) of
Theorem~\ref{thm:normaljet388} set \(S=\sum_\alpha A_\alpha A_\alpha^*\).
Then, in the order of Hermitian quadratic forms,
\[
 (1-\eta_x)s_0^2S\ \preceq\ \sum_z J(z)J(z)^*
                           \ \preceq\ (1+\eta_x)s_0^2S .
\]
For an arbitrary prescribed vector of incoming coefficients
\(Y=(Y(z))_z\), let \(\mathcal JY=\sum_zJ(z)Y(z)\).
If \(S\ne0\), write \(\mu\) for its smallest positive eigenvalue.
Then
\begin{align*}
  (1-\eta_x)s_0^2\mu\,\operatorname{dist}(Y,\ker\mathcal J)^2
    &\le \|\mathcal JY\|^2,\\
  \|\mathcal JY\|^2
    &\le (1+\eta_x)s_0^2\|S\|\sum_z\|Y(z)\|^2 .
\end{align*}
If \(S=0\), every \(J(z)\) is zero.
\end{corollary}
\begin{proof}
Define a linear map \(A\) from the direct sum of one copy of the
input space for each \(\alpha\) to the output by
\(A((v_\alpha)_\alpha)=\sum_\alpha A_\alpha v_\alpha\).
Then \(\mathcal J=A(B\otimes I)\) and \(AA^*=S\).
The previous theorem bounds \(BB^*\) above and below by the
corresponding scalar multiples of the identity. Tensoring with
the identity on the input and composing with \(A,A^*\) proves
the first display. Since \(\eta_x<1\), it also proves
\(\ker(\mathcal J\mathcal J^*)=\ker S\).
On the orthogonal complement of this kernel, the smallest
eigenvalue of \(\mathcal J\mathcal J^*\) is at least
\((1-\eta_x)s_0^2\mu\). The positive eigenvalues of
\(\mathcal J^*\mathcal J\) are the same, by the singular value
decomposition. Applying this bound to the component of \(Y\)
perpendicular to \(\ker\mathcal J\) gives the lower inequality.
The largest eigenvalue gives the upper inequality. If \(S=0\)
then \(\sum_\alpha\|A_\alpha^*u\|^2=0\) for every output vector
\(u\), and all the maps vanish.

For the actual block in Proposition~\ref{prop:markedsampling387},
take \(Y(z)=g(z)\widehat C(z)\) on its frequency support and
zero elsewhere, restoring its displayed common scalar. This
identifies the distance that occurs in a lower bound for that
block. A lower estimate for the norm of \(\widehat C\) on its
full table gives no lower estimate for this distance on its own.
In particular the signed block can vanish even when both
\(S\) and \(Y\) are nonzero. Section~\ref{sec:missing} requires a
comparison for the particular sum over the prescribed classes
and paths, as well as the local cohomology estimate there.
\end{proof}

\section{Terminal covectors and initial class tables}\label{nav:chapter-terminal-covectors}

\subsection{Summing at both marked primes before varying the terminal covector}
\label{sec:refinednormal390}

Keep the terminal prime \(T=\pi_T\bar\pi_T\) and the two marked
primes \(g_1,g_2\) chosen for the same \(x\). At \(g_1,g_2\)
write \(Z=\prod_w U_w/q_w U_w\) for the normal residue rows,
\(G=g_1g_2\), and \(Q_g=|Z|=G^D\). Let
\(\mathcal T=\{w:w\mid\pi_T\}\); it has \(f\) elements. Write
\(t_w\in U_w/TU_w\) for the raw integral normal row at \(w\),
using the intermediate lattice as source. In the decorated
lattice sum the intermediate covector at \(w\) is
\((\bar\pi_T\bmod w)t_w\), by Theorem~\ref{thm:covectorpaths352}.
This factor is a unit. Thus all nonzero \(t_w\) give exactly the
allowed intermediate covectors at this edge.

Fix the underlying intermediate lattice, the frames away from
\(T\), the terminal small class and its test, the small matrix
cells, and any other finite components outside these three
primes. Denote the resulting normal sections by
\(\Theta_{z,t}\). Initially permit zero components of \(t\);
they are removed in the theorem below. The factor
\(\Theta_{\rm top}\) and its normalization are fixed once for
all these sections. Put \(s=\|\sum_{z,t}\Theta_{z,t}\|_{\rm av}\).
The sum is over all normal rows, including zero residues. The
following form of the theta kernel calculation keeps every
coefficient of the tests at \(g_1,g_2\).

\begin{lemma}[A derivative estimate with some residue coordinates retained]
\label{lem:partialtheta390}
Let a polarized torus with period lattice \(\Gamma\) in normal
coordinates have a basis \(\Theta_{z,t}\) of the kind in
Proposition~\ref{prop:normalcells388}. The finite position group
is \(Z\times T_0\), where \(T_0\) is any product of the other
varying residue groups. Put \(q_Z=|Z|\) and
\(s=\|\sum_{z,t}\Theta_{z,t}\|_{\rm av}\). Let \(M\) be a
positive integer annihilating \(Z\) and suppose that all the
translations and modulations on \(Z\), acting trivially on
\(T_0\), have lifts to normalized theta translations by
elements of \(M^{-1}\Gamma\), up to constant phases. Distinct
pairs of translations and modulations have distinct lifts
modulo \(\Gamma\). Let \(R\) be a lower bound for the length of
a nonzero vector of \(M^{-1}\Gamma\), with
\(\pi R^2\ge\max\{64L,16D\log5,8\log2\}\).
For any subset of derivative indices of degree at most \(L\) put
\[
  H_{(\alpha,z),(\alpha',z')}
      =\sum_{t\in T_0}\ell_\alpha(\Theta_{z,t})
                              \overline{\ell_{\alpha'}(\Theta_{z',t})}.
\]
Then \(\|q_Zs^{-2}H-I\|\le2e^{-\pi R^2/8}\).
\end{lemma}
\begin{proof}
For \(a\in Z,\ b\in Z^\vee\) let
\(W_{a,b}e_z=b(z)e_{z+a}\) act on the ordinary coordinate
basis of \(\mathbb C^Z\). Here \(Z^\vee\) denotes the complex
character group. These matrices are unitary. They are mutually
orthogonal for the trace inner product, with squared norm \(q_Z\).
Indeed a nonzero shift has zero diagonal, and a nontrivial
character has sum zero. There are \(q_Z^2\) matrices, so they
form a basis of all matrices on \(\mathbb C^Z\). In the order
of indices \(z,\alpha\), expand
\[
 H=q_Z^{-1}\sum_{a,b} W_{a,b}\otimes H_{a,b}.
\]
The derivative matrix \(H_{a,b}\) is obtained by multiplying
the \(z,z'\) block of \(H\) by
\(\overline{(W_{a,b})_{zz'}}\) and summing over \(z,z'\).

Set \(Q=q_Z|T_0|\). The theta basis for probability Haar
measure is \(\sqrt Q\,\Theta_{z,t}/s\), by the
orthogonality and norm calculation of
Proposition~\ref{prop:normalcells388}. Relative to this basis,
acting on the first variable of the reproducing kernel by
\(\overline W_{a,b}\otimes I\) multiplies the coefficient of
the \(z,z'\) term by the matrix entry just displayed.
Consequently \(H_{a,b}/s^2\) is the normalized derivative matrix
of this translated kernel divided by \(Q\). Complex conjugation
of \(W_{a,b}\) replaces \(b\) by \(b^{-1}\) and gives another
member of the same family.

Let \(y\) be a lift of this operator. In the period expansion
of Lemma~\ref{lem:thetaframe389}, composition with the normalized
translation by \(y\) replaces the terms indexed by \(\Gamma\)
by terms indexed by a coset of \(\Gamma\) in \(M^{-1}\Gamma\).
For example, composing the operators
\[
     U_vF(u)=e^{-\pi u\cdot\bar v-\pi|v|^2/2}F(u+v)
\]
gives \(U_vU_w\) equal to \(U_{v+w}\) times a scalar of
absolute value one, by direct multiplication of exponentials.
Thus every derivative matrix in the coset is bounded in norm by
the same compression estimate used in that lemma. The identity
matrix comes from the zero lattice vector at \((a,b)=(0,1)\).
All other lattice vectors, including the nonzero terms of the
identity coset, occur at most once in these estimates, since
the cosets are distinct.
After multiplication by \(q_Z/s^2\), use \(\|W_{a,b}\|=1\)
and the triangle inequality in the expansion of \(H\).
The error is bounded by the sum over the nonzero vectors of
\(M^{-1}\Gamma\) of their compressed translation norms.
The shell estimate in Lemma~\ref{lem:thetaframe389} bounds this
sum by \(2e^{-\pi R^2/8}\). This proves the assertion.
\end{proof}

\begin{theorem}[The two marked primes with nonzero terminal covectors]
\label{thm:terminalmarked390}
Use the actual fields, weights and three primes on the intervals
of Section~\ref{sec:missing}. For \(\alpha,\alpha'\) in
\(\mathcal I\) of Theorem~\ref{thm:normaljet388} form the matrix
\[
  H^\times_{(\alpha,z),(\alpha',z')}
       =\sum_{\substack{t\\ t_w\ne0\ (w\mid\pi_T)}}
         \ell_\alpha(\Theta_{z,t})
                 \overline{\ell_{\alpha'}(\Theta_{z',t})}.
\]
For every sufficiently large \(x\) on these intervals,
\[
          \left\|\frac{Q_g}{s^2}H^\times-I\right\|
                                 \le e^{-x^{\varepsilon/8}} .
\]
The bound is uniform in the fixed data specified above and
in the transported complex realization with its Hermitian metric.
\end{theorem}
\begin{proof}
We treat both the full table and the table obtained by fixing
one component \(t_w=0\). Denote by \(\mathcal T'\) the set of
terminal places still varying. At every varying place the
normal test is a residue characteristic function. The fixed
matrix cells project independently onto the small and normal
rows, by the Chinese remainder argument of
Proposition~\ref{prop:markedsampling387}. At a place over
\(\pi_T\), the fixed small part must be integral. With this
condition, imposing the covector equation at \(T\) only
specifies the residue normal row, since the lifted terminal
covector is \((\bar\pi_T\bmod w)e^\vee\). The modulus \(c_B\)
has exponent one there. All other frame and coefficient
operators at this edge depend on the underlying lattice and
its frames away from \(T\). This is the finite test
\(\phi_T\) in
\srcloc{Formalism/LocalToGlobal/Columns/FiniteSumNonzeroLargeDepth}{646}
and the lifted class in lines 203--218 there.

Apply the proof of Proposition~\ref{prop:normalcells388} to
these places. The dual residue characters give orthogonality,
the position translations give equal norms, and the index of
the common period lattice is the product of the numbers of
cells. In the present coordinates that lattice is
\[
    \Gamma=\delta\Phi(\mathfrak b^{-1}\mathfrak a
                                \mathfrak c N^\vee),\qquad
    \mathfrak c=(\bar\pi_{g_1}\bar\pi_{g_2})
                                \prod_{w\in\mathcal T'} w .
\]
The frozen places contribute the norm-one ideal \(\mathfrak a\)
as before. In the table with \(t_w=0\) this includes the
change by \(w/\bar w\) at the newly frozen place. A finite
translation or modulation at the marked places acts on
the characteristic functions indexed by \(z\) as \(W_{a,b}\),
and fixes the terminal residue index. To see the asserted
lift condition of Lemma~\ref{lem:partialtheta390}, prescribe
an integral position shift at a marked half and a frequency
shift in \(g_i^{-1}U^\circ\) at its conjugate. Modulo the
common periods these shifts index the translations and
modulations, respectively. Prescribe zero at all other
places. Multiplication by \(G\) takes the whole tuple into
the local common periods. The Chinese remainder theorem for
the fractional normal module gives a lift in \(G^{-1}\Gamma\).
The normalized Weyl action of that lift agrees with the finite
one up to a unit scalar, since changing representatives by
any common period acts on all sections by the same scalar.
The two choices of ordering a position shift and a modulation
also differ by a unit scalar. Distinct pairs have distinct
classes, either by the local residue coordinates or by the
trace orthogonality in the lemma.

We bound the shortest vector even after omitting one prime
ideal above \(T\). For a fractional ideal \(\mathfrak j\) and
nonzero \(v\in\mathfrak b^{-1}\mathfrak a\mathfrak j N^\vee\),
the totally positive scalar \(h(v,v)\) belongs to
\(\mathfrak j\bar{\mathfrak j}\,\mathfrak D_{E/K}^{-1}\).
This follows locally from the integrality of the Hermitian
pairing on \(N^\vee\) and the defining property of
\(\mathfrak b\). Its absolute field norm is at least the norm
of this fractional ideal, by ideal divisibility. Applying the
arithmetic--geometric mean inequality to its values at the
\(f\) CM embeddings and using the metric normalization in
Theorem~\ref{thm:lattice55} gives
\[
   |\delta\Phi(v)|_H^2
        \ge \frac{2f}{\operatorname{rd}(E_x)}
                         (\operatorname N\mathfrak j)^{1/f}.
\]
The same calculation with the transported differential has
the same right-hand side as in Theorem~\ref{thm:normalframe389}.
For \(G^{-1}\Gamma\) take \(\mathfrak j=G^{-1}\mathfrak c\).
Since the fields split completely at the three primes,
\[
  (\operatorname N\mathfrak j)^{1/f}
            =\frac{T^{|\mathcal T'|/f}}{G}.
\]
There are either \(f\) or \(f-1\) varying places, and \(f\ge2\)
on the intervals considered. Hence a common lower bound for
the squared length is
\[
          R_*^2=\frac{2f\sqrt T}{G\,\operatorname{rd}(E_x)} .
\]
Here \(G\le4(kb)^{16}\), \(\log T\gg m x^\varepsilon\) with
\(m\to\infty\), and \(\log\operatorname{rd}(E_x)
       \ll_{\psi,\varepsilon}x^{\varepsilon/4}\), by
Proposition~\ref{prop:ray50} and the size estimates in the
proof of Theorem~\ref{thm:terminalformula347}. Consequently
\(R_*^2\) dominates \(L=\sum_\tau B_\tau\) and \(D\) as
required by Lemma~\ref{lem:partialtheta390}.
Put \(\delta_x=2\exp(-\pi R_*^2/8)\).

On the full table the lemma bounds its Gram matrix between
\((1\pm\delta_x)s^2 I/Q_g\), with the minus and plus signs
giving the lower and upper bounds, respectively. For the
single face \(t_w=0\), the squared norm of the sum of its
normal sections is \(s^2/T^d\). Indeed, in the full normal
basis the face includes exactly a fraction \(T^{-d}\) of all
positions, each of the same squared norm. Its norm can also be
computed on the torus for the frozen face, since probability
integration of a periodic metric density agrees on a finite
cover. The Gram matrix of that face is therefore bounded above
by \((1+\delta_x)s^2 I/(Q_gT^d)\). Each summand indexed by
\(t\) is positive semidefinite. Removing the union of the \(f\)
faces subtracts at most their sum in quadratic-form order.
Thus the relative error for \(H^\times\) is at most
\(\delta_x+(1+\delta_x)fT^{-d}\). Both terms together are at
most \(e^{-x^{\varepsilon/8}}\) for sufficiently large \(x\),
by the bounds on \(T,R_*^2\) and polynomial growth of \(f\).
\end{proof}

\begin{corollary}[The signed sum at both marked primes on a fixed incoming table]
\label{cor:signedterminal390}
Let \(Y(z)\) be a fixed incoming coefficient vector for each
marked residue tuple, with zero values permitted. In the fixed
complex coefficient space before imposing the preceding
terminal-incidence conditions, put
\[
       P_t(Y)=\sum_{z,\alpha}
                 \ell_\alpha(\Theta_{z,t})\,A_\alpha Y(z),
       \qquad
       M(Y)=\frac{s^2}{Q_g}\sum_{z,\alpha}\|A_\alpha Y(z)\|^2 .
\]
With \(\rho_x=e^{-x^{\varepsilon/8}}\),
\[
  (1-\rho_x) M(Y)
       \le \sum_{t_w\ne0\ \forall w}\|P_t(Y)\|^2
       \le (1+\rho_x) M(Y).
\]
In particular a nonzero \(A_\alpha Y(z)\) for any one pair
implies a nonzero projected coefficient at some tuple of
nonzero terminal covectors. This applies to the signed marked
table \(Y(z)=g(z)\widehat C(z)\) of
Proposition~\ref{prop:markedsampling387} whenever the prescribed
incoming table is common to the varying terminal refinements.
\end{corollary}
\begin{proof}
For each \(t\), summing the signed marked coefficients in the
last edge gives \(P_t(Y)\) by the substitution and branch
projection of Theorem~\ref{thm:normaljet388}. The normal cells
over \(T\) alter only the normal theta sections in that
substitution. Apply Theorem~\ref{thm:terminalmarked390} to the
conjugates of the scalar coordinates of the vectors
\(A_\alpha Y(z)\) in an orthonormal output basis, and sum the
result over that basis. This gives the two inequalities and
the nonvanishing assertion.

In the full column of Theorem~\ref{thm:covectorpaths352}, an
earlier edge with reduction \(\bar X\) requires its incoming
covector to be the pullback of the outgoing one by \(\bar X\).
Thus an initial input can select different earlier contributions
for different \(t\). The full signed period then also sums the
resulting outputs over \(t\). The displayed estimate is for
the specified common table and the sum of squared outputs.
Control of those pullbacks and of the cross terms in the
prescribed period is still required in Section~\ref{sec:missing}.
\end{proof}


\subsection{The terminal sum for a prescribed initial covector}
\label{sec:normalpullback391}

Fix the earlier matrices of a path through the penultimate
edge in Theorem~\ref{thm:covectorpaths352}. The final edge,
its two marked normal rows and its terminal normal covectors
are still to be summed. At \(w\mid\pi_T\), let \(V_{0,w}\)
be the initial residue space and \(V_w=N/wN\) the intermediate
residue space of dimension \(d\). Define
\[
 C_w=(\bar\pi_T\bmod w)\,
                 \bar x_{m-2,w}\cdots\bar x_{0,w}
                   :V_{0,w}\longrightarrow V_w,\qquad
 r_w=\dim\operatorname{im} C_w .
\]
For the terminal normal row \(t_w\in V_w^\vee\), the initial
covector is \(u_w=t_w C_w\). If \(u_w\ne0\), all the intermediate
pullbacks of \(t_w\) are nonzero. Indeed if one pullback were
zero, every further pullback towards the initial class would
also be zero. If \(u_w=0\) at any \(w\), the refined path is
absent. The other factors of each earlier matrix operator at
fixed underlying lattices and fixed frames away from \(T\)
are independent of the refinements, by the finite class formula
of Theorem~\ref{thm:covectorpaths352}. We may therefore use
\(P_t(Y)\) of Corollary~\ref{cor:signedterminal390} for a
fixed incoming vector acted on by those operators. All values
of its marked table \(Y\) enter with their original signs.

\begin{proposition}[The cost of the preceding incidence conditions]
\label{prop:incidencenorm391}
Assume \(r_w\ge1\) for every \(w\mid\pi_T\) and put
\(p_C=\sum_{w\mid\pi_T}T^{-r_w}\). With \(M(Y)\) as in
Corollary~\ref{cor:signedterminal390} and
\(\delta_x=2\exp(-\pi R_*^2/8)\) as in the proof of
Theorem~\ref{thm:terminalmarked390}, one has
\begin{align*}
 [1-\delta_x-(1+\delta_x)p_C] M(Y)
   &\le\sum_{t:\, t_w C_w\ne0\ \forall w}\|P_t(Y)\|^2\\
   &\le (1+\delta_x)M(Y).
\end{align*}
Multiplication of each summand \(P_t(Y)\) by a scalar of
absolute value one depending on \(u=(t_wC_w)_w\) preserves
both inequalities. If some \(r_w=0\), the path contributes
zero for every admissible initial covector.
\end{proposition}
\begin{proof}
With no restrictions on \(t\) the two Gram bounds with error
\(\delta_x\) occur before the faces are removed in the proof of
Theorem~\ref{thm:terminalmarked390}. The estimate on a face
there also holds for \(t_w=v\) for any fixed \(v\in V_w^\vee\).
Replacing zero by \(v\) changes the unitary character of the
frozen cell. Its period module, norm \(s/T^{d/2}\), and shortest
vector bound are the same. The sum of squared outputs on such
a face is at most \((1+\delta_x)T^{-d}M(Y)\).
There are \(T^{d-r_w}\) choices of \(v\) annihilating
\(\operatorname{im}C_w\), by rank and nullity over the residue
field. Removing all forbidden faces at \(w\) consequently costs
at most \((1+\delta_x)T^{-r_w}M(Y)\). Adding over the places
bounds the union since each squared norm is nonnegative. This
proves the inequalities. The assertion about phases follows
from the Hermitian norm. For rank zero the initial pullback
vanishes for every terminal row.

For example, start on one unrefined initial class with trivial
stabilizer as in Proposition~\ref{prop:column110}. At each
terminal place choose one representative on each nonzero
projective line of initial covectors and a nontrivial character
\(\xi_w:k_w^\times\to\mathbb C^\times\). Assign to
\(c\) times a representative the scalar \(\xi_w(c)\).
Multiply the resulting local functions and a fixed input
vector. This is a legitimate initial test on that class with
zero sum separately on every scalar orbit, by multiplicative
orthogonality. On each permitted path it contributes exactly
one of the phases in the proposition. The estimate applies
before addition of outputs with different \(t\).
\end{proof}

\begin{proposition}[Coherent summation on the fibers of the initial pullback]
\label{prop:coherentcell391}
Write \(C_w^*:V_w^\vee\to V_{0,w}^\vee\) for \(t\mapsto tC_w\).
Put \(K_w=\{y\in U_w:\bar y C_w=0\}\). For
\(u_w\in\operatorname{im} C_w^*\) choose any integral lift
\(a_w\) of a preimage of \(u_w\).
For fixed \(z\), the sum of the normal sections
\(\Theta_{z,t}\) over all \(t\) with \(t_wC_w=u_w\) is the
section with terminal test
\(\prod_{w\mid\pi_T}{\bf1}_{a_w+K_w}\). For \(u_w\ne0\)
at all these places, applying the branch map and summing the
marked inputs gives the coefficient of the original initial
covector \(u\) along this path. An initial covector outside
the image gives coefficient zero.

Including zero labels for this dimension calculation, the cells
indexed by \((z,u)\in Z\times\prod_w\operatorname{im}C_w^*\)
form an orthogonal basis of the space of normal theta sections.
This space has dimension \(Q_gT^{\sum_w r_w}\). Relative to
the periods for the integral normal vacuum, its position
lattice at each terminal half is \(K_w\) and its frequency
lattice at the conjugate half is \(U_w^\circ\).
\end{proposition}
\begin{proof}
The preimage of \(u_w\) on integral positions is a coset of
\(K_w\). Its disjoint decomposition into residue cells has
one term for each \(t_w\) mapping to \(u_w\). Addition of
their characteristic functions is exact, so the finite theta
intertwiner gives the stated section by linearity, with the
original measure and no additional multiplier. The earlier
refinements are uniquely specified by these \(t_w\), and
the nonzero initial covectors ensure that all of them pass.
This proves the statement about the coefficient. In particular
the terms in this fiber are added before taking its norm.

For the basis assertion, translations by \(K_w\) preserve
every such cell, and the frequencies of \(U_w^\circ\) act
trivially on integral positions. The characters of \(U_w/K_w\)
are represented by the dual lattice of \(K_w\) modulo
\(U_w^\circ\). They distinguish the cells and act unitarily.
Integral translations permute the \(T^{r_w}\) cells and
preserve their norms. The marked halves give \(Q_g\) further
positions exactly as in Proposition~\ref{prop:normalcells388}.
The common period lattice has covolume the product of these
indices. Its positivity and the Riemann--Roch dimension
formula give the stated number of sections. Their sum is
the nonzero normal vacuum with the frozen tests. Equal norms
and orthogonality now give a basis.
\end{proof}

\begin{theorem}[An initial-covector estimate for a common orthogonal image]
\label{thm:splitcoherent391}
Suppose \(N=N_1\perp N_2\) as Hermitian
\(\mathcal O_E\)-lattices, with both summands self-dual and
\(\operatorname{rank}N_1=a\ge1\). Suppose that at every
\(w\mid\pi_T\), \(\operatorname{im}C_w=N_1/wN_1\) in \(V_w\).
Use coordinates adapted to this splitting at the two marked
primes, and write \(z=(z_1,z_2)\), \(Q_1=G^{fa}\). Let
\(\Phi_{z_1,t_1}\) be the normal theta sections for \(N_1^\vee\),
with the residue covector \(t_{1,w}\in(N_1/wN_1)^\vee\)
at each terminal place. Let \(\Psi_{z_2}\) be the sections
for \(N_2^\vee\), with the full integral normal row at those
places. Use the product normalization from the full normal
vacuum and put \(s_1=\|\sum_{z_1,t_1}\Phi_{z_1,t_1}\|_{\rm av}\).
In orthonormal coordinates adapted to \(N_1,N_2\), put
\[
  V_{\alpha_1,z_1}
    =\sum_{\substack{\alpha_2,z_2\\
                       (\alpha_1,\alpha_2)\in\mathcal I}}
        \ell_{\alpha_2}(\Psi_{z_2})\,
                    A_{\alpha_1,\alpha_2}Y(z_1,z_2),\qquad
  M_1=\frac{s_1^2}{Q_1}\sum_{\alpha_1,z_1}\|V_{\alpha_1,z_1}\|^2 .
\]
Let \(D_u=\sum_{t:\,t_w C_w=u_w\ \forall w}P_t(Y)\) for each
admissible initial tuple \(u\). On the intervals of
Section~\ref{sec:missing} one has, with \(\rho_x=e^{-x^{\varepsilon/8}}\),
\[
        (1-\rho_x)M_1
       \le \sum_{u_w\ne0\ \forall w}\|D_u\|^2
       \le (1+\rho_x)M_1
\]
for all sufficiently large \(x\). In particular some initial
covector coefficient on this path survives its full signed
sum over the terminal fiber if and only if some
\(V_{\alpha_1,z_1}\ne0\).
\end{theorem}
\begin{proof}
The surjection \(C_w:V_{0,w}\to N_1/wN_1\) makes
its dual injective. Every \(u_w\) in its image specifies
one \(t_{1,w}\), and \(u_w\ne0\) is equivalent to
\(t_{1,w}\ne0\). The normal row on \(N_2\) ranges freely.
By Proposition~\ref{prop:coherentcell391}, its terminal test
after summation is the full integral row. Changing to coordinates
adapted to the summands at the two marked halves merely
reindexes the residue cells. The marked function \(Y(z_1,z_2)\)
retains all its entries and may depend jointly on both.

The orthogonal direct sum splits the period form, the quadratic
semicharacter and the matrix lattice. At every other frozen
split place a fixed congruence coset also splits into the two
projected cosets, since its modulus is scalar. The argument
of Proposition~\ref{prop:normalcells388} now gives the product
\(\Phi_{z_1,t_1}\Psi_{z_2}\) for the summed normal section.
In a product of orthonormal coordinate frames
\(\ell_{(\alpha_1,\alpha_2)}(\Phi\Psi)=
\ell_{\alpha_1}(\Phi)\ell_{\alpha_2}(\Psi)\), by ordinary
differentiation in disjoint variables and the product formula
for multi-index factorials. Regrouping the finite sums gives
\[
       D_u=\sum_{\alpha_1,z_1}
                 \ell_{\alpha_1}(\Phi_{z_1,t_1})\, V_{\alpha_1,z_1}.
\]
This expression is zero outside the image of the pullback.

Apply Lemma~\ref{lem:partialtheta390} to the \(fa\) normal
variables of \(N_1^\vee\). The degree cutoff for all the
\(\alpha_1\)'s is at most \(L=\sum_\tau B_\tau\). The marked
group has order \(Q_1\) and exponent dividing \(G\). Its
translations and modulations have lifts in \(G^{-1}\Gamma_1\)
by the same local congruences as in
Theorem~\ref{thm:terminalmarked390}. For the full terminal
table, and also on freezing any one terminal component, the
norm-ideal proof of that theorem gives the same lower bound
\(2f\sqrt T/(G\operatorname{rd}(E_x))\) for the squared
length. It used self-duality and integrality of the scalar
Hermitian pairing, both of which hold for \(N_1^\vee\).
Only the dimension of the individual terminal residue space
changes. Each removed zero face costs a fraction at most
\((1+\delta_x)T^{-a}\) of the reference quadratic form.
Since \(a\ge1\), the total relative error
\(\delta_x+(1+\delta_x)fT^{-a}\) is at most \(\rho_x\)
on the same intervals. Apply the resulting Gram bounds to
the vectors \(V_{\alpha_1,z_1}\) coordinate by coordinate
to obtain the displayed inequalities.

For instance, on standard unrefined lattices the residue
matrices of the earlier edges at \(T\) may be the coordinate
inclusions of successive ranks. Their composite image is the
first \(N(0)\) coordinates, an orthogonal summand at every
place. These choices of residue matrices can be imposed along
with any prescribed matrix cosets at the other selected primes
by the Chinese remainder theorem. This describes actual
indices in the finite sums. A nonzero coefficient there also
requires the displayed test involving \(V_{\alpha_1,z_1}\).
For a general path Proposition~\ref{prop:coherentcell391} still
gives its exact cells. Their positions at \(T\) can contain
undilated directions in the kernel of the pullback. A lower
bound on the full auxiliary derivative matrix alone cannot
evaluate those directions or the interference between paths
in the final Bessel period.
\end{proof}


\subsection{The number of exceptional intermediate matrix residues}
\label{sec:intermediateimages392}

Take \(E_x,g_1,g_2,T\) as in Section~\ref{sec:missing} and the
word of Theorem~\ref{thm:common111}. Use \(d=N(m-1)\),
\(D=fd\), \(G=g_1g_2\) and \(b=\exp(x^\varepsilon)\).
We first count residue images, giving each matrix residue equal
weight. The theta coefficients attached to those matrices will
enter the estimate for an individual path. The prime sizes
needed here are
\[
       \log G=(16+o(1))\log b,\qquad
       \log T=(17+o(1))m\log b
\]
by Theorem~\ref{thm:prime209}. Also \(f/m\to\infty\),
\(\log \operatorname{rd}(E_x)=o(\log b)\), and the logarithms
of the ranks are \(O(\log x)=o(\log b)\). These follow from
Theorems~\ref{thm:rangenew50} and~\ref{thm:conductor232}.

\begin{lemma}[A count of images admitting a short normal period]
\label{lem:shortimage392}
At each \(w\mid\pi_T\) choose independently and uniformly
an \(a\)-dimensional subspace \(I_w\subset V_w\), where
\(a\le d\) and \(am\ge2d\). In
Proposition~\ref{prop:coherentcell391} use the position lattice
\(K_w=\{y\in U_w:\bar y|_{I_w}=0\}\). Write \(\Gamma_I\)
for the common period lattice of these cells and the marked
residue sections at \(g_1,g_2\). Also consider, for each
\(w_0\mid\pi_T\), its period lattice on fixing the cell
at \(w_0\) and leaving the others free. The proportion of
choices for which any of these lattices, after division by
\(G\), has a nonzero vector of squared length less than
\(b^{3/2}\) is at most
\[
                    \exp(-10fd\log b)
\]
for sufficiently large \(x\). This bound also holds if all
the transported complex realizations of the CM period data
are tested together.
\end{lemma}
\begin{proof}
Work first in one realization. Let \(\Gamma_0\) be the normal
period lattice obtained by leaving the integral position row
unconstrained at the terminal places, retaining both marked
residue groups. Thus, for the fixed norm-one ideal
\(\mathfrak a\) coming from the other frozen cells,
\[
     \Gamma_0=\delta\Phi(\beta\mathfrak b^{-1}\mathfrak a N^\vee),
              \qquad \beta=\bar\pi_{g_1}\bar\pi_{g_2}.
\]
When all labels vary, \(\Gamma_I\subset\Gamma_0\). If a
single terminal label at \(w_0\) is fixed, its dual
frequencies may gain a denominator at \(\bar w_0\). They
lie in \(T^{-1}U_{w_0}^\circ\), since \(K_{w_0}\) contains
\(TU_{w_0}\). The period lattice on that face is therefore
contained in the module defining \(\Gamma_0\) enlarged by
the ideal \(\bar w_0^{-1}\). These ambient lattices are
independent of the subspaces just chosen.

Let \(e=0\) on the full table and \(e=1\) on one of the
faces. After division by \(G\), the nonzero vectors of the
ambient lattice have squared length at least
\[
    R_e^2=\frac{2f}{G\,\operatorname{rd}(E_x)}T^{-e/f}.
\]
This is the fractional-ideal norm inequality in the proof
of Theorem~\ref{thm:terminalmarked390}. If a vector also
reduces to zero at \(h\) of the unfrozen terminal halves,
its squared length is at least \(R_e^2T^{h/f}\). Indeed
membership in their prime ideals multiplies the norm of the
fractional ideal in that calculation by \(T^h\).
At a terminal half all the factors from \(G,\beta\) are units.
Thus the vanishing of this row means exactly that its
reduction is zero in the indicated normal coordinates.

Put \(R=b^{3/4}\). Disjoint open balls of radius \(R_e/2\)
about the ambient lattice vectors of norm less than \(R\)
lie in the ball of radius \(R+R_e/2\). There are at most
\((1+2R/R_e)^{2D}\) such vectors. Since
\(\log T/f=o(\log b)\), their number is at most
\[
                \exp\{(17.5+o(1))fd\log b\}.
\]
The constant \(17.5\) comes from
\(\log(R^2/R_e^2)\le(3/2+16+o(1))\log b\).
For a nonzero vector in this ball, the same inequality and
the extra factor \(T^{h/f}\) imply
\(h/f\le(17.5+o(1))\log b/\log T=O(1/m)\).

At an unfrozen place with nonzero residue row \(y\), the
probability that \(I_w\) is contained in \(\ker y\) is
\((T^{d-a}-1)/(T^d-1)\le T^{-a}\). To compute it, count
pairs consisting of an \(a\)-plane and a nonzero form
annihilating that plane. Every plane has \(T^{d-a}-1\)
such forms, and the general linear group is transitive on
the \(T^d-1\) nonzero forms. By independence, the probability
that this one ambient vector meets all the necessary conditions
is at most \(T^{-a(f-e-h)}\). The logarithm of this bound is
at most \(-(34-o(1))fd\log b\), using \(am\ge2d\). A union
bound over the ambient vectors and the \(f+1\) possible
tables proves the assertion in one realization with an
exponent \(-(16.5-o(1))fd\log b\).

For completeness we bound the number of metric lattices to
be considered after field conjugation. In product coordinates
Theorem~\ref{thm:transport59} gives a lattice of the form
\(\delta_\rho\mathfrak c\,\rho(M)\), with metric scale
\(\operatorname{area}(\delta_\rho\mathfrak c)\); here \(M\)
is the indicated fractional period module and \(\mathfrak c\)
is an ideal of the imaginary quadratic field \(K\). Changing
the representative of this ideal class and the scalar
\(\delta_\rho\) multiplies all coordinates by a complex
scalar and the metric scale by the square of its absolute
value. The resulting lengths are identical. The coordinate
change to the trace metric is unitary after transport. Its
entries and the field defining the row lattice are in the
abelian field \(E_x\). Their field embeddings have at most
\(2f\) possibilities. On torsion the further choice of
generator of the elliptic CM module multiplies a residue
module by a unit scalar. A subspace at a terminal place,
and the condition that a row annihilate it, are preserved
by this scalar. Unit phases on theta sections leave their
period lengths unchanged. Thus at most \(2f h_K\) cases of
the count suffice, where \(h_K\) is the class number of \(K\).
Each transport permutes or conjugates the residue parameters
and preserves the uniform count.

By the imaginary quadratic Minkowski bound, each ideal
class has an integral representative of norm
\(\ll\sqrt{|\Delta_K|}\). There are at most
\(\sum_{n\le B}\sum_{d\mid n}1\le B(1+\log B)\) ideals
of norm at most \(B\), using the Euler factors of a quadratic
field. Hence \(\log h_K=O(\log(2|\Delta_K|))=o(\log b)\)
on the common intervals. A union bound over the metric
cases still gives the stated exponent.
\end{proof}

\begin{theorem}[Coherent sums with a fixed intermediate pullback]
\label{thm:intermediateimage392}
Choose a vertex \(2\le j\le m-2\). Fix the underlying
lattices and frames of the path and its matrix data away from
the terminal prime. Among the remaining terminal matrix
residues for the edges from \(j\) up to \(m-1\), give equal
weight to every choice with injective residue maps at every
terminal place. Except for a proportion at most
\(\exp(-10fd\log b)\), the following assertion holds.
Write \(C_w\) for the composite from vertex \(j\), including
the unit factor \(\bar\pi_T\bmod w\). For every fixed marked
input \(Y(z)\) write
\[
 D_u=\sum_{t:\,t_wC_w=u_w\ \forall w}P_t(Y),\qquad
 M(Y)=\frac{s^2}{Q_g}\sum_{z,\alpha}\|A_\alpha Y(z)\|^2 .
\]
On the nonzero intermediate covectors \(u\) at vertex \(j\),
in each transported complex realization,
\[
 (1-\rho_x)M(Y)\ \le\
           \sum_{u_w\ne0\ \forall w}\|D_u\|^2
     \ \le\ (1+\rho_x)M(Y),\qquad
                \rho_x=e^{-x^{\varepsilon/8}} .
\]
The exceptional fraction is a counting bound for actual
matrix residues before weighting their theta coefficients.
\end{theorem}
\begin{proof}
The source and target lattices at a terminal half are free
over the local ring. The matrix quotient defining each edge
maps uniformly onto the full residue matrix space there by
the Chinese remainder theorem. The data away from the
terminal prime can be prescribed independently, as in
Proposition~\ref{prop:markedsampling387}. The intermediate
covectors have yet to be imposed. On the stated set,
the composites are injective. The image of the composite
at one place is uniform on the Grassmannian of planes of
dimension \(a=N(j)\). Indeed, multiplying the last injection
on the left by an element of \(\operatorname{GL}(V_w)\)
permutes the choices and acts transitively on the images.
The different terminal places are independent. The unit
in \(C_w\) preserves its image. The ratios of ranks meet the
requirement of the lemma, since
\[
       d=k+(m-1)(t_{\rm r}+1),\qquad
       a=k+j(t_{\rm r}+1),\qquad am\ge2d .
\]
Here \(j\ge2\) and \(m\ge2\). The bounds are uniform for all
such \(j\). The rank restriction defines a nonempty sample
at every place. For consecutive ranks \(n<n'\), for example,
its fraction in the entire matrix space is
\(\prod_{i=0}^{n-1}(1-T^{i-n'})\), by counting ordered
independent columns.

For an image outside the exceptional set, sum its terminal
cells with the same pullback. Proposition~\ref{prop:coherentcell391}
identifies the resulting test with a single coset of \(K_w\)
and supplies a complete normal basis including zero labels.
The same proof applies starting at vertex \(j\), since
covectors on the successive edges pull back by the same
matrix equation. The basis has \(T^{fa}\) terminal entries
for each \(z\). Its whole sum is the original integral
normal row, of norm \(s\). The marked translations and
modulations on \(Z\) have representatives in
\(G^{-1}\Gamma_I\). To see this, prescribe their residue
coordinates at the marked places and zero at \(T\), and
apply the Chinese remainder theorem in \(G^{-1}\Gamma_I/\Gamma_I\).
The same reasoning applies on any frozen terminal face.
The square of the shortest length in all these enlarged
lattices is at least \(b^{3/2}\) by the lemma.
The bounds \(L\ll x^{C}b\) and \(D\ll x^{C}\) of
Theorem~\ref{thm:normalframe389} now give the degree and
dimension hypotheses of Lemma~\ref{lem:partialtheta390}.
The resulting relative error on the full table is at most
\(\delta=2e^{-\pi b^{3/2}/8}\).

Fixing one zero component of \(u\) selects exactly a fraction
\(T^{-a}\) of the orthogonal equal-norm sections. On that face
the same lemma gives an upper Gram bound of
\((1+\delta)T^{-a}s^2 I/Q_g\). Removing the union of the
zero faces therefore costs at most \(f(1+\delta)T^{-a}\)
in relative norm. Together with \(\delta\) this is at most
\(\rho_x\). The result follows by applying the bounds to the
coordinates of \(A_\alpha Y(z)\) after the fiber sum, exactly
as in Corollary~\ref{cor:signedterminal390}.
\end{proof}

\begin{proposition}[A dimension obstruction at the initial vertex]
\label{prop:initialrank392}
Fix one scalar output on a path, and consider the map from
arbitrary scalar weights on its marked frequencies to its
coefficients indexed by the original nonzero initial covectors.
Its rank is at most \(T^{fk}\). If the determinant varies at
every marked place, its marked frequency support \(\mathcal Z\)
satisfies
\[
         |\mathcal Z|\ge G^{f t_{\rm r}},\qquad
    \frac{T^{fk}}{|\mathcal Z|}
            \le \exp(-8f t_{\rm r}\log b)
\]
for sufficiently large \(x\). Even on this support the kernel
in the space of arbitrary scalar weights has dimension at
least \(|\mathcal Z|-T^{fk}\).
\end{proposition}
\begin{proof}
At each terminal half the initial covector belongs to a
space of dimension \(k=N(0)\) over the field of \(T\)
elements. There are at most \(T^{fk}\) resulting scalar
coefficients, counting even those that vanish by the
pullback condition. This proves the rank bound. At a marked
place with varying determinant the support has size
\((q_w-1)q_w^{d-s_w}\) by Proposition~\ref{prop:markedsampling387}.
Here \(s_w\le r+(m-2)t_{\rm r}=d-t_{\rm r}\) by
Proposition~\ref{prop:markedcell386}. Taking the product
over the complete halves at both primes gives the first
inequality. The ranks in Theorem~\ref{thm:rangenew50}
satisfy \(km/t_{\rm r}\to0\), directly from the factor
\(\lceil x^{377}\rceil\) in \(t_{\rm r}\) and the stated
powers of \(x\) for \(k,m\). The two prime asymptotics above
give
\[
 fk\log T-ft_{\rm r}\log G
                 =-(16+o(1))ft_{\rm r}\log b .
\]
This gives the second inequality, and rank-nullity finishes
the proof. The calculation concerns arbitrary weights on the
support. For the arithmetic input, the coefficients are the
specific Fourier values of the preceding signed path sum.
Their position relative to this kernel, and their interactions
across paths with common initial covectors, enter the period
comparison of Section~\ref{sec:missing}.
\end{proof}


\subsection{The signed marked table for a fixed penultimate image}
\label{sec:markedimage393}

The rank calculation of Proposition~\ref{prop:initialrank392}
allows an independent coefficient at every frequency in the
marked support. We now determine the dependence imposed on
those coefficients by the actual matrix test. Fix one marked
place with residue field \(k\) of order \(q\). Use
\(V=k^d\), \(d=N(m-1)\), \(a=N(m-2)+1\), and the subspace
\(A=A_a\subset V\) of the first \(a\) coordinates in the
intermediate frame of Proposition~\ref{prop:markedcell386}.
Thus \(d-a=t_{\rm r}\). Retain a nonempty position cell \(D_0\)
on which some matrix passes all tests. For a varying determinant
let \(i\) be its block, \(W=A_{s_i}\), \(s=\dim W\),
\(H=\operatorname{im}\bar X^{(a)}\subset A\) and \(K_0=H\cap W\).
Here \(\bar X^{(a)}\) consists of the first \(a\) rows of the
reduced penultimate core matrix. Fix once for this \(H\) a
nonzero form \(\Lambda\) on \(A\) with kernel \(H\).

\begin{proposition}[The singular positions in the image hyperplane]
\label{prop:markedimage393}
Every position satisfying the entry congruences belongs to
\(A\), for both the plain and symmetrized tests. For a varying
determinant the cell is a coset of \(W\) in \(A\), and its
singular positions are precisely \(D_0\cap H\). These have the
form \(u+K_0\), where \(u\in H\). The class of \(u\) in
\(H/K_0\) is uniquely determined by the coset \(D_0\).
On that cell the position function, including its splitting
factor, equals \(\gamma\nu_i(\Lambda(v))\) for a root of unity
\(\gamma\). The multiplicative character is extended by zero.
Among varying paths, \(H\) also determines \(i\).
\end{proposition}
\begin{proof}
There are exactly \(a\) columns in the tested matrix
\([\,\bar X\mid v\,]\bar U\) and \(\bar U\) is invertible.
Both entry tests prescribe zero in all these columns at every
row of index at least \(a\), using zero-based indices.
Right multiplication by \(\bar U^{-1}\) shows that the same
rows of \([\,\bar X\mid v\,]\) vanish. In particular \(v\in A\).
This is an equation on the original test before any Fourier
transform.

The entry congruences in the varying case leave a coset of
\(W\) by Proposition~\ref{prop:markedcell386}. On the top
square the determinant is both the product of the diagonal
block determinants and
\(\det(\bar U)\det[\,\bar X^{(a)}\mid v_{<a}\,]\).
All the other diagonal block determinants are fixed nonzero
constants on this cell, including the unitriangular initial
block. The last displayed determinant is a nonzero homogeneous
linear form in \(v\) with kernel \(H\). Thus the varying
block determinant on the entire cell is a fixed nonzero
multiple of \(\Lambda(v)\). Its zeros are \(D_0\cap H\).
The additive superdiagonal factor is constant on a varying
cell, as in Lemma~\ref{lem:markedsingular387}. Combining
it with the characters of the fixed determinants, of the
nonzero multiplier and of the splitting unit gives \(\gamma\).

The restriction of \(\Lambda\) to \(W\) is nonzero, so
\(H+W=A\). Inclusion induces an isomorphism
\(H/K_0\longrightarrow A/W\). A coset \(D_0\) of \(W\)
therefore determines exactly one class \(u+K_0\) in \(H\),
and \(D_0=u+W\). Finally
Proposition~\ref{prop:markedcell386} makes \(\Lambda|_W\)
nonzero on block \(B_i\) and zero on all the preceding
coordinates. The first diagonal block on which \(\Lambda\)
is nonzero is consequently \(B_i\), which recovers \(i\)
from the hyperplane.
\end{proof}

\begin{theorem}[The complete reduced Fourier table at a fixed image]
\label{thm:markedimage393}
Group varying paths having this same \(H\), with fixed frame
and orientation at the marked place. For \(u\in H/K_0\)
let \(c(u)\) be the signed sum of the incoming coefficients
whose position cell is \(u+W\), including the factor \(\gamma\).
It may take values in the intermediate coefficient space.
The table before the terminal coefficient operator depends
on \(z\in V^\vee\) only through its restriction \(z_A\) to \(A\).
Its potentially nonzero restrictions are indexed bijectively by
\[
 (t,\xi)\in k^\times\times (H/K_0)^\vee,\qquad
 z_A|_W=t\Lambda|_W,\qquad z_A|_H=\xi .
\]
Here \(\xi\) on \(H\) denotes its pullback from \(H/K_0\).
In these coordinates the exact table is
\[
 F(t,\xi)=q^{s-1}G_{-1}(\nu_i)\nu_i(t)^{-1}\widehat c(\xi),
 \qquad
 \widehat c(\xi)=\sum_{u\in H/K_0} c(u)\theta(-\xi(u)).
\]
It vanishes on the other restrictions. In any Hermitian norm
on the coefficient space,
\[
       \sum_{z_A\in A^\vee}\|F(z_A)\|^2
              =q^{a+s-1}(q-1)\sum_{u\in H/K_0}\|c(u)\|^2 .
\]
In particular the original table vanishes identically if and
only if every \(c(u)\) vanishes. For fixed \(H\) the space
of scalar tables obtained by allowing arbitrary such
coefficients has dimension \(q^{a-s}\).
\end{theorem}
\begin{proof}
All the position functions are supported in \(A\), hence
their Fourier transforms are unchanged on adding an element
of \(A^\circ\) to \(z\). On a fixed cell write \(v=u+y\)
with \(u\in H\) and \(y\in W\). Then
\(\Lambda(u+y)=\Lambda(y)\). Summing
\(\nu_i(\Lambda(y))\theta(-z_A y)\) over the fibers of
\(\Lambda|_W\) gives zero unless \(z_A|_W=t\Lambda|_W\).
For \(t=0\) it is zero by multiplicative orthogonality of
the nontrivial marked character. For \(t\ne0\) it equals
\(q^{s-1}\nu_i(t)^{-1}G_{-1}(\nu_i)\), by substitution in
the one-dimensional Gauss sum. The remaining factor in this
cell is \(\theta(-z_Au)\). A form on \(A=H+W\) is determined
by its restrictions to \(H\) and \(W\) provided they agree
on their intersection. On \(K_0\) the indicated restrictions
are both zero, so each displayed pair gives exactly one \(z_A\).
Summing the cell coefficients proves the formula.

The Gauss sum satisfies \(|G_{-1}(\nu_i)|^2=q\), by
Theorem~\ref{thm:markedgauss386}. Summing over \(t\) supplies
the factor \(q-1\). Additive character orthogonality on
\(H/K_0\) gives
\(\sum_\xi\|\widehat c(\xi)\|^2=q^{a-s}\sum_u\|c(u)\|^2\).
Multiplication gives the asserted mean square. For a sum on
all \(z\in V^\vee\) the right side is multiplied by
\(q^{d-a}\). Injectivity and the dimension statement follow
from the norm identity. Actual matrix paths can use a subset
of the cells in the displayed coefficient space.
\end{proof}

\begin{corollary}[The terminal test on the last varying block]
\label{cor:lastblock393}
For each \(\xi\in(H/K_0)^\vee\) define the terminal operator
\[
 L(\xi)=\sum_{t\in k^\times}\nu_i(t)^{-1}
                \sum_{z:\,z|_A=z_A(t,\xi)}J(z).
\]
At this marked place the signed contraction of the preceding
theorem is exactly
\[
          q^{s-1}G_{-1}(\nu_i)\sum_{\xi} L(\xi)\widehat c(\xi).
\]
The inner sum defining \(L(\xi)\) replaces the fine terminal
tests by their union in the normal position model. At several
marked places these operations can be performed simultaneously
with the joint coefficient table and the original global
terminal operator.

When \(i=m-3\), the last block of the penultimate test,
\(W=A\), \(K_0=H\), and there is a single \(\xi\). Apart
from the common Gauss scalar the marked test on an
integral normal row \(y\) has value
\(\nu_i(t)^{-1}\) when \(\bar y|_A=t\Lambda\) with
\(t\ne0\), and zero otherwise. The remaining \(d-a\)
residue coordinates of \(y\) are unrestricted. If every
varying marked place is in its last block, a fixed tuple of
image hyperplanes therefore gives a single prescribed terminal
operator applied to the sum of its incoming coefficients.
\end{corollary}
\begin{proof}
Insert the Fourier formula into the original sum over \(z\),
and first collect all extensions of the same \(z_A\).
Their coefficients are identical. In the normal theta
realization their residue characteristic functions are disjoint
and add to the characteristic function for the restriction
to \(A\). Thus their operators are summed with their given
normalization, by Proposition~\ref{prop:normalcells388}.
Then collect the \(t\)'s with the same \(\xi\). This gives
the asserted expression, including the inverse multiplicative
character and the Gauss factor. At several places the finite
sums commute. Equivalently one can insert the product of the
local characteristic functions in the single global theta
section before differentiating. Neither the terminal operator
nor the incoming table has to factor over the places.

When \(i=m-3\), the endpoint \(s_i\) of the block is
\(r+(m-2)t_{\rm r}=a\), so the sole nonempty coset of
\(W\) inside \(A\) is \(A\) itself. Its singular positions
are exactly \(H\), through the origin. The quotient \(H/K_0\)
is zero and the formula gives the stated unit-weighted test
on its support. In particular the many extensions counted
in Proposition~\ref{prop:initialrank392} all enter with the
same incoming coefficient for fixed \(t\) here. The
nonvanishing of this concrete operator on the arithmetic
input after the earlier pullbacks, and the effect of the
other image hyperplanes and blocks, still enter the global
period in Section~\ref{sec:missing}.
\end{proof}


\subsection{The marked tests at the first intermediate vertex}
\label{sec:markedline394}

At vertex \(1\) the terminal prime provides about \(T^{fd/m}\)
distinct pullbacks. We compute how much of this supply is needed
by the actual test at the other two primes. Keep the common
interval, metric, frames and normal sections of
Theorem~\ref{thm:intermediateimage392}. In particular \(d=N(m-1)\),
\(D=fd\), \(b=\exp(x^\varepsilon)\), and \(Q_g=(g_1g_2)^D\).
Use \(w\) for a chosen half at \(g_1\) or \(g_2\), \(v\) for
a chosen half at \(T\), and \(k_w\) for the field of \(q_w\)
elements at \(w\). Write \(V_w\) for the marked space of dimension
\(d\), so that the integral normal position lattice \(U_w\)
reduces onto \(V_w^\vee\). For \(B\subset V_w\), let
\(S_w(B)\) denote the inverse image of its annihilator in \(U_w\).
For a full local position lattice \(S\), let \(S^\#\) be its
frequency dual for the local additive pairing of conductor the
integers. In particular \(U_w^\#\) is at the conjugate half.

\begin{proposition}[Residue sections on an affine subspace]
\label{prop:markedline394}
At each marked half choose subspaces \(K_w\subset B_w\subset V_w\),
an integral row \(z_w^0\) and a character \(\omega_w\) of the
residue row group. Put
\[
 S_w=S_w(K_w),\qquad S'_w=S_w(B_w),\qquad
 Z'=\prod_w S_w/S'_w,\qquad
 Q=|Z'|=\prod_w q_w^{\dim B_w-\dim K_w}.
\]
For a fine residue row \(z_w\) use any integral lift. On
\(z_w^0+S_w\) let \(\xi_w\) be the class of \(z_w-z_w^0\)
in \(S_w/S'_w\). Write \(\omega(z)=\prod_w\omega_w(z_w)\).
A marked input of the following form is determined by one
coefficient vector \(Y_\xi\) per \(\xi\in Z'\) by the formula
\[
 Y(z)=
 \begin{cases}
   \omega(z)Y_\xi,&z_w\in z_w^0+S_w\text{ for every }w,\\
   0,&\text{otherwise}.
 \end{cases}
\]
For a tuple of intermediate images \(I_v\), sum the normal sections
with coefficients \(\omega(z)\) over the fine positions giving
a fixed \(\xi\) and a fixed terminal pullback \(u\).
Including zero pullbacks gives a basis of a positive theta line
of dimension \(Q\prod_v T^{\dim I_v}\). At the marked halves its
common periods are \(S'_w\) in position and \(S_w^\#\) in
frequency. The squared norm of the sum of its basis sections is
\[
          s_K^2=s^2\prod_w q_w^{-\dim K_w},
          \qquad s=\Big\|\sum_{z,t}\Theta_{z,t}\Big\|_{\rm av}.
\]
Fixing one terminal label at \(v\) divides this squared norm
by \(T^{\dim I_v}\).
\end{proposition}
\begin{proof}
The marked function for one new position is \(\omega_w\)
times the indicator of one coset of \(S'_w\) in \(z_w^0+S_w\).
Addition of the fine indicators gives this identity with
coefficient one. Position translations by \(S'_w\) act on
all the functions by the same scalar from \(\omega_w\).
Frequency modulations by \(S_w^\#\) also act by one common
scalar, given by their value at \(z_w^0\). The analogous
calculation at \(v\) is Proposition~\ref{prop:coherentcell391}.
This proves the period law. For the dimension, start at \(w\)
with the position and frequency pair \(S_w,S_w^\#\).
They are exact duals, so the index gained on one side from the
integral vacuum is lost on the other. Replacing \(S_w\) by
\(S'_w\) multiplies the real covolume by \([S_w:S'_w]\).
The terminal restrictions multiply it by \(\prod_vT^{\dim I_v}\).
The frozen data are as in Proposition~\ref{prop:normalcells388}.
Riemann--Roch for the resulting positive line gives the claimed
dimension.

The sections are orthogonal already on the fine cover, because
distinct new cells have disjoint sets of indices in the original
orthogonal basis. Integration of their periodic metric products
against probability Haar measure gives the same inner product
on the common cover and on the new torus. The norm computation
can consequently be made in the fine basis. Its sections have
equal squared norm. Exactly a fraction
\(\prod_w q_w^{-\dim K_w}\) of the fine positions occurs in the
sum, each with a multiplier of absolute value one. Imposing one
terminal value fixes a fraction \(T^{-\dim I_v}\) of those cells
since pullback onto its image has equal fibers. This proves both
norm formulas. Every new cell is nonempty and has the same number
of fine terms, which proves that the sections give the asserted
basis.
\end{proof}

\begin{lemma}[Short periods after retaining the marked quotient]
\label{lem:markedshort394}
Choose independently uniform \(a_1\)-planes \(I_v\) at the
terminal halves, where \(a_1=N(1)\). On the full table of
the proposition, take the group \(\mathcal M_I\) consisting
of the common periods and all their translates representing
translations and modulations on \(Z'\) that act trivially on
the terminal labels. Make the same construction on each of
the \(f\) faces obtained by fixing one terminal label.
Suppose that \(\log Q\le(\kappa+o(1))D\log b\), where
\(0\le\kappa<16\). Fix \(\sigma>0\) with \(\kappa+\sigma<16\).
The proportion of image tuples for which any of these groups
has a nonzero vector of squared length less than \(b^{1+\sigma}\)
is at most
\[
             \exp\{-(16-\kappa-\sigma-o(1))D\log b\}.
\]
One may test all the transported complex realizations at once.
The exceptional set for fixed \(K_w,B_w\) is independent of
\(z_w^0,\omega_w\).
\end{lemma}
\begin{proof}
Leave the terminal positions unrestricted for a moment. Write
\(P=\delta\Phi(\mathfrak b^{-1}M)\) for the lattice whose
marked pairs are \(S_w,S_w^\#\), using the period coordinates of
Theorem~\ref{thm:lattice55}. Here \(M\) is a self-dual Hermitian
\(\mathcal O_E\)-lattice. Indeed one obtains its local factors
from the original self-dual normal module by replacing a
position lattice and its conjugate dual by the displayed pair.
The primes here are split and distinct from the different
at the marked places. The fixed norm-one modifications away
from them preserve the duality as well. The global module is
the intersection of these local conditions in the rational
space. It has precisely those localizations by the Chinese
remainder theorem, since only finitely many conditions change.
Duality may be checked after localization, proving the assertion.

At \(w\) the translations of the new basis are represented by
\(S_w/S'_w\). The modulations distinguishing its cells are
\((S'_w)^\#/S_w^\#\). In particular their lifts, together with
the common lattice, have local pairs
\[
                  S_w,\ (S'_w)^\# .
\]
The quotient by the common periods at these places has order
\(Q^2\). The common unit scalars in the preceding proof can
be removed on each operator. The residue pairing is perfect,
so distinct translation and modulation pairs have distinct
cosets of the period lattice. The Chinese remainder theorem
produces lifts acting trivially on the terminal labels. Adding
their lifts to the common periods gives the displayed local
pairs and preserves the period conditions at the other places.
This proves in particular that \(\mathcal M_I\) is a lattice.
Affinely shifting a test or changing its character changes
only the phases in this calculation.

Let \(e=0\) for the full table and \(e=1\) on a face at \(v_0\).
Take \(P_e=P\) for \(e=0\). For \(e=1\) multiply its fractional
module by the ideal \(\bar v_0^{-1}\).
At \(v_0\) the dual of the terminal position kernel is contained
in \(T^{-1}U_{v_0}^\#\), which justifies this enlargement.
Ignore all the remaining terminal restrictions and enlarge
the marked pairs of \(P_e\) to \(S_w,(S'_w)^\#\). Call the
result \(P_e^+\). It contains \(\mathcal M_I\), is independent
of \(I\), and is a union of exactly \(Q\) cosets of \(P_e\).
This last index is \(\prod_w[(S'_w)^\#:S_w^\#]\).
Also \(P_e^+\subset G^{-1}P_e\), where \(G=g_1g_2\),
because \(q_w S_w\subset S'_w\) at every marked place.

We use the trace lower bound with this modified self-dual module.
For any fractional ideal \(\mathfrak j\) a nonzero vector of
\(\delta\Phi(\mathfrak b^{-1}\mathfrak j M)\) has squared norm
at least \(2f(\operatorname N\mathfrak j)^{1/f}/\operatorname{rd}(E)\).
Indeed its Hermitian square belongs to
\(\mathfrak j\overline{\mathfrak j}\mathfrak D_{E/K}^{-1}\)
by local integrality. Taking norms and applying the arithmetic
and geometric mean inequality at the \(f\) embeddings gives
the bound with the same metric factor as in
Theorem~\ref{thm:lattice55}. Thus distinct vectors in a given
coset of \(P_e\) have distance at least \(R_e\), where
\[
                R_e^2=2fT^{-e/f}/\operatorname{rd}(E).
\]
Set \(R^2=b^{1+\sigma}\). Disjoint balls of radius \(R_e/2\)
about vectors of norm less than \(R\) in one coset fit inside
the ball of radius \(R+R_e/2\). There are at most
\(Q(1+2R/R_e)^{2D}\) vectors in \(P_e^+\) of norm less than
\(R\). Its logarithm is at most
\((\kappa+1+\sigma+o(1))D\log b\), since
\(\log\operatorname{rd}(E)=o(\log b)\) and
\((\log T)/f=o(\log b)\) on the common intervals.

Consider a nonzero vector counted here. Let \(h\) be the
number of unfrozen terminal halves where its position has
zero residue. It then lies in \(G^{-1}P_e\) with the additional
ideal factor given by the product of those \(h\) primes.
This uses complete splitting, so a zero position is exactly
divisibility of the local module by the prime on that half.
The trace bound gives squared norm at least
\(2fG^{-2}T^{(h-e)/f}/\operatorname{rd}(E)\).
Since \(\log G=(16+o(1))\log b\) and
\(\log T=(17+o(1))m\log b\), we have \(h/f=O(1/m)\).
At every other unfrozen half a nonzero residue must annihilate
\(I_v\). The probability of this is
\((T^{d-a_1}-1)/(T^d-1)\le T^{-a_1}\), by counting planes
in a hyperplane as in Lemma~\ref{lem:shortimage392}.
Independence gives a probability at most \(T^{-a_1(f-e-h)}\).
The identity \(a_1=k+t_{\rm r}+1\) and
\(d=k+(m-1)(t_{\rm r}+1)\) gives \(a_1m\ge d\).
The logarithm of the probability is therefore at most
\(-(17-o(1))D\log b\). Multiply by the count of ambient vectors
and sum over the \(f+1\) cases to obtain the bound.

Under transport of the CM data the lattice and module indices
undergo the embedding of \(E\) and the common elliptic ideal
change in the proof of Lemma~\ref{lem:shortimage392}. Unit
scalars preserve all the subspaces and their annihilators,
and changing an ideal representative is an isometry in the
transported metric. The at most \(2f h_K\) resulting cases
contribute \(\log(2f h_K)=o(D\log b)\) to the logarithmic
estimate by that lemma. This completes the count.
\end{proof}

\begin{theorem}[A signed mean square at vertex one]
\label{thm:firstvertex394}
Use the matrices from vertex \(1\) through vertex \(m-1\),
with the uniform counting measure on their injective reductions
at the terminal places and the fixed data of
Theorem~\ref{thm:intermediateimage392}. Let \(C_v\) be their
composite including the prescribed unit. For subspaces as in
Lemma~\ref{lem:markedshort394}, outside the exceptional fraction
there, every marked input of Proposition~\ref{prop:markedline394}
has the following bounds. Define
\[
 D_u(Y)=\sum_{\substack{z,t,\alpha\\ t_v C_v=u_v\text{ for all }v}}
              \ell_\alpha(\Theta_{z,t})A_\alpha Y(z),\qquad
 M(Y)=\frac{s^2}{Q_g}\sum_{z,\alpha}\|A_\alpha Y(z)\|^2 .
\]
Then for \(\rho_x=e^{-x^{\varepsilon/8}}\),
\[
 (1-\rho_x)M(Y)\le \sum_{u_v\ne0\ {\rm for\ all}\ v}
              \|D_u(Y)\|^2\le(1+\rho_x)M(Y).
\]
The estimate holds simultaneously for every vector in the
stated input space and for all the affine centers and characters.
\end{theorem}
\begin{proof}
The image of a uniform composite of injections is a uniform
\(a_1\)-plane by invariance under left multiplication on the
last target. The different places are independent by the
matrix residue calculation of Theorem~\ref{thm:intermediateimage392}.
For a good image take the partial derivative Gram matrix of the
basis in Proposition~\ref{prop:markedline394}, retaining the
index \(\xi\) and summing over \(u\). Expand it in the \(Q^2\)
translation and modulation matrices on \(Z'\) as in the proof of
Lemma~\ref{lem:partialtheta390}. Its identity contribution,
after multiplication by \(Q/s_K^2\), is the identity on the
retained cells and normalized derivative indices. The other
terms are compressed Fock translations at nonzero vectors in
\(\mathcal M_I\), one coset for each matrix. Their phases have
absolute value one. That proof and the shell estimate of
Lemma~\ref{lem:thetaframe389} use only distinct cosets in a
lattice whose nonzero vectors have length at least \(R\).
Apply them to \(\mathcal M_I\) itself. They give operator
error \(\delta=2e^{-\pi R^2/8}\), since \(R^2=b^{1+\sigma}\)
dominates \(D\) and the degree cutoff \(L\ll x^C b\).
The identical argument on a zero face, with its own lattice,
gives the upper bound \((1+\delta)s_K^2/(QT^{a_1})\).
Removing all zero faces therefore gives relative error at most
\(\delta+(1+\delta)fT^{-a_1}\le\rho_x\).

Apply the two bounds to the scalar coordinates of
\(A_\alpha Y_\xi\) and sum in an orthonormal output basis.
Linearity and the defining marked character put the original
signed coefficient inside the displayed norm. Its reference
norm is \((s_K^2/Q)\sum_{\xi,\alpha}\|A_\alpha Y_\xi\|^2\).
Indeed each \(\xi\) represents \(\prod_wq_w^{d-\dim B_w}\)
fine positions. Substituting the norm formula of the proposition
shows that this is exactly \(M(Y)\).
\end{proof}

\begin{corollary}[Individual determinant tests and signed sums of later blocks]
\label{cor:markedpaths394}
Fix the marked frames. Outside a proportion at most
\(\exp(-14D\log b)\) of the injective terminal matrix residues,
the estimate of Theorem~\ref{thm:firstvertex394} holds for each
individual nonzero penultimate marked test, multiplied by any
incoming vector independent of \(z\). For such an input
\(Y(z)=\prod_w\mathcal F_w(z_w)y\) its reference norm is
\[
     s^2\left(\prod_w n_w\right)\sum_\alpha\|A_\alpha y\|^2,
 \quad
 n_w=\begin{cases}
       q_w^{s_w-1}(q_w-1),&\text{on a varying cell with }
                                    s_w=\dim W_w,\\
       q_w^{\dim W_{c,w}},&\text{on a constant determinant cell}.
     \end{cases}
\]
Here \(\mathcal F_w\) includes the root of unity in
Lemma~\ref{lem:markedsingular387}. Each additional common
normalization multiplies this norm by its squared modulus.

There is also a bound after adding paths. Fix \(0<\vartheta<1\)
and set \(p=\lfloor\vartheta d\rfloor\). At every marked
half keep the varying blocks whose first row index is at least
\(p\), with row indices starting at zero in the marked frame.
Use a common terminal operator and the same pullback \(C\).
For a signed sum of these tests with arbitrary incoming vector
coefficients the theorem holds outside a proportion at most
\(\exp(-7\vartheta D\log b)\). In this case \(M(Y)\)
uses the squared norm of the signed sum at each frequency.
\end{corollary}
\begin{proof}
For a varying block take \(B_w=W=A_{s_w}\) and
\(K_w=\ker(\lambda|_W)\) in the notation of
Lemma~\ref{lem:markedsingular387}. Its Fourier value is zero
unless \(z_w|_W=t\lambda\) with \(t\ne0\). On this set it is
a scalar times \(\nu_i(t)^{-1}\theta(-z_wu)\).
Thus \(Z'_w\) has dimension one, with coordinate \(t\), and
\(\theta(-z_wu)\) is precisely an allowed \(\omega_w\).
Take \(Y_0=0\) at the omitted scalar value at this place.
For a constant determinant use \(B_w=K_w=W_c\),
\(z_w^0|_{W_c}=\rho_c\), and the same position character.
There is just one coarse position. This proves that
\(Q\le G^f\) for each individual test. Finite Fourier
orthogonality gives
\(q_w^{-d}\sum_{z_w}|\mathcal F_w(z_w)|^2
=\sum_v|h_w(v)|^2=n_w\), counting the positions on which the
original function has modulus one. Multiplication gives the
stated normalization.

For simultaneous control, \(W_c\) is a coordinate subspace of
the initial \(r\)-space by Proposition~\ref{prop:markedcell386}.
There are at most \(2^r\) choices. For a varying test the form
on \(A_{s_w}\) is supported on its one determinant block of
\(t_{\rm r}\) rows, so there are at most \(m q_w^{t_{\rm r}}\)
choices for the pair \(K_w,B_w\). Multiplying the sum of
these two bounds over the marked halves has logarithm
\(O(fr+f\log m+ft_{\rm r}\log G)=o(D\log b)\).
The centers and position characters cause no further loss by
the lemma. Apply it with \(\kappa=0,\sigma=1/2\) and take
the union to obtain the first exceptional bound.

For the signed sum take \(K_w=A_p\) and \(B_w=A_{d-t_{\rm r}}\).
The specified collection is empty if \(p>d-t_{\rm r}\).
In the other case these are nested subspaces. Each active form in
the selected blocks vanishes on \(A_p\). Its Fourier table
is supported on \(A_p^\circ\), and all its values, including
the phase from the singular coset, depend only on restriction
to \(A_{d-t_{\rm r}}\) by Proposition~\ref{prop:markedimage393}.
These two properties pass to the joint signed sum, even when
the coefficients do not factor across the marked places.
It is an input in the proposition, with centers zero and
characters one. Now \(\log Q\le(16(1-\vartheta)+o(1))D\log b\).
Use \(\kappa=16(1-\vartheta)\) and
\(\sigma=\min(8\vartheta,1/2)\) in the lemma. Since
\(16\vartheta-\sigma\ge8\vartheta\), its bound implies the
claimed estimate for all sufficiently large \(x\).
\end{proof}

The outer sum in the theorem ranges over covectors at vertex \(1\).
The first edge to vertex \(0\) further identifies these covectors,
and paths with different composites \(C\) can contribute to the
same initial coefficient. The exceptional proportion counts matrix
residues with equal weight. Section~\ref{sec:missing} requires the
integral norm and the archimedean comparison after those additional
signed contributions and their arithmetic weights are included.


\subsection{Initial covectors and the weighted first theta edge}
\label{sec:firstresidue395}

Write \(n=N(1)\), \(k=N(0)\), and \(q=T\). At one half over
the terminal prime the first residue matrix is a map
\(X:U\to V\) over \(\mathbb F_q\), with \(\dim U=k<n=\dim V\).
A later part of the theta chain can prescribe a different value
for each nonzero covector \(u\) on \(V\). For a table \(F(u)\)
of vectors in a complex Hilbert space set
\[
            (B_X F)(\eta)=\sum_{uX=\eta}F(u)
                    \qquad(0\ne\eta\in U^\vee).
\]
Both covector tables use counting measure. We first compute the
second moment when each injection \(X\) has equal weight. We then
compute the sum with the coefficients of the actual first edge.

\begin{proposition}[The mean square of the first pullback]
\label{prop:firstincidence395}
Let \(X\) be a uniform injection \(U\to V\). Put
\[
 h_1=\frac{q^{n-k}-1}{q^n-1},\qquad
 h_2=\frac{(q^{n-k}-1)(q^{n-k}-q)}
               {(q^n-1)(q^n-q)} .
\]
On scalar tables the operator \(H=\mathbb E_X B_X^*B_X\) has
eigenvalues on the three mutually orthogonal spaces given by
\begin{align*}
 1-h_1
       &\quad\text{on tables with sum zero on each covector line},\\
 1-h_1-(q-1)(h_1-h_2)
       &\quad\text{on tables constant on each line and of total sum zero},\\
 (1-h_1)q^{n-k}
       &\quad\text{on the constant table}.
\end{align*}
Here a covector line in a table contains its nonzero elements.
In particular \(H\succeq(1-q^{1-k})I\). For vector tables
these operators are tensored with the identity on the values.
\end{proposition}
\begin{proof}
The annihilator of \(XU\) is a uniform \((n-k)\)-plane in
\(V^\vee\). It contains one prescribed nonzero form with
probability \(h_1\), and two prescribed independent forms
with probability \(h_2\). Indeed, count ordered nonzero forms
or ordered independent pairs in that plane and divide by the
corresponding count in \(V^\vee\). The entry \(H_{u,u'}\)
is the probability that \(uX=u'X\ne0\). It equals \(1-h_1\)
on the diagonal. Distinct multiples of one form cannot have
the same nonzero image. For independent \(u,u'\), the
annihilator must contain \(u-u'\) and must exclude \(u\).
These forms are independent, giving probability \(h_1-h_2\).

Thus the off-diagonal part is \(h_1-h_2\) times the all-ones
matrix with every diagonal block of one line removed. It
vanishes on the separately zero-sum tables and acts as
\(-(q-1)(h_1-h_2)\) on the second subspace. Constants form
the remaining eigenspace. For each injection every nonzero
\(\eta\) has \(q^{n-k}\) preimages. The squared norm
of \(B_X\mathbf1\) divided by \(q^n-1\) is
\((q^k-1)q^{2(n-k)}/(q^n-1)=(1-h_1)q^{n-k}\).
This gives the last eigenvalue. All three are at least
\(1-qh_1\ge1-q^{1-k}\).
\end{proof}

\begin{corollary}[A second moment at the original initial vertex]
\label{cor:firstaverage395}
Fix a suffix and marked input satisfying
Theorem~\ref{thm:firstvertex394} at a good tuple of suffix
images. Put \(F(u)=D_u(Y)\) on tuples of nonzero covectors at
vertex \(1\). Keep this table fixed and let the first injections
\(X_v\) vary independently at the \(f\) terminal halves.
Write \(B_X=\bigotimes_v B_{X_v}\). Then
\[
 \mathbb E_X\sum_{\eta_v\ne0\ {\rm for\ all}\ v}
       \|(B_XF)(\eta)\|^2
   \ \ge\ (1-T^{1-k})^f(1-\rho_x)M(Y).
\]
On either table let \(\Pi=\prod_v(1-A_v)\), where \(A_v\)
averages nonzero scalar multiples in position \(v\). The exact
mean square on the separately zero-sum component is
\[
            \mathbb E_X\|\Pi B_XF\|^2
                     =(1-h_1)^f\|\Pi F\|^2 .
\]
The mean of \(\Pi B_XF\) itself is zero.
\end{corollary}
\begin{proof}
Independence makes the product Gram operator \(H^{\otimes f}\).
The lower bound of the proposition and the lower estimate
for \(\|F\|^2\) at vertex \(1\) prove the first assertion.
Each pullback commutes with scalar multiplication of the
covectors at a place, so \(\Pi B_XF=B_X\Pi F\).
The product Gram acts on the range of this projection by
\((1-h_1)^f\). For the last assertion, independently compose
each \(X_v\) on the right with an invertible endomorphism of
the initial space. This preserves its uniform distribution
and permutes all nonzero initial covectors transitively.
The mean output is consequently constant on their product.
Applying \(\Pi\) gives zero.
\end{proof}

\begin{theorem}[Rank-one frequencies of the actual first edge]
\label{thm:firstslice395}
Fix initial and first intermediate unrefined classes and their
frames. Let \(E(X):D_0\to D_1\) be the sum of the first-edge
coefficient maps with residue matrix tuple \(X=(X_v)_v\)
at \(T\). Include their prescribed finite weights away from
\(T\). Use all the integral residue matrices. Let \(F(u)\)
be the map on \(D_1\) given by the remaining edges at the
tuple \(u\) of nonzero intermediate covectors. Keep all
classes, finite signs and the coupling of the last two
edges in this sum. Assign each unrefined stabilizer mass
once, as in Theorem~\ref{thm:covectorpaths352}. For the
physical nonzero initial covectors the coefficient map is
\[
           P(\eta)=\sum_{X,u:\,u_vX_v=\eta_v\ {\rm for\ all}\ v}
                         F(u)E(X).
\]
For a tuple \(v=(v_P)_{P\mid\pi_T}\) of initial vectors define
\[
 \widehat E(v\otimes u)
   =\sum_X \left(\prod_{P\mid\pi_T}
             \psi_P\big(-u_P(X_Pv_P)\big)\right)E(X),\qquad
 K(v)=\sum_u F(u)\widehat E(v\otimes u).
\]
The transformed first-edge operator \(\widehat E(v\otimes u)\)
uses the original integral matrix-position test at each
terminal half multiplied by the indicated additive character.
Equivalently it is a frequency translate of that integral
theta test by the lift of \(-v_P\otimes u_P\) divided by \(T\),
in the dual coordinates for \(\psi_P\). The arithmetic
coefficient maps and derivatives have their original normalization.

For any compatible Hilbert norms, with Hilbert--Schmidt norm on
the coefficient maps, the actual signed row on ordinary initial
tables satisfies
\[
    \sum_{\eta}\|(\Pi P)(\eta)\|_{\rm HS}^2
         =T^{-fk}\sum_{v_P\ne0\ {\rm for\ all}\ P}
                                   \|(\Pi K)(v)\|_{\rm HS}^2 .
\]
On the left \(\eta\) ranges over tuples with no zero component.
For several first intermediate classes, the corresponding \(K(v)\)
are added before taking the squared norm in this formula.
\end{theorem}
\begin{proof}
The initial terminal test checks integrality and the equation
\(u_P X_P=\eta_P\). In the finite quotient of matrix positions
the denominator is prime to \(T\) and the sublattice in the
denominator of the quotient is contained in the kernel of
reduction modulo \(T\). Thus one may group all first-edge
terms by \(X\) prior to imposing the covectors. The other
first-edge local factors are independent of these terminal
refinements at fixed frames. This is precisely the factorization
of the test in
\srcloc{Formalism/LocalToGlobal/Columns/FiniteSumNonzeroLargeDepth}{646}
and lines 1115--1120 there. For a local integral matrix
\(x\),
\[
 \sum_{X_P}\psi_P\big(-u_P(X_Pv_P)\big)
                       \mathbf1_{\bar x=X_P}
                   =\psi_P\big(-u_P(\bar x v_P)\big).
\]
This equality holds with both sides zero outside the integral
lattice if the right side is extended by zero. Summing higher
digits of a matrix cell and applying the linear theta
distribution therefore gives the claimed operator. In split
position coordinates a frequency translation acts by multiplication
by its pairing character, proving the equivalent description.
The remaining path can depend on \(u\) and on its unrefined
initial class at vertex \(1\) and uses the value of the preceding
edge in that class. This gives the displayed formula for \(P\)
by composition of the actual maps.

Extend \(P\) by the same fiber formula to include all zero
faces. Its unnormalized additive Fourier transform at \(v\)
is \(K(v)\), because substituting \(\eta=uX\) in the additive
character and interchanging finite sums gives exactly that
expression. Scalar multiplication on a covector corresponds
to inverse scaling on the dual vector under Fourier transform,
so scalar averaging, and hence \(\Pi\), commute with the
transform. On a face with a zero component the corresponding
factor \(1-A_P\) is zero. Thus the choice of the formal
extension on those faces has no effect on the projected identity.
Plancherel on the additive group of order \(T^{fk}\), applied
to each scalar matrix entry in orthonormal bases, proves the
formula. Class addition commutes with all these finite linear
operations. This also proves the assertion about the order
of summation.
\end{proof}

The first mean square distinguishes the initial pullbacks when
the first edge is sampled separately. The expression \(K(v)\)
retains the dependence of its theta coefficient on that edge.
In particular the signed first-edge sum required in
Section~\ref{sec:missing} samples frequencies with at most rank
one at each terminal half. Its size depends on how those
translated coefficients pair with the fixed suffix maps.


\subsection{Integral initial tests and the values of the signed column}
\label{sec:initialmoments396}

Fix one ordinary initial class \(a\), its rank \(k\) module and
its terminal residue spaces \(U_P\) at the \(f\) chosen halves.
Use the frames and coefficient lattice \(D_a\) of
Theorem~\ref{thm:pointlattice353}. Enlarge the coefficient
field \(J\) to contain the character values used below. At
every \(P\mid\pi_T\) choose a nontrivial multiplicative character
\(\chi_P:\mathbb F_T^\times\to\mathcal O_J^\times\), extended
by zero at zero. For a tuple \(v\) of nonzero initial vectors
and a coefficient vector \(y\), set
\[
     \phi_{v,y}(\eta)
         =\left(\prod_P\chi_P(\eta_P(v_P))\right)y
\]
on the covectors over \(a\), and give it zero values over
other classes. We compute the cost of these tests in the
ordinary column and then evaluate them on the original signed
period. Both calculations use the same \(E_x,T\) and word.

\begin{proposition}[Integral domain and complex norm of the character tests]
\label{prop:initialtest396}
For every \(v\) the exact domain of \(y\) giving an integral
initial test is \(D_*=z_{\max}D_a\) of
Theorem~\ref{thm:pointlattice353}. Each such test has zero
sum on each local scalar orbit, with the other positions fixed.
Choose one nonzero representative of each vector line and
each covector line at every \(P\). Put \(N=(T^k-1)/(T-1)\)
and let \(M_P\) be the \(N\)-by-\(N\) matrix with entry
\(\chi_P(\ell(v))\) in the covector row \(\ell\) and vector
column \(v\). Then
\[
        M_P^*M_P=T^{k-1}I,\qquad
        M^*M=T^{f(k-1)}I,\qquad M=\bigotimes_P M_P .
\]
The matrix \(M\) maps the line coefficients of a sum of the
tests to its values on the covector representatives. In the
Hilbert norm on all nonzero covectors, the squared norm acquires
the additional factor \((T-1)^f\) from the scalar orbits.
\end{proposition}
\begin{proof}
The stabilizer of the unrefined ordinary initial class at
the prescribed level is trivial by Proposition~\ref{prop:column110}.
At a tuple of lines the coefficient lattice is constant
throughout its scalar orbit. The scalar in the test is a
unit precisely when all \(\eta_P(v_P)\) are nonzero.
The intersection of the actual coefficient lattices over those
line tuples is \(z_{\max}D_a\), by the intersection calculation
in Theorem~\ref{thm:pointlattice353}. This proves both directions
of the domain assertion. Scaling at one place multiplies each
nonzero value by a nontrivial character, whose sum is zero.

At one place compute the inner product of two columns by
summing first over all nonzero forms and dividing by \(T-1\).
The summand is constant on each covector line, since the two
unit characters of the scalar cancel. For equal columns exactly
\((T-1)T^{k-1}\) forms have nonzero pairing. For distinct
vector lines the two prescribed vectors are independent.
Every pair of field values for their evaluations occurs
\(T^{k-2}\) times, and the resulting product of the two
character sums is zero. For \(k=1\) only the diagonal case
occurs. Tensoring the one-place identities gives the formula
for \(M\). Each covector representative has \((T-1)^f\)
scalar multiples of equal norm, giving the last assertion.
\end{proof}

\begin{theorem}[The integral index in the prescribed scalar characters]
\label{thm:initialindex396}
Let \(\mathcal L_\chi\) consist of the actual integral
initial tables supported over \(a\) satisfying
\(\phi(c\eta)=\prod_P\chi_P(c_P)\phi(\eta)\).
Let \(\mathcal J_\chi\) be the sublattice generated by
\(\phi_{v,y}\) for \(y\in D_*\) and for all vector representatives.
At every coefficient prime away from \(T\) these lattices coincide.

At a coefficient prime \(\mathfrak p\mid T\) normalize
\(v_\mathfrak p(T)=1\). For lattices \(J_0\subset L_0\)
of equal rank let \(\iota_\mathfrak p(L_0,J_0)\) denote the
valuation of the determinant of the inclusion in local bases.
Write \(r_D=\dim_J(D_a\otimes J)\) and
\(\chi_P=\omega^{j_P}\) in this completion, where \(\omega\)
is the Teichm\"uller character and \(1\le j_P\le T-2\).
For the actual highest weights \(\lambda_\tau\) one has
\[
\begin{split}
 \frac{\iota_\mathfrak p(\mathcal L_\chi,\mathcal J_\chi)}
                 {r_D N^f}
   =&\sum_{P\mid\pi_T}
       \left(\frac{k(N-1)}{2N}
                    -\left(1-\frac{k}{N}\right)\frac{j_P}{T-1}\right)\\
    &+\sum_{\tau:\,v_\mathfrak p(\tau(\bar\pi_T))=1}
       \left(\lambda_{\tau,1}
                      -\frac1k\sum_{i=1}^k\lambda_{\tau,i}\right).
\end{split}
\]
\end{theorem}
\begin{proof}
Specifying a value on each tuple \(L\) of covector lines
identifies \(\mathcal L_\chi\) with \(\bigoplus_L D_{a,L}\).
For each \(L\) one can choose \(v\) with all its pairings
nonzero, so \(D_*\subset D_{a,L}\) by the proposition.
Inside \(\bigoplus_L D_*\) the generating map is exactly
\(M\otimes I_{r_D}\). Away from \(T\) the three coefficient
lattices \(D_*,D_{a,L},D_a\) coincide, and \(M^{-1}=
T^{-f(k-1)}M^*\) is integral there. This proves the assertion
at all such primes, including those dividing \(T-1\).

At \(\mathfrak p\), consider the unnormalized additive Fourier
transform from covector functions \(g(c\ell)=\chi_P(c)^{-1}g(\ell)\)
to vector functions. In projective coordinates its matrix is
\(G_P M_P^t\), with
\(G_P=\sum_{s\ne0}\chi_P(s)^{-1}\psi_P(-s)\). Indeed scaling
the summation variable gives \(G_P\chi_P(\ell(v))\), also
when \(\ell(v)=0\) by character orthogonality. In the notation
of Theorem~\ref{thm:fouriersmith338}, scaling the basis
covectors gives these input functions eigencharacter
\(\omega^{j_P}\). Changing the sign or unit of the additive
character permutes the vector positions up to unit factors.
The theorem therefore gives the valuation of the determinant
of \(G_PM_P^t\) as \(k((N-1)/2+j_P/(T-1))\).
In rank one the same theorem gives
\(v_\mathfrak p(G_P)=j_P/(T-1)\), since there is one line
and its nonzero pairing is a unit. Subtracting \(N\) copies
gives the determinant of \(M_P\). The determinant of the
\(f\)-fold tensor has valuation \(N^{f-1}\) times the sum
of these one-place valuations. This accounts for the first
line of the index formula.

It remains to compute the inclusion of \(\bigoplus_L D_*\)
in the full line lattice. In the frame of
Theorem~\ref{thm:fourierslopes354} the exponents of
\(D_{a,L}\) relative to \(D_a\) are sums of one diagonal-coordinate
weight per affected embedding. The average of that coordinate
weight in \(W_{\lambda_\tau}\), counted with its multiplicity,
is \(k^{-1}\sum_i\lambda_{\tau,i}\). Indeed permutation
matrices preserve the multiset of weights and the sum of all
coordinates on every weight is the central weight
\(\sum_i\lambda_{\tau,i}\). On \(D_*\) the common scalar
has valuation \(\sum_{\tau\ {\rm affected}}\lambda_{\tau,1}\).
Subtracting the two determinant valuations and multiplying
by the number \(N^f\) of line tuples gives the second line
after normalization. This determinant calculation also applies
to fractional coefficient frames.
\end{proof}

\begin{theorem}[The full signed period modulo powers of the terminal prime]
\label{thm:initialmoment396}
Take \(y\in D_*\) and a fixed algebraic projected small test
\(F\). For each path \(\gamma\) from \(a\) let
\(w_\gamma=W_{\gamma,F}(y)\) be its coefficient in
Theorem~\ref{thm:covectorpaths352}, with the prescribed
masses and signs. Write \(\eta_{\gamma,P}\) for the composite
terminal covector at the initial vertex. Let \(R\) be the
completed valuation ring at \(\mathfrak p\mid T\), and
choose \(0\ne\Delta\) in its fraction field such that
\(\Delta w_\gamma\in R\) for every path in this finite sum.
For \(h\ge1\) set \(e_P=j_P T^{h-1}\). Use integral
\(\mathbb Z_T\)-lifts of the reduced edge matrices, the
initial vectors and the marked final covector including its
unit \(\bar\pi_T\bmod P\). Their composite row is denoted
\(\widetilde\eta_{\gamma,P}\). Then
\[
 \Delta\,\mathcal P(\phi_{v,y},F)
     \equiv
       \sum_\gamma (\Delta w_\gamma)
           \prod_{P\mid\pi_T}
              \big(\widetilde\eta_{\gamma,P}
                              (\widetilde v_P)\big)^{e_P}
                       \pmod {T^h R}.
\]
All choices of integral lifts give the same residue.
The right side can be formed by attaching at each edge matrix
\(x_{i,P}\) its ordinary symmetric power of degree \(e_P\)
in the lifted terminal coordinates, in tensor product over
\(P\). These maps are composed in the order of the theta chain.
The original last two edge tests retain their common coupling
variable throughout the multiplication.
\end{theorem}
\begin{proof}
Inserting this single admissible test into
Theorem~\ref{thm:covectorpaths352} gives the exact equality
\[
   \mathcal P(\phi_{v,y},F)
     =\sum_\gamma w_\gamma
               \prod_P\chi_P\big(\eta_{\gamma,P}(v_P)\big).
\]
Paths with a zero pullback can be included with value zero.
The terminal primes split completely, so after choosing local
frames each residue matrix may be lifted over \(\mathbb Z_T\).
For \(z\in\mathbb Z_T^\times\), write
\(z=\omega(\bar z)(1+Tb)\). The binomial theorem implies
\((1+Tb)^{T^r}\in1+T^{r+1}\mathbb Z_T\) by induction on \(r\).
Also \(T\equiv1\bmod T-1\). Hence
\(z^{j_PT^{h-1}}\equiv\chi_P(\bar z)\bmod T^h\).
For \(T\mid z\) this congruence holds with value zero because
\(j_PT^{h-1}\ge h\). Apply it to each scalar in the exact
equality. Integrality of the finitely many \(\Delta w_\gamma\)
permits addition of the congruences in \(R\), with no loss
of precision. This proves the formula and independence of lifts.

For a map \(X\) of free \(\mathbb Z_T\)-modules its induced
map on the symmetric algebra sends \(v^e\) to \((Xv)^e\).
The functional induced by a row \(a\) sends this last element
to \(a(Xv)^e\). Both assertions follow on monomials by ordinary
multiplication without division by factorials. Apply these
identities through all the matrices in a path, tensor over
the places, and take the terminal evaluation. This produces
the displayed product. Distributivity allows one to sum the
tensor maps over an edge while its shared class and coupling
indices stay fixed. The original coefficient maps compose
on their own factors. One can carry a single common denominator
outside the final sum, so this description does not require
individual edge maps to be integral in a new normalization.
\end{proof}

\begin{corollary}[A residue test on the prescribed signed matrices]
\label{cor:initialresidue396}
Let \(\kappa(\mathfrak p)\) be the residue field of \(R\).
In initial vector coordinates form the multihomogeneous polynomial
\[
     H((X_P)_P)=
       \sum_\gamma\overline{\Delta w_\gamma}
                     \prod_P\big(\eta_{\gamma,P}(X_P)\big)^{j_P}
                  \quad\text{over }\kappa(\mathfrak p).
\]
It is nonzero if and only if at least one of the integral tests
\(\phi_{v,y}\) above has normalized period
\(\Delta\mathcal P(\phi_{v,y},F)\) a unit at \(\mathfrak p\).
At \(\chi_P=\omega\) for every \(P\), its coefficient tensor is
\(\sum_\gamma\overline{\Delta w_\gamma}\bigotimes_P
\eta_{\gamma,P}\).
\end{corollary}
\begin{proof}
Take \(h=1\) and reduce the preceding congruence to the residue
field. This identifies \(H(v)\) with the required period residue.
The degree in each individual variable is at most
\(j_P\le T-2\). Evaluation of polynomials of degree less than
\(T\) per variable on the full product of copies of \(\mathbb F_T\)
is injective even after coefficient extension. In one variable
this follows from the root bound, and induction on the number
of variables gives the general assertion. If some vector
\(v_P\) is zero then the polynomial vanishes there since
\(j_P>0\). A nonzero polynomial therefore has nonzero value
at an allowed tuple. Scaling a representative scales that
value by a unit. Finally expansion of the product of linear
forms at \(j_P=1\) gives the stated tensor.
\end{proof}

The index theorem computes the cost of using these tests in the
actual initial coefficient lattice. The last two results retain
the entire signed sum of paths for each test. A bound in
Section~\ref{sec:missing} still requires an estimate for those
values with the prescribed arithmetic weights and period
normalization. The common denominator \(\Delta\) and a nonzero
coefficient of the residue polynomial have to be controlled
for that application.


\subsection{Integral moments of the composite terminal covectors}
\label{sec:teichmoments400}

Fix a field \(E_x\), a word and a translate on the interval of
Section~\ref{sec:missing}. Write \(p=T\) for its terminal prime,
\(d=p-1\), and \(R\) for the completed integer ring at a chosen
coefficient prime above \(p\). Normalize \(v(p)=1\). Write
\([a]\in\ZZ_p\) for the Teichm\"uller representative of
\(a\in\mathbb F_p\), with \([0]=0\). We use \(0^0=1\) in monomials.
The arithmetic weights, signs, projections and measures throughout
are those of the path formula in Theorem~\ref{thm:covectorpaths352}.

\begin{lemma}[The coefficients of a multiplicative character of a linear form]\leavevmode
\label{lem:teichrow400}
Take \(1\le j<d\) and a row \(a=(a_1,\ldots,a_k)\) over
\(\mathbb F_p\). Let \(\mathcal M_j\) consist of the multi-indices
\(0\le m_i\le d\) with \(|m|=\sum_i m_i\equiv j\pmod d\).
There are coefficients \(c_{j,m}\in\ZZ_p\), independent of the row,
such that for every \(v\in\mathbb F_p^k\) one has
\[
    [a(v)]^j=\sum_{m\in\mathcal M_j}c_{j,m}
          \prod_i[a_i]^{m_i}[v_i]^{m_i},\qquad
    v_p(c_{j,m})=\frac{|m|-j}{d}.
\]
Evaluation of the monomials \(\prod_i[v_i]^{m_i}\),
\(m\in\mathcal M_j\), on one representative per nonzero vector
line gives a square matrix invertible over \(\ZZ_p\).
\end{lemma}
\begin{proof}
In one variable, evaluation of degrees \(0,\dots,p-1\) on the
\(p\) Teichm\"uller representatives is an isomorphism over
\(\ZZ_p\). Its Vandermonde determinant is a product of differences
which are all nonzero modulo \(p\). Tensoring proves the same
assertion for \(k\) variables with these individual degree bounds.
Under scaling \(v\mapsto u v\), a monomial is multiplied by
\([u]^{|m|}\). Thus the monomials in \(\mathcal M_j\) are an
integral basis of functions with scaling character \([u]^j\).
Indeed distinct scalar characters differ by a unit at a generator
of \(\mathbb F_p^\times\). Such a function vanishes at zero, and
its values on the representatives are arbitrary. This proves
the invertibility assertion.

For the coefficients adjoin a primitive \(p\)th root \(\zeta\)
and put \(G_m=\sum_{u\ne0}[u]^{-m}\zeta^u\) for \(0<m<d\).
The polynomial of degree at most \(d\) giving \(\zeta^x\) on
\([x]\) has coefficients
\[
  b_0=1,\qquad b_m=G_m/d\quad(0<m<d),\qquad b_d=-p/d .
\]
Its value at zero determines \(b_0\). On the nonzero residues
multiplicative character orthogonality gives the interior
coefficients, while the mean equals \(-1/d=b_0+b_d\).
For \(u\ne0\), substitution of \([u][x]\) gives \(\zeta^{ux}\).
Finite Fourier inversion here is just the identity
\(\sum_{u\ne0}[u]^{-j}\zeta^{ut}=G_j[t]^j\), obtained by
changing the summation variable when \(t\ne0\) and by character
orthogonality when \(t=0\). Apply it to \(t=a(v)\). Expanding
\(\prod_i\zeta^{u a_i v_i}\) by the polynomials just computed
and summing over \(u\) gives precisely the stated monomials with
\[
                 c_{j,m}=\frac{d}{G_j}\prod_i b_{m_i} .
\]
The computation includes \(a_i=0\), since then the corresponding
polynomial factor has value one. These coefficients lie in
\(\ZZ_p\) by uniqueness of interpolation applied to the row with
all coordinates one and its \(\ZZ_p\)-valued function.

For completeness put \(\lambda=\zeta-1\). The identity
\(p=\prod_{u=1}^{p-1}(1-\zeta^u)\) gives \(v(\lambda)=1/d\),
since \((\zeta^u-1)/\lambda\) reduces to \(u\) modulo \(\lambda\).
Expand \(G_m\) using \(\zeta^u=(1+\lambda)^u\) with
\(1\le u\le p-1\). For \(n<m\) the coefficient of \(\lambda^n\)
is a multiple of \(p\) in \(\ZZ_p\). Indeed its reduction is
\(\sum_{u\ne0}u^{-m}\binom{u}{n}\), a sum of nontrivial
powers on \(\mathbb F_p^\times\). For \(n=m\) that coefficient
reduces to \(-1/m!\). Terms with \(n>m\) have strictly larger
valuation than \(m/d\), as do the terms with \(n<m\) because
\(m<d\). Thus \(v(G_m)=m/d\), in particular \(G_m\ne0\).
The exponents of \(b_0,b_d\) are zero and one. The formula for
\(c_{j,m}\) now gives the claimed valuation.
\end{proof}

\begin{theorem}[All determinantal ideals of the signed character tests]\leavevmode
\label{thm:periodminors400}
Fix an ordinary initial class \(a\), an \(R\)-basis \(y_s\) of
\(D_*\otimes R\) as in Proposition~\ref{prop:initialtest396}, and
a finite list of projected algebraic tests \(F_i\). At the chosen
terminal places \(P\), take \(1\le j_P<d\) and the characters
\(u\mapsto[u]^{j_P}\). Denote the matrix of periods on the tests
\(\phi_{v,y_s}\) by \(P_\chi\), with rows \(i\) and columns
\((s,v)\), where \(v\) runs through tuples of vector line
representatives. For a path \(\gamma\) starting at \(a\), write
\(a_{\gamma,P,h}\) for the coordinates of its composite terminal
covector in the initial frame. Form the moment matrix
\[
 S_{i,(s,m)}=\sum_{\gamma} W_{\gamma,F_i}(y_s)
                          \prod_{P,h}[a_{\gamma,P,h}]^{m_{P,h}},
 \qquad m_P\in\mathcal M_{j_P}.
\]
Put \(b_m=\sum_P(|m_P|-j_P)/d\). For a matrix \(U\) over
\(\operatorname{Frac}(R)\) let \(I_q(U)\) be the fractional ideal
generated by its minors of size \(q\), including the zero ideal.
For every \(q\ge1\),
\[
 I_q(P_\chi)=
  \sum_{\substack{L\text{ a set of }q\text{ columns }(s,m)}}
           p^{\sum_{(s,m)\in L}b_m} I_q(S_{:,L}) .
\]
In particular, if both matrices have rank at least \(q\), the
valuation loss from \(I_q(S)\) to \(I_q(P_\chi)\) lies between
the sum of the \(q\) smallest and the sum of the \(q\) largest
of the displayed column exponents \(b_m\), counted once per
column. The two matrices have the same rank.
\end{theorem}
\begin{proof}
Apply the lemma at every terminal place to the compact scalar in
Theorem~\ref{thm:covectorpaths352}. Each full path contributes
its original signed coefficient, evaluated at \(y_s\), multiplied
by the characters of its composite covectors on the initial
vectors. If a composite covector is zero at some place all
terms for that path vanish, since \(|m_P|>0\). Thus summation
over the actual paths gives the matrix factorization
\[
    P_\chi = S\,\operatorname{diag}_{(s,m)}
             \Bigl(\prod_P c_{j_P,m_P}\Bigr)\, E .
\]
Here \(E\) is block diagonal by \(s\). Each block is the tensor
product of the square monomial evaluation matrices from the
lemma, with rows indexed by monomials. In particular \(E\) is
invertible over \(R\). All entries of the diagonal matrix are
nonzero, with the specified valuations. The factorization uses
the same coupled last-edge variables and the same projection
inside every \(W_{\gamma,F_i}\).

Multiplication by an invertible integral matrix preserves every
minor ideal by Cauchy--Binet applied in both directions. Scaling
a column scales every minor containing it by the same scalar.
Applying these two facts to the factorization proves the formula
for fractional ideals. Taking the least valuation of a nonzero
minor gives the lower bound. A minor of \(S\) attaining its
least valuation gives the upper bound after scaling its \(q\)
columns. Invertibility over the fraction field gives equality
of the ranks.
\end{proof}

\begin{corollary}[The full integral initial table at the terminal prime]\leavevmode
\label{cor:fullperiodmoments400}
Keep the projected tests \(F_i\), and allow all ordinary initial
classes and all integral initial tables with separate scalar sums
zero. For each initial class \(a\) and each tuple of covector
lines \(L\), choose an \(R\)-basis \(y_{a,L,s}\) of its actual
coefficient lattice \(D_{a,L}\otimes R\). For every tuple of
nontrivial scalar characters \(j=(j_P)_P\) and every
\(m_P\in\mathcal M_{j_P}\), form the column
\[
  B_{i,(a,s,j,m)}=\sum_{\gamma\text{ from }a}
       W_{\gamma,F_i}\bigl(y_{a,[\eta_\gamma],s}\bigr)
                       \prod_{P,h}[a_{\gamma,P,h}]^{m_{P,h}} .
\]
Paths having a zero terminal pullback at any place are omitted.
For each \(q\), the ideal of all minors of size \(q\) of the
projected period map on this full integral domain is exactly
\(I_q(B)\). In particular each moment uses powers at most
\(p-1\) per residue coordinate, at every coefficient precision.
\end{corollary}
\begin{proof}
The lattice over an initial class depends only on the tuple
of lines, and the values in different refinements can be
chosen independently by Theorem~\ref{thm:pathnorm352}. Fix the
bases in the statement. They give integral coordinates for
functions valued in these lattices and intertwine local scalar
multiplication, since the scalars preserve each tuple of lines.
The scalar group at each place has order \(d\), a unit in \(R\),
and all its characters take values in \(\ZZ_p^\times\). Its
character projectors are integral and split the function space
as a direct sum. Summation along a scalar orbit acts, after
regarding the sum as constant on that orbit, by \(d\) on the
trivial character and by zero on the other characters. The
separate zero-sum condition therefore retains exactly the tuples
with every scalar character nontrivial.

For a fixed such tuple, the lemma and its tensor product show
that the monomials in the reduced coordinates of a covector
form an integral basis of scalar functions in that character.
Multiplying each by the basis section \(\eta\mapsto
 y_{a,[\eta],s}\) and extending by zero at the other classes
gives an integral basis of the full domain when all indices
vary. Evaluating the original period formula in this basis gives
exactly the displayed matrix. Changes of the chosen coefficient
bases or of the initial residue frames change the domain basis
by an invertible integral matrix. All its minor ideals are
therefore intrinsic. The degree bound is the bound on each
individual Teichm\"uller variable in the lemma.
\end{proof}

The last formula includes the line-dependent coefficient lattices
even at embeddings of large weight. It locates all possible
cancellation for a fixed list of projected tests in explicit
sums of the original arithmetic path coefficients. A uniform
bound for these signed minors and their comparison to the analytic
incoming row on the same intervals are still required in
Section~\ref{sec:missing}.


\subsection{The signed initial table at primes dividing the scalar group order}
\label{sec:signedscalarcharacters432}

Fix one of the words, translates and terminal primes \(T\) on
the common intervals of Section~\ref{sec:missing}. Put \(d=T-1\)
and let \(f\) be the number of terminal halves. Use the actual
initial coefficient lattices \(D_{a,L}\) of
Corollary~\ref{cor:fullperiodmoments400}. Here \(a\) is an
ordinary initial class and \(L\) a tuple of covector lines.
Choose a coefficient field \(J\) containing \(\mu_d\) and the
coefficients of the projected tests. Fix a representative of each
line. Over \(a,L\) a table is a function on \(G^f\) with values in
\(D_{a,L}\), where \(G=\mathbb F_T^\times\). The lattice is
constant on this orbit. A table belongs to \(\Lambda\) if its sum
in each scalar coordinate is zero with the others fixed. Take the
direct sum over all \(a,L\). For a tuple of characters
\(\boldsymbol\chi=(\chi_1,\ldots,\chi_f)\), all nontrivial, let
\(\Lambda_{\boldsymbol\chi}\) be the lattice of integral tables
transforming by their product on \(G^f\). Let
\[
        \Lambda_{\rm ch}=\sum_{\boldsymbol\chi}
                          \Lambda_{\boldsymbol\chi}\ \subseteq\Lambda .
\]
The summands are independent over \(J\). We compute this
inclusion integrally and apply it after contraction with the
prescribed signed period.

\begin{lemma}[Consecutive nodes of prime-power order]\leavevmode
\label{lem:scalarnodes432}
Let \(w\mid\ell\) be a finite place of \(J\), let
\(v_w(\ell)=1\), and suppose \(J\) contains a primitive
\(P\)-th root \(\zeta\), where \(P=\ell^a\) with \(a\ge1\).
For \(1\le n\le P\) take \(n\) consecutive nodes
\(\zeta^b,\ldots,\zeta^{b+n-1}\). The elementary divisors
of evaluation at these nodes on monomials of degrees
\(0,\ldots,n-1\) have the valuations
\[
     s_0=0,\qquad
     s_j=\frac{1}{\varphi(P)}\sum_{u=1}^j\ell^{v_\ell(u)}
                                 \quad(1\le j<n).
\]
\end{lemma}
\begin{proof}
If \(1\le u<P\), the order of \(\zeta^u\) is
\(\ell^{a-v_\ell(u)}\). For a primitive \(\ell^t\)-th root
\(\eta\) the product of \(1-\eta^c\) over unit classes \(c\)
modulo \(\ell^t\) is \(\Phi_{\ell^t}(1)=\ell\).
The reduction of \(\eta\) is one, so the ratio
\((1-\eta^c)/(1-\eta)\) has reduction \(c\ne0\).
All factors have the same valuation. It follows that
\(v_w(1-\zeta^u)=\ell^{v_\ell(u)}/\varphi(P)\).

Write \(x_i=\zeta^{b+i}\). Replace monomials by the monic
Newton polynomials \(Q_j(X)=\prod_{i<j}(X-x_i)\), starting
with \(Q_0=1\). This is an integral triangular change of basis
with determinant one. The matrix of the new evaluations is
lower triangular with diagonal \(Q_j(x_j)\) of valuation \(s_j\).
For \(i\ge j\), the valuation of \(Q_j(x_i)\) is
\(\varphi(P)^{-1}\) times the sum of \(\ell^{v_\ell(u)}\)
over \(u=i-j+1,\ldots,i\), with the empty sum for \(j=0\).
Every interval of \(j\) consecutive integers contains at least
\(\lfloor j/\ell^t\rfloor\) multiples of \(\ell^t\). Expanding
\(\ell^{v_\ell(u)}=1+\sum_{t=1}^{v_\ell(u)}
(\ell^t-\ell^{t-1})\) shows that the minimum of the sum occurs
on \(1,\ldots,j\). Dividing column \(j\) by \(Q_j(x_j)\)
therefore produces an integral matrix with diagonal one.
It is invertible. The \(s_j\) are nondecreasing, which proves
the assertion.
\end{proof}

\begin{proposition}[Elementary divisors of the integral character tables]\leavevmode
\label{prop:augmentationcharacters432}
For \(w\mid\ell\) with \(a=v_\ell(d)>0\), put
\(P=\ell^a\) and \(b=d/P\). At one scalar coordinate the
valuation list of the inclusion of the nontrivial character
lines into the integral zero-sum lattice is
\[
    (b-1)\text{ copies of }(s_0,\ldots,s_{P-1})
          \quad\text{and one copy of }(s_0,\ldots,s_{P-2}).
\]
For \(\Lambda_{\rm ch}\subseteq\Lambda\) on one tuple of lines,
the valuations are the sums of \(f\) entries independently
chosen from this list, counted with multiplicities and
repeated for each coefficient basis vector. At \(w\nmid d\)
all the elementary divisors of this inclusion are units.
\end{proposition}
\begin{proof}
Write \(G=C_b\times C_P\), with \(C_j\) cyclic of order \(j\).
The \(b\)-character projectors have denominator \(b\), a unit
in \(\mathcal O_{J,w}\). For each nontrivial \(C_b\)-character
the sum over \(C_b\) vanishes, so its part of the zero-sum
lattice is a full copy of \(\mathcal O_{J,w}^{P}\).
Its pure character vectors make the Fourier matrix
\((\zeta^{ij})_{0\le i,j<P}\). This matrix, or its transpose,
is evaluation on \(P\) consecutive nodes as in the lemma.

For the trivial \(C_b\)-character the remaining \(P\) values
have sum zero, since \(b\) is a unit. Forget their last
coordinate. This identifies the zero-sum lattice integrally
with \(\mathcal O_{J,w}^{P-1}\). The nontrivial
\(C_P\)-characters have columns
\((\zeta^{ij})_{0\le i<P-1}\), now with \(1\le j<P\).
Its transpose is evaluation at the \(P-1\) nodes starting at
\(\zeta\) on degrees \(0,\ldots,P-2\). The lemma gives the
shorter list.

For \(f\) scalar coordinates the zero-sum lattice on an orbit
is the tensor product of the \(f\) one-coordinate augmentation
lattices with \(D_{a,L}\). For example, split each function
space as its sum-zero kernel plus the span of the delta
function at the identity. In the tensor product the simultaneous
kernel consists precisely of the summands using the kernel in
every coordinate. The pure character lattices tensor in the
same way. Their inclusion is consequently diagonal in the
tensor products of the local bases from the lemma. Away from
\(d\) every character projector is integral and the nontrivial
tuples give the full simultaneous kernel. Taking the direct
sum over the orbits finishes the computation.
\end{proof}

\begin{theorem}[The loss in each projected signed minor]\leavevmode
\label{thm:signedcharacterloss432}
Let \(\mathcal P:\Lambda\to J^r\) evaluate the original signed
period at any fixed list of algebraic projected tests, with all
the weights \(W_{\gamma,F_i}\) of
Theorem~\ref{thm:periodminors400}. Enlarge to a \(J\)-linear
map on the span. For an integral source lattice \(H\), let
\(I_q(\mathcal P,H)\) be the fractional ideal generated by
all coordinate \(q\)-minors on \(H\). If the map has rank
\(\ge q\ge1\), then for every \(w\mid\ell\),
\[
 0\le v_w I_q(\mathcal P,\Lambda_{\rm ch})
             -v_w I_q(\mathcal P,\Lambda)
                  \le q f v_\ell(d).
\]
More precisely the two bounds can be replaced by the sums of
the smallest and largest \(q\) entries of the valuation list
of Proposition~\ref{prop:augmentationcharacters432} for the
whole source. At every complex embedding, in the induced
automorphic norms, the Gram operators on the common target
satisfy
\[
       \mathcal P\mathcal P^*
        =\sum_{\boldsymbol\chi}
              \mathcal P_{\boldsymbol\chi}
                   \mathcal P_{\boldsymbol\chi}^* ,
       \qquad
       \mathcal P_{\boldsymbol\chi}
          =\mathcal P|_{\Lambda_{\boldsymbol\chi}\otimes\CC}.
\]
\end{theorem}
\begin{proof}
Choose local bases diagonalizing the inclusion and write \(u_i\)
for the diagonal multipliers. In the basis for the larger lattice
let \(B\) be the actual signed period matrix. In the smaller
basis its matrix is \(B\operatorname{diag}(u_i)\).
Consequently
\[
 I_q(\mathcal P,\Lambda_{\rm ch})\mathcal O_{J,w}
   =\sum_{|L'|=q}\Bigl(\prod_{i\in L'}u_i\Bigr)
                      I_q(B_{:,L'})\mathcal O_{J,w}.
\]
Changes of either source basis preserve the minor ideals by
the two directions of the Cauchy--Binet formula. The minimum
among the valuations of the nonzero \(q\)-minors for the larger
lattice is attained. The displayed equality gives the lower
bound for all minors, and using a minor attaining that minimum
gives the upper bound.

In the full \(P\)-node list one has \(s_{P-1}=a\).
Indeed \(\prod_{j=1}^{P-1}(1-\zeta^j)=P\), by evaluating
\((X^P-1)/(X-1)\) at one. The shorter list is bounded by
the same endpoint. The tensor product therefore has entries
between zero and \(fa\), which proves the stated inequality.
This argument uses only a change of initial tests on the very
same map. In particular the arithmetic path sums have already
been taken when \(B\) is formed.

At a complex embedding the scalar action is unitary. On
each orbit it permutes values with equal class mass and
the same coefficient norm, as in
Proposition~\ref{prop:initialtest396}. Distinct character
spaces are orthogonal, by averaging their inner product over
the finite group. Their span in the zero-sum subspace is
everything, by character orthogonality at each coordinate.
Writing its identity as the sum of the orthogonal projections
and composing with \(\mathcal P\) and \(\mathcal P^*\) gives
the formula. It holds for the chosen tests after their full
signed contraction.
\end{proof}

\begin{corollary}[The total cost on the common period intervals]\leavevmode
\label{cor:jointscalarbudget432}
For these same tests of rank at least \(q\), the globally
normalized logarithmic increase in the finite content of their
\(q\)-minors, divided by \(q\), is between zero and
\[
                     f\log(T-1).
\]
In the period family of Theorems~\ref{thm:common111} and
\ref{thm:prime210} this is
\(O_{\psi,\varepsilon}(x^{5/2+143\varepsilon/16})\),
uniformly for all words and translates on the interval.
For the inclusion of the entire domain before applying the
period, its exact global logarithmic index per retained
dimension, normalized by \([J:\QQ]\), is
\[
             f\,\frac{d-2}{2(d-1)}\log d .
\]
\end{corollary}
\begin{proof}
Weight a local valuation over \(\ell\) by
\([J_w:\QQ_\ell]\log\ell/[J:\QQ]\) and sum. At each
\(\ell\) the sum of the local degrees equals \([J:\QQ]\).
The upper bound from the theorem gives \(f\log d\) per
minor dimension. The bounds \(f\ll_\psi x^{2+8\varepsilon}\),
\(\log T\ll m x^\varepsilon\) and
\(m\asymp x^{1/2-\varepsilon/16}\) give the specified power.
In particular this cost is \(o(\exp(x^\varepsilon)/x^A)\)
for every fixed \(A>0\).

For the exact domain index the full \(P\)-node Fourier
determinant has valuation \(Pa/2\). Indeed the Fourier matrix
times the inverse-character matrix is \(P I\), and replacing
\(\zeta\) by \(\zeta^{-1}\) merely permutes the columns.
The shorter list omits \(s_{P-1}=a\), so the total of the
one-coordinate list is \(a(d-2)/2\). Its length is \(d-1\).
Independent tensor choices add their averages, giving local
mean \(fa(d-2)/(2(d-1))\). All orbit and coefficient
multiplicities repeat this same mean. Summing with the local
degrees proves the exact formula.
\end{proof}

At \(T\), Corollary~\ref{cor:fullperiodmoments400} evaluates the
full table and Theorem~\ref{thm:initialindex396} computes the
cost of generating a fixed character lattice by vector tests.
At primes dividing \(T-1\), the latter tests generate each
individual character lattice there, and the calculation above
covers their sum inside the full table. The signed Gram formula
still involves the sizes of the restricted operators
\(\mathcal P_{\boldsymbol\chi}\). Comparing them on the
algebraic and analytic incoming rows remains the period
estimate in Section~\ref{sec:missing}.



\subsection{One normalization for the incoming automorphic row}
\label{sec:commonphase441}

Fix \(x\), the field \(E_x\), a translate and the marked finite level
from Section~\ref{sec:missing}. For each word \(u\) let \(\Phi_u(t)\)
be its automorphic Poincar\'e column with class data \(t\). This is the
column before pairing with the incoming row. We allow the signed
projections in the class data. For a column initially defined by
truncation, its pairing means the limit after orthogonal projection
as in Theorem~\ref{thm:global158}. We compare two ways of evaluating
these columns on the \emph{same} automorphic multiplicity space.
The calculation also applies to minors with columns from different
words.

Write \(G_\infty\) for the real group of the row, \(\pi\) for its
irreducible unitary discrete series and \(H_\pi\) for its Hilbert
space. On the adelic quotient by the rational group take the measure
induced by an adelic Haar measure \(\nu\) and counting measure on the
rational group. Let \(M\) consist of the bounded intertwiners
\(S:H_\pi\to L^2_{\rm cusp}(\nu)\) at the prescribed finite level
and in the selected finite types. Finite right translations commute
with \(G_\infty\), so these conditions can be imposed on the
multiplicity space. Throughout, inner products are linear in the
second argument.

The comparison uses a second Haar measure on an isomorphic adelic
group, a rational-point-preserving isomorphism \(\iota\) and a fixed
finite right translate \(t_0\). Pulling that Haar measure back gives
\(c\nu\) for some \(c>0\). Use \(\mu\) for its quotient measure.
For \(v\in H_\pi\) define \(U_vS\) by transporting \(Sv\) through
\(\iota\) and evaluating it at \(\iota^{-1}(g)t_0\). This defines
an element of \(L^2(\mu)\) modulo equality almost everywhere. In
this notation the level and its character on the second group
are transported by \(k\mapsto t_0^{-1}\iota^{-1}(k)t_0\).

\begin{lemma}[Projection between two evaluation lines]\leavevmode
\label{lem:rowoverlap441}
There is a positive definite inner product \(h_\nu\) on \(M\) such that
\[
  \langle U_w S,U_v T\rangle_\mu
       =c\,h_\nu(S,T)\langle w,v\rangle_\pi .
\]
For a unit vector \(w\), let \(W_w=U_w M\), with \(M\) the closed
multiplicity space, and let \(P_w\) be orthogonal projection onto
\(W_w\). Then \(P_w U_v T=\langle w,v\rangle_\pi U_w T\), and
\[
  \|(1-P_w)U_v T\|_\mu^2
     =c\,h_\nu(T,T)
         \bigl(\|v\|_\pi^2-|\langle w,v\rangle_\pi|^2\bigr).
\]
\end{lemma}
\begin{proof}
Since the group acts unitarily, \(S^*T\) is a bounded operator
commuting with the irreducible representation on \(H_\pi\).
Its real and imaginary self-adjoint parts have invariant spectral
projections, each zero or the identity. Both parts are scalars.
Define \(h_\nu(S,T)\) by \(S^*T=h_\nu(S,T)I\). For a unit vector
\(v\), \(h_\nu(S,S)=\|Sv\|_\nu^2\), and if this scalar vanishes
then \(S^*S=0\) and \(S=0\). This proves the inner product assertion
before transport. The norm on intertwiners is their operator norm,
so the closed conditions defining \(M\) give a Hilbert space.

The Haar scalar is common to all pairs of functions. Indeed the
quotient integration formula first sums a compactly supported
function over the rational subgroup and then integrates the sum on
the quotient. The rational subgroup and its counting measure are
preserved by \(\iota\), so multiplying adelic Haar by \(c\) multiplies
the quotient measure by \(c\). Right translation preserves that
measure by unimodularity of the adelic unitary group. The transported
inner product is therefore \(c\langle Sw,Tv\rangle_\nu\), which is
the asserted formula. For unit \(w\), \(U_w\) is a scaled isometry
with closed image. The formula shows that
\(U_vT-\langle w,v\rangle_\pi U_wT\) pairs to zero with every
element of this image. Subtraction of the squared norms gives the
second identity.
\end{proof}

Let \(P\) be the finite subgroup giving the character level in the
row model and put \(n=|P|\). The multiplicity coordinates of the
incoming row are endowed with the Petersson form for
\(\nu/n\). Write \(eF\in M\) for the intertwiner of a row \(F\)
in these coordinates. The form for \(\nu/n\) equals \(h_\nu/n\):
one obtains it by integrating the same product of evaluated functions
with the quotient measure divided by \(n\). Thus
\[
                    h_\nu(eF,eF)=n\|F\|^2 .
\]
The formulas for \(n\) and this normalization can be located in
\srcloc{Formalism/LocalToGlobal/IndefiniteReal/LinksDefs}{675}
and
\srcloc{Development/AbelianVarietyTheory/Proofs/IwasawaDecomposition/LevelRowsReadingEmbeddingGoodRowsModelWithNorm}{79}.

Take \(w\) to be the tensor product of the prescribed unit vectors
at the real places. In the common \(G_\infty\)-type of
Theorem~\ref{thm:global158}, the row at its line is \(W_w\).
The character and type conditions have the same meaning in both
evaluations. In fact
\((U_w S)(gk)\) evaluates \(S w\) at
\(\iota^{-1}(g)t_0(t_0^{-1}\iota^{-1}(k)t_0)\).
This intertwines every finite right translation with its transported
translation. It therefore intertwines the finite character
averages, and complex conjugation conjugates all their coefficients.
Let \(C\) denote pointwise complex conjugation to the conjugate
row. It is an antiunitary isomorphism.

\begin{theorem}[A common phase on the prescribed periodized columns]\leavevmode
\label{thm:commonphase441}
In the multiplicity coordinates above, the normalized conjugate row is
\[
       F\longmapsto \frac{\delta}{\sqrt{cn}}\ C U_w(eF),
                        \qquad |\delta|=1.
\]
Let a second isometric evaluation of the same rows in \(W_w\) be
\(JF=a U_v(eF)\), with the vector \(v\) and scalar \(a\) fixed
before the word is selected. In particular the choice
\(J=(cn)^{-1/2}U_w e\) is available in these coordinates.
Let \(\tau\) be the identity or complex conjugation, according to
the chosen embedding convention. For a fixed class table \(t\)
put
\[
 \begin{aligned}
    L_u(F,t)&=\langle CJF,\Phi_u(t)\rangle_\mu,\\
    B_u(F,t)&=\tau\!\left(
       \left\langle\frac{\delta}{\sqrt{cn}}\, C U_w(eF),
                        \Phi_u(t)\right\rangle_\mu\right).
 \end{aligned}
\]
There is a single \(\kappa\), with \(|\kappa|=1\), for which
\[
                       B_u(F,t)=\tau\bigl(\kappa L_u(F,t)\bigr)
\]
for every word, row and test. If the row space is nonzero, put
\(b=a\langle w,v\rangle_\pi\). Then
\(\kappa=\overline\delta/(\sqrt{cn}\,b)\).
\end{theorem}
\begin{proof}
On \(M\) the evaluation \(U_w\) multiplies squared norms by \(c\).
On incoming coordinates \(e\) multiplies them by \(n\). The unique
positive factor that makes the composite an isometry on a nonzero
row is \(1/\sqrt{cn}\). A choice of phase gives the displayed
normalization. For an arbitrary second evaluation in the theorem,
\(JF\in W_w\), so the projection formula gives
\(JF=b\,U_w(eF)\). If \(F\ne0\), its norm gives
\(|b|^2cn=1\). Moreover, \((1-P_w)U_v(eF)=0\) for a nonzero \(F\).
The second identity of the lemma then gives equality in
Cauchy--Schwarz for \(w,v\), so \(v\) lies in \(\mathbb C w\).

It remains to check that passing to coordinates on the columns
introduces the same test into both inner products. At fixed word,
record the real coefficient section on each finite class
representative, with the given transformation under its stabilizer.
These data determine the section on the whole class by the level
character and change of frame. Consistency for two representatives
of one argument follows from the stabilizer condition. In a basis
of the space of such data, write \(t=\sum_i z_i t_i\).
Reconstruction gives precisely the section with data \(t\).
In \(\tau\)-coordinates its entries are \(\tau(z_i)\), which become
\(z_i\) on return to the original coordinates. Periodizing the
reconstructed section gives \(\sum_i z_i\Phi_u(t_i)=\Phi_u(t)\).
This argument applies after the original signed projections and on
their integral tables. For the projected \(L^2\) limits it follows
first on truncations and then by continuity of the row pairing.
The coordinate maps used here are specified in
\srcloc{Formalism/LocalToGlobal/IndefiniteReal/SpOfRwRange}{87}.

Finally \(CJF=\overline b\,C U_w(eF)\).
Conjugate linearity of the first argument gives the factor \(b\)
in \(L_u\) and \(\overline\delta/\sqrt{cn}\) in the unconjugated
value defining \(B_u\). Their ratio proves the formula. Each factor
was fixed on the row, so it also works on zero columns. If the row
space is zero take \(\kappa=1\).
\end{proof}

\begin{corollary}[Mixed minors in the two analytic evaluations]\leavevmode
\label{cor:mixedphase441}
Choose any list of rows and any list of integral projected class
tests, allowing different words. Denote the resulting two matrices
by \(L,B\). Their column Gram matrices satisfy
\(B^*B=\tau(L^*L)\). Every \(j\)-rowed minor of \(B\) is
\(\tau(\kappa^j)\) times \(\tau\) of the corresponding minor of
\(L\). For integral tables over a coefficient ring \(O\subset
\mathbb C\), the \(\tau(O)\)-module in \(\mathbb C\) generated by
these minors of \(B\) equals \(\tau(\kappa^j)\) times the image
under \(\tau\) of the \(O\)-module generated by the minors of \(L\).
\end{corollary}
\begin{proof}
Entrywise the theorem gives \(B=\tau(\kappa L)\). Both possible
maps \(\tau\) commute with conjugation. Matrix multiplication
cancels \(\overline\kappa\kappa=1\) in each term of the Gram
matrix, including entries from different words. Multilinearity
of the determinant gives the formula for the minors. Taking all
finite linear combinations over the indicated rings proves the
module identity without assigning a finite-place valuation to a
complex phase.
\end{proof}

The two matrices here evaluate the same Poincar\'e columns. The
separate arithmetic column in Theorem~\ref{thm:archcolumn313}
comes from contracting the chain of CM theta sections. To apply
that theorem to these analytic minors one must still identify
this contraction on the specified tests with the Poincar\'e
pairing, with its arithmetic scale. In the archive the equality
on the independently specified coefficient vectors occurs as the
hypothesis on the period values in
\srcloc{Development/AbelianVarietyTheory/Proofs/DecoratedLatticeProofs/JoinSubspaceRealDataExists}{101}.

\clearpage
\section{References}\label{nav:bibliography-XXIV}
\begingroup\def\section*#1{}\begin{thebibliography}{99}\small
\bibitem{XXIV-XIX-bayer447} C. Bayer and J. Teichmann, \emph{The proof of Tchakaloff's theorem}, arXiv:math/0502473 (2005). \href{https://arxiv.org/abs/math/0502473}{Authors' preprint}.
\bibitem{XXIV-XIX-gurvits489} L. Gurvits, \emph{Van der Waerden/Schrijver--Valiant like conjectures and stable (aka hyperbolic) homogeneous polynomials}, Electron. J. Combin. \textbf{15} (2008), R66. \href{https://arxiv.org/abs/0711.3496}{Author's preprint}.
\bibitem{XXIV-XIX-oscillator519} S. Gurevich, R. Hadani and N. Sochen, \emph{The finite harmonic oscillator and its applications to sequences, communication, and radar}, IEEE Trans. Inform. Theory \textbf{54} (2008), 4239--4253, equation (I.1) and \S I-B. \href{https://people.math.wisc.edu/~sgurevich/IEEE-Journal.pdf}{Authors' paper}.
\end{thebibliography}\endgroup
\end{document}
